\chapter{The geometric bar complex}
\label{chap:bar-cobar}\label{ch:cobar}
\label{chap:bar-construction}

% Local macros (provide-command idempotent; defined for cross-chapter
% references to chap:climax-platonic and Theorem~A in
% theorem_A_infinity_2.tex, both of which use these macros internally).
\providecommand{\KZ}{\mathrm{KZ}}
\providecommand{\Arn}{\mathrm{Arnold}}
\providecommand{\dbar}{d_{\mathrm{bar}}}
\providecommand{\BRST}{\mathrm{BRST}}
\providecommand{\ConnConf}{\mathsf{ConnConf}}

\index{bar construction!as integral transform|textbf}
\index{non-abelian Fourier transform|textbf}

An augmented algebra has exactly one canonical input beyond its
multiplication: the augmentation ideal
$\bar\cA=\ker(\varepsilon)$, the part invisible to the counit.
Desuspend it (the shift $s^{-1}$ with $|s^{-1}v|=|v|-1$), then form
the cofree conilpotent coalgebra
$T^c(s^{-1}\bar\cA)$ with deconcatenation coproduct. This ordered
tensor coalgebra is the level-$2$ bar object before averaging. Its coderivations are
freely determined by their projection to $s^{-1}\bar\cA$, so the
chiral product and the internal differential determine a unique bar
coderivation. At genus~$0$ the Arnold relation and the Borcherds
identity force this coderivation to square to zero. The symmetric Ran
bar appears only after passing to coinvariants; the ordered kernel is
recoverable only under an explicit PBW/Koszul acyclicity hypothesis.
The geometric incarnation has two faces on algebraic curves: the
ordered tensor coalgebra and its symmetric shadow.

Locality of the chiral product forces a nilpotent coalgebra.
The OPE has poles of finite order along the diagonal; the logarithmic
normal form \(\eta_{ij}\), represented on an affine/formal genus-zero
screen by \(d\log(z_i-z_j)\), extracts their residues on the
Fulton--MacPherson compactification \(\FM_n(X)\); and the Arnold
relation makes \(d^2=0\) on that genus-zero screen.
The primitive construction is the ordered bar
$B^{\mathrm{ord}}(\cA) = T^c(s^{-1}\bar\cA)$: the cofree
tensor coalgebra with deconcatenation coproduct, carrying the
$R$-matrix, the KZ associator, and the Yangian deformation before
averaging. The factorization coalgebra
$\barB(\cA)=\mathrm{Sym}^c(s^{-1}\bar\cA)$ on the Ran space
$\Ran(X)$ (the ind-scheme of nonempty finite subsets of~$X$)
is its $\Sigma_n$-coinvariant quotient. It is the input to the
modular Maurer--Cartan element
$\Theta_\cA \in \MC(\gAmod)$, to Theorem~A, and to the shadow
obstruction tower $\Theta_\cA^{\leq r}$.
At genus~$g \geq 1$ the fiberwise differential acquires curvature
$d^2_{\mathrm{fib}} = \kappa(\cA) \cdot \omega_g$ from the Hodge bundle
(Chapter~\ref{chap:higher-genus}); on the Koszul locus, the cobar
functor inverts the symmetric bar:
$\Omega \circ \barB \xrightarrow{\sim} \mathrm{id}$
(Theorem~\ref{thm:bar-cobar-inversion-qi}). This is bar--cobar
inversion. It is not Verdier duality, not the construction of
$\cA^!$, and not a closed-sector bulk construction.
The three-step bar differential
$d = d_{\mathrm{internal}} + d_{\mathrm{residue}}
+ d_{\mathrm{de\,Rham}}$ that we built by hand for the Heisenberg
algebra in Chapter~\ref{ch:heisenberg-frame} has the following general
form. The three components
correspond to three types of geometry on the compactified
configuration space $\overline{C}_n(X)$:
\begin{center}
\small
\begin{tabular}{lll}
\textbf{Component} & \textbf{Geometric origin} & \textbf{Extracts} \\
\hline
$d_{\mathrm{internal}}$ & Internal differential of $\cA$ & Algebraic structure \\
$d_{\mathrm{residue}}$ & Residues along collision divisors $D_{ij}$ & OPE coefficients \\
$d_{\mathrm{de\,Rham}}$ & de~Rham differential on $\overline{C}_n(X)$ & Configuration geometry
\end{tabular}
\end{center}

\noindent
Theorem~\ref{thm:bar-nilpotency-complete} assembles these three
components into the full genus-$0$ collision differential
$\dzero = d_{\mathrm{internal}} + d_{\mathrm{residue}} + d_{\mathrm{de\,Rham}}$
and proves $\dzero^2 = 0$; the key input is the Arnold relation
(Theorem~\ref{thm:arnold-three},
Remark~\ref{rem:costello-gwilliam-factorization}).

The three-component decomposition is not the structural endpoint.
On the genus-$0$ affine/tangent KZ window, the finite-window ordered
bar superconnection is the pullback of one universal operator-valued
flat connection along the KZ functor:
\[
\nabla^{B,0}_{\cA,V,n}
=d_\cA+d_{\mathrm{dR}}
-\sum_{i<j}\widetilde\rho_{\cA,V,n}(t_{ij})\,\eta_{ij}
=\KZ_{\cA,V,n}^{*}(\nabla_{\Arn})
\]
in the category $\ConnConf_{\mathbb{P}^1}$ of meromorphic flat
superconnections on configuration spaces with KZ-type coefficients
\textup{(}Theorem~\ref{thm:climax-vol1-platonic} of
Chapter~\ref{chap:climax-platonic}, equation~\eqref{eq:climax-equation-G2}\textup{)}.
Arnold's universal flat connection
$\nabla_{\Arn} = d - \sum_{i<j} t_{ij}\eta_{ij}$ on the trivial
$U(\mathfrak{t}_n)$-bundle over $\Conf_n^{\mathrm{ord}}(X)$ is the
initial object of $\ConnConf_{\mathbb{P}^1}$
\textup{(}Definition~\ref{def:connconf-platonic}\textup{)};
every chiral algebra $\cA$ supplies the structure morphism
\(\KZ(\cA)\) that pulls \(\nabla_{\Arn}\) back to the bar
superconnection; the ordered bar differential on
\(B^{\mathrm{ord}}(\cA)\) is obtained by FM-boundary residue
realization. The Arnold three-term identity
\eqref{eq:arnold-three-term} is the codimension-two compatibility that
forces $\dbar^2 = 0$ \textup{(}Case~(ii) of
Theorem~\ref{thm:bar-nilpotency-complete}\textup{)}; the classical
Yang--Baxter equation on the family
$\{r^{(ij)}(z_{ij}) = \KZ^{*}(\eta_{ij})\}$ is its $r$-matrix-valued
incarnation. The three-component form is computationally explicit;
Theorem~\ref{thm:climax-vol1-platonic} identifies it with the
residue-realized form of a single universal pullback.

\smallskip
\noindent\textit{Notation.}
Throughout this chapter, the unadorned differential $d$ on the bar complex denotes the genus-$0$ collision differential $\dzero$ of Convention~\ref{conv:higher-genus-differentials}; it satisfies $\dzero^2 = 0$ (Theorem~\ref{thm:bar-nilpotency-complete}). The fiberwise curved differential $\dfib$ and the total corrected differential $\Dg{g}$ appear only in Chapter~\ref{chap:higher-genus}.

\medskip

\begin{convention}[Set notation and ordering; \ClaimStatusDefinitional]\label{conv:set-notation}
Throughout this chapter: collision divisors are $D_{ij}$ with $i < j$ (indices in increasing order); the hat notation $\widehat{ij}$ denotes \emph{omission} of both factors $\phi_i$ and $\phi_j$ after applying the OPE (we write $\widehat{ij}$ for the collision pattern itself, and $\widehat{\phi_i, \phi_j}$ when listing omitted terms in a tensor product).
\end{convention}

\begin{convention}[Bar coalgebra and Koszul dual algebra; \ClaimStatusDefinitional]
\label{conv:bar-coalgebra-identity}
\index{bar construction!as bar coalgebra|textbf}
\index{Koszul dual!coalgebra vs.\ algebra|textbf}
The primitive bar construction of an augmented chiral algebra~$\cA$ is
the ordered tensor coalgebra
\[
B^{\mathrm{ord}}_X(\cA)=T^c(s^{-1}\bar\cA),
\qquad \bar\cA=\ker(\varepsilon),
\]
with deconcatenation coproduct
\[
\Delta_{\mathrm{dec}}(x_1\otimes\cdots\otimes x_n)
=\sum_{i=0}^{n}
(x_1\otimes\cdots\otimes x_i)\otimes
(x_{i+1}\otimes\cdots\otimes x_n).
\]
The Ran-space bar is the coinvariant shadow
\[
\barB_X(\cA)=
\bigl(B^{\mathrm{ord}}_X(\cA)\bigr)_{\Sigma_\bullet}
\cong \mathrm{Sym}^c(s^{-1}\bar\cA),
\]
with coshuffle coproduct. The quotient
$\pi_\Sigma\colon B^{\mathrm{ord}}_X(\cA)\twoheadrightarrow
\barB_X(\cA)$ is the averaging/coinvariant map.  Its fibre is the
non-symmetric ordered kernel; a separate acyclicity theorem governs
reconstruction of that fibre. Throughout this chapter, the unqualified notation
$\barB_X(\cA)$ refers to this symmetric factorization coalgebra on
unordered $\operatorname{Ran}(X)$; the ordered tensor coalgebra is
always written $B^{\mathrm{ord}}_X(\cA)$.

Fix a quadratic presentation
$\cA=T_X(\mathcal V)/(R)$.  It defines the presentation coalgebra and
its canonical comparison
\[
\cA^i:=C_X(s^{-1}\mathcal V,s^{-2}R),
\qquad
q_\cA\colon\cA^i\longrightarrow\barB_X(\cA).
\]
The full bar coalgebra has cohomology
$H^\bullet(\barB_X(\cA))$.  Quadratic Koszulness is the assertion that
$q_\cA$ is a quasi-isomorphism; equivalently,
$H^\bullet(q_\cA)$ identifies $H^\bullet(\cA^i)$ with bar cohomology.
Diagonal concentration refers to the bar-weight/cohomological
bigrading.  Concentration in a single bar degree is the stronger
Gaussian special case.

On a finite-type quadratic lane, the presentation-dual algebra
$\mathbb D_{\Ran}(\cA^i)$ has generators $\mathcal{V}^\vee =
\mathcal{H}om_{\mathcal{O}_X}(\mathcal{V}, \omega_X)$
($\mathcal{D}_X$-module dual of the generators of~$\cA$) and
relations~$R^\perp$ (orthogonal complement under the residue
pairing).  Verdier duality on the full bar coalgebra forms the algebra
\[
K_X(\cA):=\mathbb{D}_{\operatorname{Ran}}\barB_X(\cA),
\qquad
\mathbb D(q_\cA)\colon
K_X(\cA)\longrightarrow\mathbb D_{\operatorname{Ran}}(\cA^i).
\]
A chosen Koszul-pair datum supplies
$\nu_\cA\colon K_X(\cA)\to\cA^!$; it is an equivalence on the named
finite or completed Koszul surface.
\begin{remark}[Verdier side of Theorem~A]
The Verdier construction uses the symmetric bar
$\barB_X(\cA)$ on unordered $\operatorname{Ran}(X)$.
The functor $\mathbb{D}_{\operatorname{Ran}}$ is defined on
factorization $\mathcal{D}$-modules on $\operatorname{Ran}(X)$;
the ordered bar lives on ordered configuration spaces.
The coinvariant map
$\pi_\Sigma\colon B_X^{\mathrm{ord}}(\cA)\to\barB_X(\cA)$
supplies the passage from the ordered configuration model to the
symmetric Ran ambient before Verdier duality is applied.
\end{remark}
\begin{remark}[Chiral and line duals]
Two duals must remain separated:
$\cA^!_{\mathrm{ch}}$ (the chiral dual, arising from Verdier duality on
the symmetric bar $\barB_X(\cA)$) and $\cA^!_{\mathrm{line}}$ (the line
dual, arising from linear duality on the ordered
bar $B^{\mathrm{ord}}_X(\cA)$). Throughout this volume, bare~$\cA^!$
denotes $\cA^!_{\mathrm{ch}}$ unless otherwise specified.
\end{remark}
The cobar--bar counit
$\Omega_X(\barB_X(\cA)) \xrightarrow{\sim} \cA$
(Theorem~\ref{thm:bar-cobar-inversion-qi})
is a \emph{resolution} of~$\cA$: the cobar of the bar coalgebra
recovers~$\cA$.  The Verdier algebra is formed separately by~$K_X$.

The bar-dual identity has a precise witness: the \emph{canonical
twisting morphism}
$\tau \colon \barB_X(\cA) \to \cA$,
the degree-$+1$ map projecting bar generators to their algebra
images. The graded vector space
$\mathfrak{g}_\tau := \Hom(\barB_X(\cA),\, \cA)$ carries a
natural dg Lie structure, the \emph{convolution Lie algebra}, whose
bracket is $[f, g] = \mu_\cA \circ (f \otimes g) \circ
\Delta_{\barB}$ and whose differential is
$\partial f = d_\cA \circ f - (-1)^{|f|} f \circ d_{\barB}$.
(This dg~Lie algebra is the truncation of a full $L_\infty$-structure
arising from the operadic convolution; the higher brackets
$\ell_n$ for $n \ge 3$ encode homotopy-transfer data and are
recovered via graph sums, see
Theorem~\ref{thm:modular-quantum-linfty}.)
The twisting morphism~$\tau$ is a \emph{Maurer--Cartan element}
of this Lie algebra:
\[
\partial\tau + \tfrac{1}{2}[\tau, \tau] = 0
\qquad
\bigl(\text{equivalently: } \partial\tau + \tau \star \tau = 0\bigr),
\]
and the bar-cobar adjunction is the pair
$(\varepsilon_\tau,\, \eta_\tau)$ it determines.
At higher genus, $\tau$ extends to a curved MC element
$\tau_{\mathrm{mod}} = \tau + \Theta_\cA$
in the modular convolution algebra
$\mathfrak{g}_\tau \widehat{\otimes} R\Gamma(\overline{\mathcal{M}}_{g,\bullet})$
(Remark~\ref{rem:theta-modular-twisting}), whose curvature
$\kappa(\cA)\cdot\omega_g$ controls the genus tower.
The four main theorems are four properties of~$\tau$:
existence~(A), invertibility~(B), branch structure~(C),
leading coefficient~(D).

\smallskip\noindent
\begin{center}
\small
\renewcommand{\arraystretch}{1.15}
\begin{tabular}{lll}
\textbf{Object} & \textbf{Type} & \textbf{Relationship} \\
\hline
$\barB_X(\cA)$ & factorization coalgebra
 & bar coalgebra \\[2pt]
$\cA^i=C_X(s^{-1}\mathcal V,s^{-2}R)$ & factorization coalgebra
 & quadratic presentation coalgebra; $q_\cA\colon\cA^i\to\barB_X(\cA)$
 \\[2pt]
$H^\bullet(\barB_X(\cA))$ & graded factorization coalgebra
 & bar cohomology; identified with $H^\bullet(\cA^i)$ by
   $H^\bullet(q_\cA)$ on the Koszul locus \\[2pt]
$\cA^!$ & chiral algebra
 & chosen Koszul partner; finite presentation model
   $\mathbb D_{\Ran}(\cA^i)$ has generators $\mathcal{V}^\vee$ and
   relations $R^\perp$ \\[2pt]
$\mathbb{D}_{\operatorname{Ran}}\barB_X(\cA)$
 & Verdier--Koszul chiral algebra
 & $=:K_X(\cA)$; maps by $\mathbb D(q_\cA)$ to
   $\mathbb D_{\Ran}(\cA^i)$ and by $\nu_\cA$ to $\cA^!$ \\[2pt]
$\Omega_X(\barB_X(\cA))$ & chiral algebra
 & $\simeq \cA$ (Ran reconstruction)
\end{tabular}
\end{center}
\end{convention}

\begin{convention}[Chiral product vs.\ chiral bracket; \ClaimStatusDefinitional]
\label{conv:product-vs-bracket}
\index{chiral product|textbf}
\index{chiral bracket|textbf}
The \emph{chiral product}
$\mu \colon j_* j^*(\cA \boxtimes \cA) \to \Delta_!(\cA)$
(Definition~\ref{def:chiral-algebra}) encodes all OPE pole
orders simultaneously. The \emph{chiral bracket}
$\{a,b\} = a_{(0)}b$ extracts the simple-pole residue of~$\mu$;
the \emph{invariant pairing} $\kappa(a,b) = a_{(1)}b$ extracts
the double-pole residue on quadratic current-type sectors.
The bar differential $d_{\mathrm{res}}$ uses the full product~$\mu$,
not only these two projections. Equivalently,
\[
d_{\mathrm{res}}=\sum_{m\ge0} d_{(m)},\qquad
d_{(m)}\text{ applies the OPE mode }a_{(m)}b.
\]
Thus $d_{\mathrm{bracket}}=d_{(0)}$, while the scalar curvature
projection is the relevant part of $d_{(1)}$ for affine and
Heisenberg current algebras and of higher modes in Virasoro/W-type
families. No individual pole-order summand is a differential in
general; the Borcherds identity makes their signed sum square to zero
(\S\ref{sec:selection-principle},
Proposition~\ref{prop:pole-decomposition}).
\end{convention}

\section{The geometric bar complex}
\label{sec:bar-cobar}

\subsection{OPE coefficients as residues}

\begin{remark}[Physical origin]\label{rem:physical-genesis}
In two-dimensional conformal field theory, the operator product expansion (OPE)
\[\phi_i(z)\phi_j(w) = \sum_{k} \frac{C_{ij}^k}{(z-w)^{h_i+h_j-h_k}} \phi_k(w) + \text{(less singular)}\]
encodes the structure constants $C_{ij}^k$ (fusion rules) and conformal weights $h_i$; associativity of multiple OPEs constrains the $C_{ij}^k$.
The bar construction geometrizes this data: $\overline{C}_n(X)$ parametrizes insertion points, $D_{ij}$ encodes the collision $z_i \to z_j$, the logarithmic normal form \(\eta_{ij}\) carries the singularity (locally \(d\log(z_i-z_j)\) on an affine/formal screen), and $\operatorname{Res}_{D_{ij}}$ extracts the OPE coefficients.
\end{remark}

\begin{remark}[Parallel track: the BV complex]\label{rem:bv-parallel-track}
\index{BV complex!parallel track}
The bar complex has a second life as the
BV complex of a chiral field theory. In the BV formalism
(Costello--Gwilliam~\cite{CG17}), a chiral algebra~$\cA$ determines a
classical field theory whose fields are generators of~$\cA$ and whose
\emph{antifields} are cogenerators of~$\cA^!$; the BV bracket arises from
the residue pairing on configuration spaces. The identification
\[
\text{bar complex of } \cA
\;=\; \text{BV complex of the chiral field theory of } \cA
\]
is made precise in Theorem~\ref{thm:brst-bar-genus0}. The dictionary
is complete:

\begin{center}
\small
\begin{tabular}{ll}
\textbf{This chapter} & \textbf{BV reading} \\ \hline
Bar differential $d_{\mathrm{res}}$ & BRST operator \\
Logarithmic form $\eta_{ij}$ & Propagator \\
Arnold relation $\Rightarrow d^2 = 0$ & Quantum master equation \\
Bar cohomology $H^*(\bar{B}(\cA))$ & Physical observables \\
Curvature $m_0$ (Chapter~\ref{chap:higher-genus}) & Conformal anomaly
\end{tabular}
\end{center}
\noindent
This is not an analogy; it is Theorem~\ref{thm:brst-bar-genus0}.
The algebraic-geometric side comes first because it admits complete
proofs; the physical identifications follow as theorems, not
heuristics.
\end{remark}

\begin{remark}[Holomorphic factorization on $\FM_k(X)$]
\label{rem:costello-gwilliam-factorization}
\index{Costello--Gwilliam!factorization algebra}%
\index{bar differential!as factorization structure}%
The bar differential $d_{\mathrm{res}}$, constructed from OPE
residues along collision divisors in $\overline{\Conf}_{k}(X)$,
is the factorization algebra structure of
Costello--Gwilliam~\cite{CG17} restricted to codimension-$1$
boundary strata of the Fulton--MacPherson
compactification~$\FM_k(X)$. Each boundary
stratum~$D_S \cong \FM_{|S|} \times \FM_{k-|S|+1}$ encodes a
cluster of $|S|$ points colliding, and the residue along~$D_S$
extracts the factorization product on the cluster.
The Arnold relation $\eta_{ij} \wedge \eta_{jk} + \text{cyc.} = 0$
is the codimension-$2$ cancellation that makes $d^2_{\mathrm{res}} = 0$
a consequence of the face relations on~$\FM_k(X)$.
Thus, at the level of operadic boundary strata of~$\FM_k(X)$,
the bar complex realizes holomorphic factorization in the sense
of Costello--Gwilliam~\cite{CG17}. This chapter uses only that
single-colour ordered construction and its Ran coinvariants. Any
open--closed or Swiss-cheese enhancement is a separate structure on
the derived-centre pair, not proof input for the bar differential
constructed here.
\end{remark}

\begin{example}[OPE as residue: the Heisenberg current]\label{ex:ope-to-residue}
The Heisenberg current $J(z)$ with OPE $J(z)J(w) = k/(z-w)^2 + \text{reg.}$ gives the bar element $J(z_1) \otimes J(z_2) \otimes \eta_{12} \in \bar{B}^1(\mathcal{H})$, where $\eta_{12} = d\log(z_1 - z_2)$. The bar differential extracts the full chiral product~$\mu$ (Convention~\ref{conv:product-vs-bracket}): since $J_{(0)}J = 0$ (no simple pole) and $J_{(1)}J = k$ (double pole), the result is
\[
d_{\mathrm{res}}(J \otimes J \otimes \eta_{12}) = k \cdot \mathbf{1}.
\]
The entire bar differential comes from the curvature component $d_{\mathrm{curvature}}$ (Proposition~\ref{prop:pole-decomposition}). One must not multiply the OPE pole by the propagator~$\eta_{12}$ and take a combined residue; the logarithmic form accounts for the bar degree, not an additional pole (Computation~\ref{comp:deg1-general}).
\end{example}

\begin{remark}[Logarithmic forms are forced]\label{rem:why-log-forced}
On an affine genus-zero collision screen, three requirements
distinguish the logarithmic representative
\(\eta_{ij}=d\log(z_i-z_j)\) among meromorphic \(1\)-forms with poles
along diagonals, and together determine the local normal form
uniquely. Globally, the same object is the logarithmic normal form
along \(D_{ij}\) together with its coordinate-change cocycle.

\emph{Conformal invariance.}
Under a conformal transformation $z \mapsto f(z)$, a direct computation gives
$d\log(f(z_i) - f(z_j)) = \frac{f'(z_i)dz_i - f'(z_j)dz_j}{f(z_i) - f(z_j)}$,
which near the diagonal $z_i \approx z_j$ reduces to $(dz_i - dz_j)/(z_i - z_j)$
by cancellation of $f'$. Thus $\eta_{ij}$ is conformally invariant up to regular terms.

\emph{Well-defined residues.}
The residue $\operatorname{Res}_{D_{ij}}$ requires a \emph{simple pole} along $D_{ij}$: forms with higher-order poles (e.g.\ $dz_i/(z_i - z_j)^2$) do not admit canonical residues. A logarithmic form $\omega = \frac{df}{f} \wedge \alpha + \beta$ with $f = z_i - z_j$ has $\operatorname{Res}_{D_{ij}}(\omega) = \alpha|_{D_{ij}}$, which is canonical.

\emph{Arnold relations.}
The cyclic identity
$\eta_{ij} \wedge \eta_{jk} + \eta_{jk} \wedge \eta_{ki} + \eta_{ki} \wedge \eta_{ij} = 0$
is the form-level expression of $d^2 = 0$ on configuration spaces; it holds for $d\log$ forms and fails for other meromorphic 1-forms with simple poles.

Conformal invariance, canonical residues, and differential nilpotence together force the logarithmic form uniquely.
\end{remark}

\subsection{Non-abelian Poincar\'e realization of the bar complex}

Factorization homology integrates a chiral algebra~$\cA$ over a space of configurations rather than over the space itself: each marked point contributes an insertion of~$\cA$, and the compactification~$\overline{C}_n(X)$ encodes how those insertions collide. Non-abelian Poincar\'e duality (NAP) identifies this integral with the bar complex.

\begin{definition}[Bar as factorization homology]\label{def:bar-fh}
The geometric bar construction is factorization homology of the
chiral algebra, following the Beilinson--Drinfeld factorization
framework~\cite[\S3.4]{BD04}:
\[\bar{B}^{\text{geom}}_n(\mathcal{A}) = \int_{\overline{C}_{n+1}(X)/X} \mathcal{A}\]
where we integrate over configuration spaces relative to $X$.
\end{definition}

\begin{remark}[Configuration spaces]\label{rem:why-config-NAP}
Non-abelian Poincar\'e duality~\cite[\S3.4]{BD04} integrates
over configuration spaces rather than the manifold itself. The
compactification $\overline{C}_n(X)$ adds boundary divisors encoding
collision data; the key inputs are~\cite[Theorem~3.4.22, \S3.6]{BD04}
(Ran space as the correct geometric setting). Residue extraction is
the NAP analogue of the cup product.
\end{remark}

\begin{theorem}[Bar construction as NAP homology; \ClaimStatusProvedHere]\label{thm:bar-NAP-homology}\label{thm:bar-computes-chiral-homology}
\index{factorization homology!via bar complex}
Let $\mathcal{A}$ be a chiral algebra on a smooth curve $X$ satisfying:
\begin{enumerate}[label=\textup{(\roman*)}]
\item \emph{Holonomicity}: $\mathcal{A}$ is a holonomic
 $\mathcal{D}_X$-module (equivalently, locally finitely generated
 with regular singularities). Concretely, the solution space of the
 $\mathcal{D}_X$-equations on any disk is finite-dimensional, the
 algebraic counterpart of finite-dimensional state spaces at each
 point.
\item \emph{Excision}: the factorization algebra associated
 to~$\mathcal{A}$ satisfies the Weiss cosheaf condition
 (Definition~\ref{def:weiss-cover}), i.e., the
 factorization structure glues over Weiss covers. This
 provides the descent data for factorization homology.
\end{enumerate}
Then factorization homology is computed via the geometric bar complex:
\[H_*(\bar{B}^{\text{geom}}(\mathcal{A})) \cong \int_{C_*(X)} \mathcal{A}\]

This is factorization homology of $X$ with coefficients in $\mathcal{A}$ in the sense of Ayala--Francis~\cite{AF15}. Holonomicity guarantees finite-dimensionality of stalks, which is needed for the Ayala--Francis identification~\cite[Theorem~5.1]{AF15}; excision provides the descent data.

The coalgebra structure on $\bar{B}^{\text{geom}}(\mathcal{A})$ arises from the coproduct in factorization homology:
\[\int_X A \to \int_{X_1} A \otimes \int_{X_2} A\]
when X decomposes as $X = X_1 \sqcup X_2$.
\end{theorem}

\begin{proof}
The identification follows from \cite{AF15} (for the factorization homology framework) and \cite{BD04} (Section~3.4, for the chiral algebra case). The bar differential $d = d_{\text{int}} + d_{\text{res}} + d_{dR}$ corresponds to the three components of the factorization structure:
\begin{itemize}
\item $d_{\text{int}}$: Internal operations in $\mathcal{A}$ (factorization algebra structure)
\item $d_{\text{res}}$: Residues at collisions (the full chiral product, realizing the NAP cup product)
\item $d_{dR}$: de Rham differential on logarithmic forms
\end{itemize}
The key verification is that Weiss covers of $C_n(X)$ compute factorization homology (\cite{AF15}, Theorem~5.1), and the FM compactification provides a canonical system of Weiss covers whose \v{C}ech complex is precisely the bar complex (Definition~\ref{def:weiss-cover}).
\end{proof}

\begin{remark}[BD comparison]\label{rem:bd-comparison-bar}
The comparison with Beilinson--Drinfeld is on the symmetric bar
$\barB_X(\cA)$ on $\Ran(X)$, not on the ordered tensor coalgebra
$B_X^{\mathrm{ord}}(\cA) = T^c(s^{-1}\bar\cA)$. The ordered object is an
$\Eone$ refinement that keeps the ordered collision data and the
$R$-matrix; BD's Ran-space formalism is permutation-invariant from the
start. Passing from $B_X^{\mathrm{ord}}(\cA)$ to $\barB_X(\cA)$ is the
convention bridge from the ordered $\Eone$-presentation to the
symmetric factorization-coalgebra formalism of
\cite[Theorem~3.4.9, §§3.4.11--3.4.12]{BD04}.

At genus~$0$, the differential
$d_{\mathrm{internal}} + d_{\mathrm{residue}} + d_{\mathrm{de\,Rham}}$
is the Fulton--MacPherson/logarithmic-form model for the same Cousin
differential that Beilinson--Drinfeld define abstractly on the Ran
stratification. The codimension-two compatibility forcing $d^2 = 0$ is
written here as the Arnold relation on logarithmic forms; in BD it is
packaged in the compatibility of the Cousin differential with
factorization. This identification is an inference from
\cite[§§3.4.11--3.4.12]{BD04} together with
Theorem~\ref{thm:geometric-equals-operadic-bar}; BD do not isolate the Arnold
relation under that name.

Chapter~4 of \cite{BD04} uses ``commutative'' for commutative
$\cD_X$-algebras, equivalently the pole-free subclass where the chiral
product extends across the diagonal; see \cite[Theorem~4.6.1]{BD04}.
This is strictly smaller than the $\Einf$-chiral regime used in this
monograph. The Heisenberg, affine Kac--Moody, and Virasoro algebras are
$\Einf$-chiral here because their OPEs are local and permutation
compatible, but they are not BD-commutative because their OPEs still
have poles. For the pole-free BD-commutative subclass, the symmetric bar
is the whole story. For the with-poles $\Einf$ case, the
ordered-to-unordered descent is $R$-twisted.

This also matches the chiral-homology computations of
\cite[Theorem~4.8.1]{BD04}: on the symmetric side, BD identify chiral
homology of chiral envelopes with a Chevalley complex of derived global
sections, while our Theorem~A/Theorem~B package uses the same symmetric
input and then adds the ordered bar, the descent datum, and the
bar-cobar inversion.
\end{remark}

\begin{remark}[Literature confirmation: Beilinson--Drinfeld]
The present bar construction agrees with the Beilinson--Drinfeld
factorization-coalgebra formalism at the symmetric level. Three points
match exactly.
\begin{enumerate}[label=\textup{(\roman*)}]
\item \textup{(Augmentation.)} The reduced bar used here,
 $T^c(s^{-1}\bar\cA)$ with $\bar\cA = \ker(\varepsilon)$, is the
 ordered refinement of the augmented symmetric object that appears in
 BD. On the BD side, the augmentation ideal enters through the kernel
 of the augmentation morphism for the chiral envelope in
 \cite[Theorem~4.8.1]{BD04}; our ordered bar keeps that reduced input
 before passing to $\Sigma_n$-coinvariants.
\item \textup{(Ran space and factorization.)} BD's principal theorem in
 \cite[Theorem~3.4.9]{BD04} identifies factorization algebras with
 chiral algebras on the Ran space, and the explicit factorization
 compatibility in \cite[(3.4.11.5)]{BD04} is the symmetric counterpart
 of our descent from $B_X^{\mathrm{ord}}(\cA)$ to $\barB_X(\cA)$ on
 $\Ran(X)$.
\item \textup{(Arnold relation.)} The codimension-two compatibility that
 BD package inside the Cousin/factorization differential
 \textup{(}\cite[(3.4.11.4)--(3.4.11.5), Lemma~3.4.12]{BD04}\textup{)}
 is the same nilpotence mechanism written here in logarithmic-form
 language as the Arnold relation. This is a reformulation, not a new
 differential.
\end{enumerate}
Hence there is no discrepancy with BD: our bar complex is an
$\Eone$-ordered lift of the BD symmetric factorization coalgebra, with
the same augmentation, the same Ran-space output, and the same
genus-$0$ square-zero relation after translating conventions.
\end{remark}

\begin{remark}[The nilpotence-periodicity correspondence]
\label{rem:nilpotence-periodicity}
\index{nilpotence-periodicity correspondence|textbf}
\index{periodicity!as exponential of nilpotence}
The logarithmic representative~$\eta_{ij}=d\log(z_i-z_j)$ is not an
incidental choice of kernel on the affine genus-zero screen: the bar construction is the
\emph{categorical logarithm} for chiral algebras, mapping
multiplicative structure (the OPE) to additive/nilpotent structure
(the bar differential with $\dzero^2 = 0$). The cobar construction
is the corresponding \emph{exponential}, reconstructing the algebra
from the nilpotent data. The structural principle governing the
entire theory is:
\begin{center}
\emph{Nilpotence is the logarithm of periodicity.}
\end{center}
At genus~$0$, the logarithm~$\eta_{ij}$ is single-valued
on~$C_2(\mathbb{P}^1)$, the bar differential is nilpotent
($\dzero^2 = 0$), and the monodromy is trivial. At
genus~$g \geq 1$, the logarithm $d\log E(z_i,z_j)$ acquires
monodromy around the $B$-cycles of~$\Sigma_g$, the fiberwise bar
differential fails nilpotence by
$\dfib^{\,2} = \kappa(\cA)\cdot\omega_g$, and the curvature
$\kappa(\cA)$ is the \emph{infinitesimal generator} of this
monodromy: the residue at the branch point of the logarithm.
This is the chiral Riemann--Hilbert correspondence: the nilpotent
datum~$\kappa$ determines the periodic datum (the monodromy of the
bar family over~$\overline{\mathcal{M}}_g$) via exponentiation.

The four theorems are four aspects of this correspondence:
Theorem~A asserts that the categorical logarithm and exponential
exist; Theorem~B that they are mutually inverse on the
convergence domain (the Koszul locus); Theorem~C that the failure
off the convergence domain decomposes into complementary
Lagrangians (the branch structure of the logarithm); and
Theorem~D that the leading coefficient~$\kappa(\cA)$ of the
logarithm determines the scalar genus expansion
$F_g^{\mathrm{sc}}(\cA)=\kappa(\cA)\cdot\lambda_g^{\mathrm{FP}}$.
On the uniform-weight exact lane the full coefficient satisfies
$F_g(\cA)=F_g^{\mathrm{sc}}(\cA)$ at all genera. For a multi-weight
algebra and $g\geq2$, the full coefficient is
$F_g(\cA)=F_g^{\mathrm{sc}}(\cA)+\delta F_g^{\mathrm{cross}}(\cA)$,
with the cross-channel correction supplied by
Theorem~\ref{thm:multi-weight-genus-expansion}.

The modular periodicity profile~$\Pi(\cA)$ is the
global/exponential counterpart of the infinitesimal/nilpotent
characteristic~$\kappa(\cA)$. At the scalar level, the T-matrix
eigenvalues $\lambda = \exp(2\pi i(h - c/24))$ are the
\emph{exponentials} of conformal weights shifted by the central
charge; the modular period~$N$ is the smallest positive integer for
which $\exp(2\pi i N c/24) = 1$. At the geometric level, the
nilpotency $\kappa(\lambda)^{g-1} = 0$
(Theorem~\ref{thm:geometric-depth-smooth}) makes the operator
$\exp(\kappa(\lambda))$ \emph{unipotent}, not periodic: geometry
provides a stabilization threshold, not a period. True periodicity
arises from modular and quantum sources whose nilpotent generators
have finite-order exponentials.
\end{remark}

\subsection{Construction of the geometric bar complex}

\begin{definition}[Genus-graded geometric bar complex, preliminary form]
\label{def:geom-bar-preview}
The bar complex at genus $g$ and bar degree~$n$ is, in first
approximation:
\[\bar{B}^{(g),n}(\mathcal{A}) = \Gamma\left(\overline{C}_{n+1}^{(g)}(\Sigma_g), j_*j^*\mathcal{A}^{\boxtimes(n+1)} \otimes \Omega^n(\log D^{(g)})\right)\]
The total bar complex is
$\bar{B}(\mathcal{A}) = \bigoplus_{g \geq 0} \bar{B}^{(g)}(\mathcal{A})$.
The complete definition, incorporating the bigrading by bar
degree~$p$ and form degree~$q$ together with the genus-graded
orientation bundle $\mathrm{or}_{p+1}^{(g)}$, the sign line that records
orientations of the codimension-$p$ collision strata at genus~$g$, is given in
Definition~\ref{def:geom-bar} below.
\end{definition}

\begin{convention}[Global sections vs.\ derived global sections]\label{conv:gamma-vs-rgamma}
\index{global sections!underived vs.\ derived}
Throughout this chapter, $\Gamma(-)$ denotes \emph{underived} global sections:
for a sheaf $\mathcal{F}$ on a variety $Y$, the symbol $\Gamma(Y, \mathcal{F})$
is the vector space $H^0(Y, \mathcal{F})$. When derived global sections are
needed (e.g., over moduli stacks or in derived $\mathcal{D}$-module arguments),
we write $R\Gamma$ explicitly. In particular, the bar complex definition above
uses honest global sections of a coherent sheaf on the smooth variety
$\overline{C}_{n+1}^{(g)}(\Sigma_g)$.
\end{convention}

\begin{remark}[Components of the definition]\label{rem:unpacking-bar-def}
The FM compactification $\overline{C}_{n+1}^{(g)}(\Sigma_g)$ parametrizes $(n+1)$ points on $\Sigma_g$ with collision patterns (Chapter~\ref{chap:config-spaces}). The degree-$n$ bar complex has $(n+1)$ insertions $\phi_0(z_0) \otimes \cdots \otimes \phi_n(z_n)$, where $\phi_0$ is the ``output'' (operadic structure).

The external tensor product $j_*j^*\mathcal{A}^{\boxtimes(n+1)}$ extends $\mathcal{A}^{\boxtimes(n+1)}$ from the open locus of distinct points across collisions, with the OPE controlling the singularity structure (BD~\cite[\S3.4.14--3.4.22]{BD04}).

The logarithmic forms $\Omega^n(\log D^{(g)})$ are spanned at genus~$0$ by wedge products of $\eta_{ij} = d\log(z_i - z_j)$; at genus $g \geq 1$, theta functions ($g=1$) and prime forms ($g \geq 2$) contribute.

A global section $\alpha \in \bar{B}^{(g),n}(\mathcal{A})$ is a ``correlation function'':
\[\alpha = \sum_I a_I(z_0, \ldots, z_n) \cdot \phi_{i_0}(z_0) \otimes \cdots \otimes \phi_{i_n}(z_n) \otimes \omega_I(z_0, \ldots, z_n)\]
the geometric incarnation of an $(n+1)$-point function.
\end{remark}

\begin{example}[Genus zero, degree 1]\label{ex:bar-genus0-deg1}
At genus 0, degree 1:
\[\bar{B}^{(0),1}(\mathcal{A}) = \Gamma\left(\overline{C}_2(\mathbb{P}^1), j_*j^*(\mathcal{A} \boxtimes \mathcal{A}) \otimes \Omega^1(\log D_{12})\right)\]

\emph{Configuration space.} $C_2(\mathbb{P}^1) \subset (\mathbb{P}^1)^2$ has $\dim_\mathbb{C} = 2$. The $\text{PSL}_2$ action is transitive on ordered pairs of distinct points, so the quotient $C_2(\mathbb{P}^1)/\text{PSL}_2$ is a point. After fixing the $\text{PSL}_2$ gauge (e.g., $z_2 = \infty$), the residual coordinate is $z_1 \in \mathbb{C}$.

\emph{Boundary divisor.} $D_{12} = \{z_1 = z_2\}$ is a single point in $\overline{C}_2(\mathbb{P}^1)$.

\emph{Logarithmic 1-forms.} $\Omega^1(\log D_{12})$ consists of forms:
\[\omega = f(z_1, z_2) \cdot \eta_{12}\]
where $\eta_{12} = \frac{dz_1 - dz_2}{z_1 - z_2}$ and $f$ is a meromorphic function.

\emph{Elements.} Typical element is:
\[\phi_i(z_1) \otimes \phi_j(z_2) \otimes \eta_{12}\]

\emph{Dimension.} If $\mathcal{A}$ has $N$ generators, then:
\[\dim \bar{B}^{(0),1}(\mathcal{A}) = N^2 \cdot \dim H^0(\overline{C}_2(\mathbb{P}^1), \Omega^1(\log D_{12}))\]

For $\mathbb{P}^1$, $\dim H^0(\overline{C}_2, \Omega^1(\log D)) = 1$ (only constant coefficient functions after fixing $\text{PSL}_2$).

Hence $\dim \bar{B}^{(0),1}(\mathcal{A}) = N^2$.
\end{example}

\begin{example}[Genus zero, degree 2]\label{ex:bar-genus0-deg2}
At genus 0, degree 2:
\[\bar{B}^{(0),2}(\mathcal{A}) = \Gamma\left(\overline{C}_3(\mathbb{P}^1), j_*j^*(\mathcal{A}^{\boxtimes 3}) \otimes \Omega^2(\log D)\right)\]

\emph{Configuration space.} $\overline{C}_3(\mathbb{P}^1)$ has $\dim_\mathbb{C} = 3$. $\text{PSL}_2(\mathbb{C})$ acts transitively on ordered triples of distinct points, so the quotient $\overline{C}_3(\mathbb{P}^1)/\text{PSL}_2 \cong \overline{\mathcal{M}}_{0,3}$ is a single point: the standard gauge $z_3 = 0,\ z_2 = 1,\ z_1 = \infty$ exhausts all three complex parameters of $\text{PSL}_2$. (Cross-ratios appear first at $4$ points: $\overline{C}_4(\mathbb{P}^1)/\text{PSL}_2 \cong \overline{\mathcal{M}}_{0,4} \cong \mathbb{P}^1$.)

\emph{Boundary divisor.} $D = D_{12} \cup D_{23} \cup D_{13}$ (three divisors, one for each pair of points colliding).

\emph{Logarithmic 2-forms.} $\Omega^2(\log D)$ is spanned by:
\[\eta_{12} \wedge \eta_{23}, \quad \eta_{23} \wedge \eta_{31}, \quad \eta_{31} \wedge \eta_{12}\]
subject to Arnold relation:
\[\eta_{12} \wedge \eta_{23} + \eta_{23} \wedge \eta_{31} + \eta_{31} \wedge \eta_{12} = 0\]

So the space of 2-forms is 2-dimensional (three generators, one relation).

\emph{Elements.} Typical element is:
\[\sum_{i,j,k} c_{ijk} \cdot \phi_i(z_1) \otimes \phi_j(z_2) \otimes \phi_k(z_3) \otimes (\eta_{12} \wedge \eta_{23})\]

\emph{Dimension.}
\[\dim \bar{B}^{(0),2}(\mathcal{A}) = N^3 \cdot 2\]

This grows rapidly with $n$.
\end{example}

\subsubsection{The bar differential}

\begin{definition}[Bar differential; \ClaimStatusDefinitional]\label{def:bar-differential-complete}%
\label{def:geometric-bar}\label{def:geometric-bar-definition}\label{def:bar-geometric}%
\label{def:bar-diff-detailed}\label{def:diff-total}%
\index{bar complex!differential|textbf}
The bar complex carries a bigrading
$\barB_{p,q}(\cA)$ by \emph{bar degree}~$p$ (number of
desuspended tensor factors minus one) and
\emph{form degree}~$q$. The \emph{cohomological
degree} of a homogeneous element
$s^{-1}\phi_0 \otimes \cdots \otimes s^{-1}\phi_p
\otimes \omega \in \barB_{p,q}$ is
\begin{equation}\label{eq:bar-cohomological-degree}
n \;=\; \sum_{i=0}^{p}\bigl(|\phi_i| - 1\bigr) + q
\;=\; \Bigl(\sum_{i=0}^{p} |\phi_i|\Bigr) - (p+1) + q,
\end{equation}
where $|\phi_i|$ is the internal (algebra) cohomological
degree and $-1$ is the desuspension shift. The total
differential $d = d_{\mathrm{internal}} +
d_{\mathrm{residue}} + d_{\mathrm{form}}$ has
$|d| = +1$, verified component by component:
\begin{itemize}
\item $d_{\mathrm{internal}}$ preserves $(p,q)$ and
 increases some $|\phi_i|$ by $1$: net $n \mapsto n+1$.
\item $d_{\mathrm{residue}}$ maps
 $\barB_{p,q} \to \barB_{p-1,q}$: bar degree
 decreases by $1$ (two factors merge via~$\mu$). The
 Poincar\'e residue removes the normal logarithmic form
 $d\log(z_i-z_j)$, while the codimension-one Gysin shift in
 $\Delta_!$ supplies the compensating cohomological degree; in this
 totalization the $q$-index is therefore unchanged. The
 desuspension contribution changes by $-(p) - (-(p+1)) = +1$:
 net $n \mapsto n+1$.
\item $d_{\mathrm{form}}$ maps
 $\barB_{p,q} \to \barB_{p,q+1}$: form degree
 increases by $1$ via~$d_{\mathrm{dR}}$:
 net $n \mapsto n+1$.
\end{itemize}
\noindent
On the primitive ordered bar write
\[
[\phi_0|\cdots|\phi_p]
:=s^{-1}\phi_0\otimes\cdots\otimes s^{-1}\phi_p.
\]
The ordered residue component is defined by first moving the two
colliding desuspended factors to adjacent positions using the Koszul
rule on the degrees $|\phi|-1$, then applying the desuspended chiral
product and the Poincar\'e residue along the corresponding collision
divisor. For adjacent factors this gives
\[
d_{\mathrm{residue}}
[\phi_0|\cdots|\phi_p]
=
\sum_{i=0}^{p-1}
(-1)^{\sum_{r<i}(|\phi_r|-1)+\epsilon_{\mathrm{or}}(D_{i,i+1})}
[\phi_0|\cdots|\mu(\phi_i,\phi_{i+1})|\cdots|\phi_p],
\]
with the logarithmic form factor restricted by
$\operatorname{Res}_{D_{i,i+1}}$. For a non-adjacent collision
$D_{ij}$ the same formula is conjugated by the unique order-preserving
shuffle that brings $i$ and $j$ together; the shuffle contributes only
the Koszul sign on the desuspended factors and the boundary-orientation
parity from Convention~\ref{conv:orientations-enhanced}. The symmetric
bar differential is the $\Sigma_{p+1}$-coinvariant image of this
ordered coderivation.

Thus
\[
d_B=d_{\cA}+d_{\mathrm{dR}}+d_{\mathrm{res}},
\]
and the signs are not independent notation. They are fixed by the
desuspended ordered tensor coalgebra and by the orientation convention
for Fulton--MacPherson boundary divisors. In the adjacent ordered
chart, for
\([a_1|\cdots|a_n]\otimes\alpha
=s^{-1}a_1\otimes\cdots\otimes s^{-1}a_n\otimes\alpha\), the local
formulas are
\[
d_{\cA}([a_1|\cdots|a_n]\otimes\alpha)
=
\sum_i
(-1)^{\sum_{q<i}(|a_q|-1)}
[a_1|\cdots|-d_{\cA}a_i|\cdots|a_n]\otimes\alpha,
\]
\[
d_{\mathrm{dR}}([a_1|\cdots|a_n]\otimes\alpha)
=
(-1)^{\sum_i |a_i|}
[a_1|\cdots|a_n]\otimes d_{\mathrm{dR}}\alpha,
\]
and
\[
d_{\mathrm{res}}([a_1|\cdots|a_n]\otimes\alpha)
=
\sum_{i=1}^{n-1}\sum_{r\ge0}
(-1)^{\sum_{q<i}(|a_q|-1)+|a_i|}
[a_1|\cdots|a_i{}_{(r)}a_{i+1}|\cdots|a_n]
\otimes
\operatorname{Res}^{\mathrm{form}}_{D_{i,i+1}}(\alpha).
\]
For a non-adjacent full-FM face, the adjacent formula is conjugated by
the order-preserving shuffle on the desuspended factors and multiplied
by the boundary-orientation parity of
Convention~\ref{conv:orientations-enhanced}. The consecutive-block
\(\chirAss\) quotient has only adjacent word-contractions; it is not
the full pairwise FM screen. These local formulas and their
anticommutation identities are proved in
Proposition~\ref{prop:ordered-bar-local-differential-identities}. The
finite-output and product-completion hypotheses under which the
aritywise sums define a complete conilpotent functor are part of
Theorem~\ref{thm:ordered-bar-complete-conilpotent-functor}.

\end{definition}

\begin{definition}[Ordered chiral bar term with logarithmic coefficients;
\ClaimStatusDefinitional]
\label{def:ordered-chiral-bar-term-log}
\index{ordered chiral bar!logarithmic coefficient complex}
\index{tensor powers!external, chiral, factorization}
Let $\varepsilon\colon\cA\to\omega_X$ be an augmentation of a
dg chiral algebra on a smooth curve~$X$, and put
$\bar\cA=\ker\varepsilon$. For $n\ge0$ let
$U_n(X)=X^n\setminus\bigcup_{i<j}\Delta_{ij}$, let
$j_n\colon U_n(X)\hookrightarrow X^n$ be the inclusion, and let
$q_n\colon\FM_n(X)\to X^n$ be the Fulton--MacPherson resolution with
normal-crossings boundary $D_n=\FM_n(X)\setminus U_n(X)$.

There are three tensor powers in the bar construction.
\begin{enumerate}[label=\textup{(\alph*)}]
\item The \emph{external tensor power} is
 $(s^{-1}\bar\cA)^{\boxtimes n}$ on $X^n$.
\item The \emph{chiral tensor power} is the meromorphic extension
 $j_{n,*}j_n^*((s^{-1}\bar\cA)^{\boxtimes n})$; the chiral product
 is read from its restrictions to the diagonals.
\item The \emph{factorization tensor power} is the restriction of this
 object to products of pairwise disjoint opens; on such a product it
 is the external product of the corresponding lower-arity factors.
\end{enumerate}
The logarithmic Fulton--MacPherson coefficient complex is
\[
\cK^{\ord}_n(\cA)
=
q_n^*j_{n,*}j_n^*((s^{-1}\bar\cA)^{\boxtimes n})
\otimes_{\cO_{\FM_n(X)}}
\Omega^\bullet_{\FM_n(X)}(\log D_n).
\]
Before taking global sections, the arity-$n$ ordered bar object is the
right $\cD_{X^n}$-module
\[
\mathbb B^{\ord}_{X,n}(\cA)
:=
(q_n)_+
\cK^{\ord}_n(\cA)
\in D(\cD_{X^n}),
\]
where $(q_n)_+$ is the de~Rham direct image. The ordered Ran object is
the collection
\[
\mathbb B^{\ord}_X(\cA)
:=\bigoplus_{n\ge0}\mathbb B^{\ord}_{X,n}(\cA)
\quad\text{on}\quad
\Ran^{\ord}(X)=\coprod_{n\ge0}U_n(X),
\]
with factorization restrictions along products of pairwise disjoint
opens. Thus the slogan
\[
B_X^{\ord}(\cA)=T^c(s^{-1}\bar\cA)
\]
means the underlying ordered tensor coalgebra of this aritywise
$\cD$-module system after forgetting the FM logarithmic coefficient
complex.
The $n$-th ordered chiral bar term is
\[
B^{\ord}_{X,n}(\cA)
=
R\Gamma_{\mathrm{dR}}\bigl(\FM_n(X),\cK^{\ord}_n(\cA)\bigr),
\qquad
B_X^{\ord}(\cA)=\bigoplus_{n\ge0}B^{\ord}_{X,n}(\cA).
\]
After forgetting the logarithmic coefficient model, the underlying
coalgebra is the ordered tensor coalgebra
$T^c(s^{-1}\bar\cA)$; the logarithmic forms record the
configuration-space residues that define the differential.
\end{definition}

\begin{proposition}[Ordered bar differential and local identities;
\ClaimStatusProvedHere]
\label{prop:ordered-bar-local-differential-identities}
\index{ordered chiral bar!local differential identities}
\index{bar differential!desuspended ordered formula}
\textup{[Type signature: Open quadrant; ordered
$\cD$-module/Fulton--MacPherson presentation; Beilinson level~$2$;
finite-window augmented complete $\chirAss$ or chiral
$A_\infty$ hypothesis package.]}

On the affine ordered collision screen, the differential on
$B_X^{\ord}(\cA)$ is
\[
d_B=d_\cA+d_{\mathrm{dR}}+d_{\mathrm{res}}.
\]
For a homogeneous word
$[a_1|\cdots|a_n]=s^{-1}a_1\otimes\cdots\otimes s^{-1}a_n$
and a logarithmic form $\alpha$, the three components are:
\begin{align}
d_\cA([a_1|\cdots|a_n]\otimes\alpha)
&=
\sum_{i=1}^{n}
(-1)^{\sum_{q<i}(|a_q|-1)}
[a_1|\cdots|-d_\cA a_i|\cdots|a_n]\otimes\alpha,
\label{eq:ordered-bar-dA}\\
d_{\mathrm{dR}}([a_1|\cdots|a_n]\otimes\alpha)
&=
(-1)^{\sum_i |a_i|}
[a_1|\cdots|a_n]\otimes d_{\mathrm{dR}}\alpha,
\label{eq:ordered-bar-ddR}\\
d_{\mathrm{res}}([a_1|\cdots|a_n]\otimes\alpha)
&=
\sum_{i=1}^{n-1}\sum_{r\ge0}
(-1)^{\sum_{q<i}(|a_q|-1)+|a_i|}
[a_1|\cdots|a_i{}_{(r)}a_{i+1}|\cdots|a_n]
\otimes\Res_{D_{i,i+1}}\alpha .
\label{eq:ordered-bar-dres}
\end{align}
Here $a_{(r)}b$ is the $r$-th OPE mode of the chiral product, and
$\Res_{D_{i,i+1}}$ is the Poincar\'e residue of the logarithmic form
factor. Thus the length-$2$ and length-$3$ residue computations are
\[
d_{\mathrm{res}}([a|b]\otimes\eta_{12})
=
\sum_{r\ge0}(-1)^{|a|}[a_{(r)}b],
\]
and
\[
\begin{aligned}
d_{\mathrm{res}}([a|b|c]\otimes\alpha)
&=
\sum_{r\ge0}(-1)^{|a|}
[a_{(r)}b|c]\otimes\Res_{D_{12}}\alpha\\
&\quad
-\sum_{r\ge0}(-1)^{|a|+|b|}
[a|b_{(r)}c]\otimes\Res_{D_{23}}\alpha .
\end{aligned}
\]
On the full pairwise Fulton--MacPherson screen, rather than the
strict adjacent $\chirAss$ quotient, the same convention includes the
non-adjacent face
\[
\sum_{r\ge0}
(-1)^{|a|+(|b|-1)(|c|-1)+\epsilon_{\mathrm{or}}(D_{13})}
[a_{(r)}c|b]\otimes\Res_{D_{13}}\alpha,
\]
where the block $\{1,3\}$ is ordered by its least original label and
$\epsilon_{\mathrm{or}}(D_{13})$ is the boundary-orientation parity of
Convention~\ref{conv:orientations-enhanced}. In the consecutive-block
ordered $\chirAss$ quotient this $D_{13}$ term is absent; in the full
FM model it is one of the three Arnold faces.
With the signs of
Convention~\ref{conv:A-infinity-2-ordered-bar-signs} and
Computation~\ref{comp:bar-signs-2-1},
\[
d_\cA^2=d_{\mathrm{dR}}^2=d_{\mathrm{res}}^2=0,\qquad
d_\cA d_{\mathrm{dR}}+d_{\mathrm{dR}}d_\cA=
d_\cA d_{\mathrm{res}}+d_{\mathrm{res}}d_\cA=
d_{\mathrm{dR}}d_{\mathrm{res}}+d_{\mathrm{res}}d_{\mathrm{dR}}=0.
\]
Consequently $d_B^2=0$. The deconcatenation coproduct
\[
\Delta_{\mathrm{dec}}[a_1|\cdots|a_n]
=
\sum_{p=0}^{n}
[a_1|\cdots|a_p]\otimes[a_{p+1}|\cdots|a_n]
\]
is the consecutive-cut convention, hence has no additional sign. For an
ordered factorisation cut obtained by first permuting the desuspended
word into blocks \(I\sqcup J=\{1,\ldots,n\}\), the convention is
\[
\Delta_{I\mid J}[a_1|\cdots|a_n]
=
\epsilon_{\mathrm{Kos}}\!\bigl(\pi_{I,J};
s^{-1}a_1,\ldots,s^{-1}a_n\bigr)
[a_I]\otimes[a_J],
\]
where \(\pi_{I,J}\) is the unique shuffle putting the labels of \(I\)
before the labels of \(J\), and \(\epsilon_{\mathrm{Kos}}\) is the usual
Koszul sign for permuting the desuspended homogeneous factors; \(a_I\)
and \(a_J\) denote the corresponding subwords in their induced orders.
With these signs the coproduct
is coassociative, $d_B$ is a coderivation for it, and
$B_X^{\ord}(\cA)$ is conilpotent by word length. Its completion
$\widehat B_X^{\ord}(\cA)=\prod_{n\ge0}B^{\ord}_{X,n}(\cA)$ is complete
for the length filtration, and every finite weight window is
conilpotent after quotienting by sufficiently long words. If
$\varepsilon$ is a morphism of dg chiral algebras, then
$\bar\cA=\ker\varepsilon$ is preserved by $d_\cA$ and by the reduced
OPE differential, so the reduced bar removes all vacuum degeneracies;
the length-$0$ summand is the coaugmentation and the projection to it
is the counit.
\end{proposition}

\begin{proof}
The formulas are the desuspended coderivation formulas recorded in
Convention~\ref{conv:A-infinity-2-ordered-bar-signs}. The unary part is
$b_1(s^{-1}a)=-s^{-1}(d_\cA a)$; the binary OPE-mode part is
$b_2^{(r)}(s^{-1}a\otimes s^{-1}b)=(-1)^{|a|}s^{-1}(a_{(r)}b)$, and
extension as a coderivation gives the prefix signs in
\eqref{eq:ordered-bar-dA} and~\eqref{eq:ordered-bar-dres}. The
de~Rham sign in~\eqref{eq:ordered-bar-ddR} is the geometric
totalization sign of Definition~\ref{def:bar-differential-complete}.
The displayed length-$2$ and length-$3$ identities are the special
cases computed in
Computation~\ref{comp:bar-signs-2-1}. For the full pairwise
length-$3$ screen, the $D_{13}$ term is obtained by moving the
desuspended factor $s^{-1}c$ across $s^{-1}b$, giving the shuffle sign
$(-1)^{(|b|-1)(|c|-1)}$, then applying the same binary projection
$b_2^{(r)}(s^{-1}a\otimes s^{-1}c)=(-1)^{|a|}s^{-1}(a_{(r)}c)$ and the
boundary-orientation sign of $D_{13}$.

The diagonal identities follow from the three source complexes:
$d_\cA^2=0$ because $\cA$ is dg, $d_{\mathrm{dR}}^2=0$ for logarithmic forms,
and $d_{\mathrm{res}}^2=0$ from the codimension-two calculation of
Theorem~\ref{thm:bar-nilpotency-complete}. In the strict ordered
$\chirAss$ quotient this is the consecutive-block Stasheff identity of
Construction~\ref{constr:bar-diff-collision}; in the full pairwise
Fulton--MacPherson screen it is the Borcherds identity multiplied by
Arnold's three-term logarithmic relation
Theorem~\ref{thm:arnold-three}. The mixed anticommutators are the three
compatibilities proved in the last paragraph of
Theorem~\ref{thm:bar-nilpotency-complete}: the chiral product is a
chain map, logarithmic de~Rham forms obey the standard tensor-product
sign rule, and Poincar\'e residue commutes with $d_{\mathrm{dR}}$ up to the
Koszul sign fixed by Convention~\ref{conv:orientations-enhanced}.

Coassociativity of $\Delta_{\mathrm{dec}}$ is the equality of the two ways of
cutting an ordered word into three consecutive pieces. Since $d_B$ is
defined as the coderivation extending its projections to
$s^{-1}\bar\cA$, the coderivation identity follows from the universal
property of the cofree tensor coalgebra. Iterating the reduced
coproduct on a word of length $n$ gives zero after more than $n$
iterations, which proves conilpotence; the completed statement is the
finite-window form of the length filtration in
Convention~\ref{conv:bar-length-filtration}. Finally, a morphism of
augmented dg chiral algebras satisfies
$\varepsilon d_\cA=d_{\omega_X}\varepsilon$ and
$\varepsilon\mu=\mu_{\omega_X}(\varepsilon\boxtimes\varepsilon)$, so
the reduced differential preserves $\ker\varepsilon$ after projecting
away the scalar vacuum term. This is exactly the reduced-bar
normalization.
\end{proof}

\begin{definition}[Ordered Fulton--MacPherson quotient and symmetric
descent datum; \ClaimStatusDefinitional]
\label{def:ordered-fm-sigma-descent-datum}
\index{Fulton--MacPherson compactification!ordered quotient}
\index{symmetric bar!descent datum}
For \(n\ge0\), the ordered Fulton--MacPherson compactification is the
labelled compactification
\[
\FM^{\ord}_n(X):=\FM_n(X),
\]
whose open stratum is \(U_n(X)\). The symmetric group \(\Sigma_n\)
acts by relabelling markings and collision trees; it preserves the
normal-crossings boundary \(D_n=\FM_n(X)\setminus U_n(X)\). The
unordered Fulton--MacPherson target is the finite quotient stack
\[
\FM^\Sigma_n(X):=[\FM^{\ord}_n(X)/\Sigma_n],
\]
or its coarse quotient when the ambient category has already been
restricted to coarse spaces. The quotient morphism is denoted
\[
q_n\colon \FM^{\ord}_n(X)\longrightarrow \FM^\Sigma_n(X).
\]

A \(\Sigma_\bullet\)-descent datum on the ordered bar system
\(\mathbb B^{\ord}_{X,\bullet}(\cA)\) consists, for every \(n\) and
every \(\sigma\in\Sigma_n\), of isomorphisms
\[
\rho_{\sigma,n}\colon
\sigma^*\mathbb B^{\ord}_{X,n}(\cA)
\xrightarrow{\ \sim\ }
\mathbb B^{\ord}_{X,n}(\cA)
\]
satisfying the usual cocycle
\[
\rho_{\sigma\tau,n}
=\rho_{\sigma,n}\circ\sigma^*(\rho_{\tau,n})
\]
after the standard identification of pullbacks. These isomorphisms
must commute with \(d_\cA\), \(d_{\mathrm{dR}}\),
\(d_{\mathrm{res}}\), the ordered deconcatenation coproduct, and the
factorisation restriction maps for disjoint opens. In the pole-free
case \(\rho_{\sigma,n}\) is the Koszul-signed permutation action. In
the with-poles ordered \(\Eone\)-chiral case it is the
strong-unitary \(R\)-twisted action generated by adjacent
\(\tau_{i,i+1}\circ R_{i,i+1}(z_i-z_{i+1})\), when the
Yang--Baxter and unitarity relations of
Lemma~\ref{lem:R-twisted-descent} are satisfied.
\end{definition}

\begin{proposition}[Symmetric bar descent criterion;
\ClaimStatusConditional]
\label{prop:symmetric-bar-descent-criterion}
\index{symmetric bar!descent criterion}
\index{ordered-to-symmetric bar descent}
Type signature: \textup{(}Open quadrant, ordered logarithmic
Fulton--MacPherson/Ran presentation descended to the symmetric Ran
quotient, Beilinson level~$2$, hypothesis package: the descent datum of
Definition~\ref{def:ordered-fm-sigma-descent-datum}, or equivalently
the pole-free equivariant datum / strong-unitary \(R\)-twisted datum in
the relevant case\textup{)}.
The symmetric bar
\[
B^\Sigma_X(\cA):=\bigl(B_X^{\ord}(\cA)\bigr)_{\Sigma_\bullet}
\]
exists as a dg factorization coalgebra on $\Ran(X)$ precisely when the
ordered coefficient system on $\Ran^{\ord}(X)$ carries a
\(\Sigma_\bullet\)-descent datum of
Definition~\ref{def:ordered-fm-sigma-descent-datum}. In the pole-free
symmetric case this is ordinary equivariant descent. In the with-poles
$\Eone$-chiral case it is the strong-unitary
$R$-twisted descent of Lemma~\ref{lem:R-twisted-descent}; without that
locality/descent datum, $B^\Sigma_X(\cA)$ is only a formal quotient of
graded objects and not a well-defined dg factorization coalgebra.
\end{proposition}

\begin{proof}
The quotient map
$\Ran^{\ord}(X)\to\Ran(X)=\Ran^{\ord}(X)/\Sigma_\bullet$ is finite
etale on each configuration stratum. Descent of a $\cD$-module with
differential across this quotient is equivalent to a compatible
$\Sigma_\bullet$-equivariant structure on each arity, with the cocycle
condition on double overlaps. Compatibility with $d_\cA$,
$d_{\mathrm{dR}}$, $d_{\mathrm{res}}$, and $\Delta_{\mathrm{dec}}$ is exactly the assertion
that the quotient differential and coshuffle coalgebra structure are
well-defined. The pole-free case is the usual Koszul-signed
permutation action. With poles, analytic continuation around
diagonals supplies braid monodromy, so a symmetric quotient exists only
after the braid action factors through $\Sigma_\bullet$ by the
strong-unitary $R$-twisted descent package. This is the criterion
already isolated in Definition~\ref{def:ordered-fact-dmod},
Theorem~\ref{thm:three-bar-complexes}, and
Lemma~\ref{lem:R-twisted-descent}.
\end{proof}

\begin{proposition}[Ordered bar carrier and ordered Ran factorisation;
\ClaimStatusProvedHere]
\label{prop:ordered-bar-carrier-ran-factorisation}
\index{ordered chiral bar!tensor carrier}
\index{ordered Ran space!bar factorisation coalgebra}
Type signature: \textup{(}Open quadrant, ordered
$\cD$-module/Fulton--MacPherson/Ran presentation, Beilinson level~$2$,
hypothesis package: augmented complete finite-window chiral
$\Eone$ algebra, holonomic/excision factorisation descent, and the sign package of
Theorem~\ref{thm:bar-sign-coherence}\textup{)}.
Let $\cA$ be an augmented chiral algebra satisfying the stated
finite-window hypotheses. Then:
\begin{enumerate}[label=\textup{(\roman*)}]
\item The ordered bar object
$\mathbb B_X^{\ord}(\cA)$ of
Definition~\ref{def:ordered-chiral-bar-term-log} is, after forgetting
the logarithmic FM coefficient system, exactly the cofree ordered
tensor coalgebra
$T^c(s^{-1}\bar\cA)$ with deconcatenation.
\item On the ordered Ran space
\[
\Ran^{\ord}(X)=\coprod_{n\ge0}U_n(X),
\]
the aritywise objects $\mathbb B^{\ord}_{X,n}(\cA)$ form a
factorisation coalgebra. For pairwise disjoint opens
$V_1,\ldots,V_k\subset X$ and $n=n_1+\cdots+n_k$, the restriction map is
the composite
\[
\mathbb B^{\ord}_{X,n}(\cA)|_{U_{n_1}(V_1)\times\cdots\times U_{n_k}(V_k)}
\longrightarrow
\boxtimes_{\alpha=1}^k
\mathbb B^{\ord}_{V_\alpha,n_\alpha}(\cA|_{V_\alpha}),
\]
obtained by restricting external tensor factors and logarithmic forms
and then applying the ordered deconcatenation cut at the block
boundaries.
\item Residues are compatible with this factorisation product. A
collision divisor meeting
$U_{n_1}(V_1)\times\cdots\times U_{n_k}(V_k)$ lies entirely inside one
block, since the opens are disjoint; hence the Poincare residue acts
on that block and commutes with the external product on the other
blocks, with the Koszul sign already fixed in
Proposition~\ref{prop:ordered-bar-local-differential-identities}.
\end{enumerate}
\end{proposition}

\begin{proof}
Part~\textup{(i)} is the definition of the ordered reduced bar term
together with the universal property of the cofree tensor coalgebra:
the coefficient object is \(s^{-1}\bar\cA\), the coproduct is
deconcatenation, and the logarithmic FM factor records the collision
forms without changing the ordered tensor carrier.

For~\textup{(ii)}, restrict the FM space to the open locus where the
first $n_1$ labels lie in $V_1$, the next $n_2$ labels lie in $V_2$,
and so on. Since the opens are pairwise disjoint, no cross-block
diagonal meets this locus. The FM boundary and logarithmic form system
therefore factor as the product of the FM boundary systems in the
blocks. The coefficient system
$j_{n,*}j_n^*((s^{-1}\bar\cA)^{\boxtimes n})$ restricts to the
external product of the lower-arity coefficient systems, and
deconcatenation at the block boundaries gives the displayed
factorisation map. The cosheaf/factorisation assertion follows from
the holonomic/excision hypothesis, as in
Theorem~\ref{thm:bar-NAP-homology}.

For~\textup{(iii)}, on a product of disjoint opens the divisor
$D_{ij}$ is empty unless labels $i$ and $j$ lie in the same block. If
they do lie in the same block, the local logarithmic normal parameter
is a coordinate on that block's FM factor. Poincare residue extracts
that normal factor and leaves all other blocks untouched; moving the
residue past the other coefficient factors gives exactly the Koszul
sign appearing in~\eqref{eq:ordered-bar-dres}. Hence
$d_{\mathrm{res}}$ is compatible with the factorisation product.
\end{proof}

\begin{corollary}[\texorpdfstring{$\chirAss$}{chiral associative}
dual-cooperad carrier; \ClaimStatusConditional]
\label{cor:ordered-bar-chirass-dual-carrier}
\index{associative chiral operad!Koszul dual cooperad}
\index{ordered chiral bar!conditional operadic carrier}
Type signature: \textup{(}Open quadrant, ordered
$\cD$-module/Fulton--MacPherson/Ran presentation, Beilinson level~$2$,
hypothesis package:
Proposition~\ref{prop:ordered-bar-carrier-ran-factorisation} plus the
definition of the quadratic dual cooperad
\((\chirAss)^{!,c}\) in
Definition~\ref{def:chirass-dual-cooperad} and the
finite locally free quadratic self-duality package of
Proposition~\ref{prop:chirAss-self-dual}\textup{)}.
Under the finite locally free hypotheses of
Proposition~\ref{prop:chirAss-self-dual}, the associative chiral
operad has quadratic dual cooperad
\[
(\chirAss)^{!,c}\cong(\chirAss\otimes\operatorname{sgn})^c,
\]
and the ordered bar object \(\mathbb B_X^{\ord}(\cA)\) is a coalgebra
over this cooperad: arity \(n\) is
\((s^{-1}\bar\cA)^{\otimes n}\), and every cooperadic decomposition
map is the ordered consecutive deconcatenation cut, with the Koszul
signs of the desuspended factors. Without those finite locally free
hypotheses, the
proved carrier is the ordered tensor coalgebra and ordered Ran
factorisation of
Proposition~\ref{prop:ordered-bar-carrier-ran-factorisation}; the
\((\chirAss)^{!,c}\)-identification is not asserted.
\end{corollary}

\begin{proof}
Proposition~\ref{prop:ordered-bar-carrier-ran-factorisation} gives the
ordered tensor coalgebra \(T^c(s^{-1}\bar\cA)\) with deconcatenation.
Definition~\ref{def:chirass-dual-cooperad} defines
\((\chirAss)^{!,c}\) only in the finite quadratic model where dual
arity pieces exist.
Proposition~\ref{prop:chirAss-self-dual}, under its finite locally
free hypotheses, identifies the relevant quadratic dual cooperad with
the sign-twisted coassociative chiral cooperad. The cofree coalgebra
over that cooperad is exactly the ordered tensor coalgebra on
\(s^{-1}\bar\cA\); its structure maps are the consecutive cuts of
Proposition~\ref{prop:ordered-bar-carrier-ran-factorisation}. Hence
the cooperadic carrier follows under the stated hypotheses and only
under those hypotheses.
\end{proof}

\begin{theorem}[Ordered bar descent to the symmetric Ran quotient;
\ClaimStatusConditional]
\label{thm:ordered-bar-operadic-ran-descent}
\index{ordered chiral bar!symmetric descent}
\index{symmetric bar!conditional factorisation coalgebra}
Type signature: \textup{(}Open quadrant, ordered
$\cD$-module/Fulton--MacPherson/Ran presentation compared with the
symmetric Ran quotient, Beilinson level~$2$, hypothesis package:
Proposition~\ref{prop:ordered-bar-carrier-ran-factorisation} plus the
ordered-to-symmetric descent datum of
Definition~\ref{def:ordered-fm-sigma-descent-datum} and
Proposition~\ref{prop:symmetric-bar-descent-criterion}\textup{)}.
The labelled Fulton--MacPherson compactification maps to the unordered
compactification by the finite quotient
\[
q_n\colon
\FM^{\ord}_n(X)\longrightarrow \FM^\Sigma_n(X)
\]
of Definition~\ref{def:ordered-fm-sigma-descent-datum}.
If the descent criterion of
Proposition~\ref{prop:symmetric-bar-descent-criterion} is satisfied,
then the symmetric bar
\[
B^\Sigma_X(\cA)=(B_X^{\ord}(\cA))_{\Sigma_\bullet}
\]
is a dg factorisation coalgebra on $\Ran(X)$ and the quotient
$B_X^{\ord}(\cA)\to B^\Sigma_X(\cA)$ is a morphism of dg coalgebras.
Without the descent datum, only the ordered factorisation coalgebra of
Proposition~\ref{prop:ordered-bar-carrier-ran-factorisation} is
asserted.
\end{theorem}

\begin{proof}
The labelled carrier and ordered factorisation structure are
Proposition~\ref{prop:ordered-bar-carrier-ran-factorisation}. The
action of $\Sigma_n$ relabels the markings and preserves the
normal-crossings boundary, giving the finite map from the labelled
compactification to its unordered quotient. Descent of the bar
differential and coalgebra structure through this quotient is precisely
the content of
Proposition~\ref{prop:symmetric-bar-descent-criterion}. Under that
criterion the quotient map commutes with $d_\cA$, $d_{\mathrm{dR}}$,
$d_{\mathrm{res}}$, and deconcatenation, hence is a morphism of dg
coalgebras. If the criterion is not supplied, the coinvariant quotient
has not been shown to carry a compatible dg factorisation coalgebra
structure.
\end{proof}

\begin{proposition}[First averaging kernels and the associator
class; \ClaimStatusConditional]
\label{prop:ordered-bar-low-arity-averaging-kernels}
\index{averaging map!low arity kernel}
\index{ordered bar!arity two and three averaging kernel}
\index{Drinfeld associator!arity-three averaging kernel}
Type signature: \textup{(}Open quadrant, ordered tensor/Fulton--
MacPherson presentation, Beilinson level~$2$, hypothesis package:
characteristic zero finite-window ordered bar plus, for the
associator clause, the affine KZ/Drinfeld--Kohno trace-window
package\textup{)}.
Let $W=s^{-1}\bar\cA$ be a finite homogeneous window of the ordered bar
and let
\[
\operatorname{Av}_n=\frac1{n!}\sum_{\sigma\in\Sigma_n}\rho(\sigma)
\]
be the Reynolds projection for the Koszul-signed
\textup{(}or $R$-twisted, when the descent datum is present\textup{)}
$\Sigma_n$-action on $W^{\otimes n}$. Then:
\begin{enumerate}[label=\textup{(\roman*)}]
\item In arity~$2$,
\begin{equation}\label{eq:arity-two-averaging-kernel}
\ker(\operatorname{Av}_2)
=\operatorname{im}(1-\operatorname{Av}_2)
=\operatorname{im}\bigl(1-\rho(s_1)\bigr),
\qquad s_1=(12).
\end{equation}
In the pole-free Koszul-signed specialization this is
\[
\ker(\operatorname{Av}_2)
=
\operatorname{span}\{\,x\otimes y
-(-1)^{|x||y|}y\otimes x\,\},
\]
with $|x|,|y|$ the desuspended degrees. In the \(R\)-twisted case the
generator is instead
\[
x-\rho_R(s_1)x
=x-\tau_{12}R_{12}(z_1-z_2)x .
\]
This is the antisymmetric or \(R\)-skew residue channel discarded by
the symmetric bar.
\item In arity~$3$,
\begin{equation}\label{eq:arity-three-averaging-kernel}
\ker(\operatorname{Av}_3)
=
\operatorname{im}(1-\operatorname{Av}_3)
=
\bigoplus_{\lambda\vdash3,\;\lambda\ne(3)}
(W^{\otimes3})_{[\lambda]},
\end{equation}
the sum of the standard and sign isotypic components. Equivalently it
is generated by differences \(u-\rho(\sigma)u\), and, in characteristic
zero, is the direct sum
\[
(W^{\otimes3})_{[(2,1)]}\oplus(W^{\otimes3})_{[(1,1,1)]}.
\]
\item Under the affine KZ/Drinfeld--Kohno package, the linearised
associator class $[t_{12},t_{23}]$ belongs to the arity-$3$ kernel:
it changes sign under the transposition exchanging the outer labels
and therefore has zero Reynolds average. Its symmetric shadow is not
the Lie commutator but the associative degree-$3$ part of
$\operatorname{Av}_3(\Phi_{\mathrm{KZ}})$.
\item Every antisymmetric residue class is killed by averaging. More
precisely, for an ordered residue tensor
\[
r^-_{12}:=r_{12}-\rho(s_1)r_{12}
\]
one has \(\operatorname{Av}_2(r^-_{12})=0\). More generally,
\(\operatorname{Av}_n\) kills every nontrivial
$\Sigma_n$-isotypic component and retains exactly the trivial
isotypic component, in agreement with
Proposition~\ref{prop:sn-irrep-decomposition-bar}.
\end{enumerate}
\end{proposition}

\begin{proof}
Over characteristic zero, the group algebra of $\Sigma_n$ is
semisimple and the Reynolds operator is the idempotent projector onto
the trivial isotypic component. For \(n=2\), the group has two
elements, so
\[
\operatorname{Av}_2=\tfrac12(1+\rho(s_1)),\qquad
1-\operatorname{Av}_2=\tfrac12(1-\rho(s_1)),
\]
and \(\ker(\operatorname{Av}_2)=\operatorname{im}(1-\rho(s_1))\).
In the pole-free specialization \(\rho(s_1)\) is the Koszul-signed
flip, giving the displayed antisymmetriser. In the \(R\)-twisted case
\(\rho(s_1)=\tau_{12}R_{12}(z_1-z_2)\), giving the displayed
\(R\)-skew generator. This proves~\textup{(i)}.
For \(n=3\), semisimplicity gives the decomposition into the three
partitions $(3)$, $(2,1)$, and $(1,1,1)$; removing the trivial
component gives~\textup{(ii)}.

For~\textup{(iii)}, the Drinfeld--Kohno generators satisfy
$t_{ij}=t_{ji}$ and the KZ associator has linearised commutator
$[t_{12},t_{23}]$ in degree~$3$. The transposition $(1\,3)$ sends this
commutator to $[t_{32},t_{21}]=[t_{23},t_{12}]
=-[t_{12},t_{23}]$. Hence the Reynolds sum pairs this term with its
negative and annihilates it. The remaining symmetric contribution comes
from the associative products in the full KZ associator, as recorded in
Remark~\ref{rem:lie-associative-dichotomy} and
Corollary~\ref{cor:av-kernel-deg3-split}. Part~\textup{(iv)} is the
same Reynolds-projection calculation in arbitrary arity; the displayed
residue tensor \(r^-_{12}\) is in the image of
\(1-\rho(s_1)\), hence has zero Reynolds average by~\textup{(i)}.
\end{proof}

\begin{convention}[Length filtration; \ClaimStatusDefinitional]
\label{conv:bar-length-filtration}
\index{bar complex!length filtration}
The ordered bar carries the decreasing complete word-length filtration
\[
F^rB^{\mathrm{ord}}_X(\cA)
=\prod_{n\ge r}(s^{-1}\bar\cA)^{\otimes n}.
\]
The deconcatenation coproduct satisfies
$\Delta_{\mathrm{dec}}F^r\subset
\sum_{a+b\ge r}F^a\otimes F^b$, the internal and form differentials
preserve~$F^r$, and the residue differential sends
$F^r$ to~$F^{r-1}$. Hence the completed ordered convolution algebra
$\prod_{n\ge0}\Hom((s^{-1}\bar\cA)^{\otimes n},\cA)$ is complete for
the induced arity filtration, and every Maurer--Cartan equation has a
finite arity-$N$ truncation. The symmetric Ran bar inherits the quotient
filtration through $\pi_\Sigma$; it does not remember the filtered
kernel of $\pi_\Sigma$.
\end{convention}

\begin{lemma}[$d_{\mathrm{form}}$ preserves
\ClaimStatusProvedHere
 logarithmic forms]\label{lem:ddr-preserves-log}
The de~Rham differential restricts to a map
$d_{\mathrm{dR}} \colon
 \Omega^q_{\overline{C}_{p+1}}(\log D) \to
 \Omega^{q+1}_{\overline{C}_{p+1}}(\log D)$,
so $d_{\mathrm{form}}$ is well-defined on the bar
complex.
\end{lemma}

\begin{proof}
A logarithmic $q$-form on a smooth variety~$Y$ with
normal-crossings divisor~$D$ is a section of
$\Omega^q_Y(\log D) = \bigwedge^q \Omega^1_Y(\log D)$.
The sheaf $\Omega^1_Y(\log D)$ is locally free,
generated over~$\cO_Y$ by $dz_1/z_1, \ldots, dz_r/z_r,
dz_{r+1}, \ldots, dz_n$ where $D = \{z_1 \cdots z_r = 0\}$.
Since $d_{\mathrm{dR}}(f \cdot dz_i/z_i) = df \wedge
dz_i/z_i$ and $df$ is a section of~$\Omega^1_Y(\log D)$
(regular functions have logarithmic differentials), the
wedge product remains in
$\Omega^{q+1}_Y(\log D)$. This applies with
$Y = \overline{C}_{p+1}(\Sigma_g)$ and
$D = \partial\overline{C}_{p+1}(\Sigma_g)$, which is
normal-crossings by the Fulton--MacPherson construction
(\cite[Theorem~2]{FM94}).
\end{proof}

\begin{remark}[Three components]\label{rem:three-components}
\index{boundary divisor!in bar differential}
Each component has a geometric source and a physical identity:
$d_{\text{internal}}$ is the BRST operator on~$\cA$;
$d_{\text{residue}}$ extracts OPE coefficients via residues
at collision divisors~$D_{ij}$;
$d_{\text{form}}$ is the de~Rham differential on configuration space,
enforcing Ward identities
(Lemma~\ref{lem:ddr-preserves-log}).

Nilpotency $d^2 = 0$ is \emph{not} automatic. Among the nine
cross-terms in $(d_1+d_2+d_3)^2$, the critical residue-residue terms
require the Borcherds identity for OPE modes together with the Arnold
relation for logarithmic forms and the signed residue-orientation
package. Pure mode Jacobi with the form and boundary signs suppressed
does not prove the bar differential square-zero. In the
consecutive-block \(\chirAss\) quotient this is the Stasheff
associativity cancellation for adjacent word-contractions; in the full
FM screen it uses all pairwise collision divisors. See
Proposition~\ref{prop:pole-decomposition},
Proposition~\ref{prop:ordered-bar-local-differential-identities}, and
Theorem~\ref{thm:bar-sign-coherence}. The remaining cross-terms vanish
by Stokes ($\{d_{\text{form}},d_{\text{residue}}\}=0$) and the
derivation property.
\end{remark}

\begin{remark}[Bar differential from the KZ--Arnold superconnection]
\label{rem:dbar-as-kz-pullback}
\index{bar differential!as KZ pullback@as $\KZ^*(\nabla_{\mathrm{Arnold}})$}%
\index{climax theorem!bar differential identification}%
The three-component differential
$d = d_{\mathrm{internal}} + d_{\mathrm{residue}} + d_{\mathrm{form}}$ of
Definition~\ref{def:bar-differential-complete}, restricted to the
genus-$0$ affine/tangent KZ window, is the residue realization of the
bar superconnection obtained by pulling back one universal flat
operator-valued connection along a faithful structure morphism induced
by \(\cA\). This is the structural identity \textup{(}G2\textup{)} of
Chapter~\ref{chap:climax-platonic}:
\[
\nabla^{B,0}_{\cA,V,n}
=d_\cA+d_{\mathrm{dR}}
-\sum_{i<j}\widetilde\rho_{\cA,V,n}(t_{ij})\,\eta_{ij}
=\KZ_{\cA,V,n}^{*}(\nabla_{\mathrm{Arnold}}),
\]
where $\nabla_{\mathrm{Arnold}} = d - \sum_{i<j} t_{ij}\,\eta_{ij}$ is
Arnold's universal flat connection on the trivial
$U(\mathfrak{t}_n)$-bundle over $\Conf_n^{\mathrm{ord}}(X)$
\textup{(}Definition in
\S\ref{ssec:climax-platonic-arnold} of
Chapter~\ref{chap:climax-platonic}\textup{)}, and
the KZ functor of
Construction~\ref{constr:kz-functor-platonic} sends the universal
infinitesimal braid generator $t_{ij}$ to the OPE-derived
$r$-matrix $r^{(ij)}(z_{ij})\cdot z_{ij}$ on
$\End\bigl(B^{\mathrm{ord}}(\cA)\bigr)$. Applying the
Fulton--MacPherson boundary/residue realization to this
superconnection gives the three components of \(d\):
\begin{itemize}
\item the residue component
$d_{\mathrm{residue}} = \sum_{i<j} \mathrm{Res}_{D_{ij}}\,\mu(\phi_i,\phi_j)$
is the chain-level expression of $\KZ^{*}(t_{ij}\eta_{ij})$ at the
collision divisor (\S\ref{ssec:climax-platonic-path-A-fm}, Path~A);
\item the de~Rham component $d_{\mathrm{form}}$ is induced by the
exterior differential in $\nabla_{\mathrm{Arnold}}$;
\item the internal component $d_{\mathrm{internal}}$ is the lift of
$d$ on the coefficient bundle to the family parameter and contributes
no new term beyond the $\cA$-internal differential.
\end{itemize}
The slogan \(d_{\mathrm{bar}}=\KZ^*(\nabla_{\mathrm{Arnold}})\)
means precisely this residue-realized superconnection identity; it is
not a literal equality between a chain differential and a connection.
The Arnold three-term relation~\eqref{eq:arnold-three-term} of
Theorem~\ref{thm:arnold-cohomology-platonic} is precisely
the codimension-two compatibility that forces $d^2 = 0$
\textup{(}Case~(ii) of Theorem~\ref{thm:bar-nilpotency-complete}\textup{)};
the classical Yang--Baxter equation on the family
$\{r^{(ij)}(z_{ij}) = \KZ^{*}(\eta_{ij})\}$ is its $r$-matrix-valued
incarnation. The per-family genus-$0$ affine KZ differentials are not
independent data: each is obtained from this pulled-back
superconnection and its residue realization along the structure
morphism induced by \(\cA \to \End(B^{\mathrm{ord}}(\cA))\). The genus
split and the positive-genus corrections are stated in
Theorem~\ref{thm:kz-arnold-pullback-genus-package}.
\end{remark}

\begin{theorem}[KZ/Arnold pullback package by genus;
\ClaimStatusConditional]
\label{thm:kz-arnold-pullback-genus-package}
\index{KZ connection!bar pullback}
\index{Arnold connection!genus restriction}
\index{infinitesimal braid Lie algebra!bar representation}
Type signature: \textup{(}Open quadrant, ordered logarithmic
configuration presentation, Beilinson level~$2$, hypothesis package:
finite current/KZ trace-window with invariant pairing, faithful
infinitesimal-braid action on the displayed commutator span, and, for
coordinate changes, a conformal stress tensor implementing the Virasoro
action\textup{)}.
Let $V\subset\bar\cA$ be a finite homogeneous current window with
invariant form $\ell_V$ and tensor
$\Omega_{\ell,V}=\sum_{\alpha,\beta}\ell_V(e_\alpha,e_\beta)
e^\alpha\otimes e^\beta$. Let
\[
\cK^{\ord}_{n,V}(\cA)
=
q_n^*j_{n,*}j_n^*((s^{-1}V)^{\boxtimes n})
\otimes_{\cO_{\FM_n(X)}}
\Omega^\bullet_{\FM_n(X)}(\log D_n)
\subset \cK^{\ord}_n(\cA)
\]
be the corresponding finite ordered bar window. For each $n$ define
\begin{equation}\label{eq:kz-arnold-tn-representation}
\rho_{\cA,V,n}\colon\mathfrak t_n\longrightarrow
\End(V^{\otimes n}),
\qquad
\rho_{\cA,V,n}(t_{ij})=\Omega_{\ell,V}^{ij},
\end{equation}
and let
\begin{equation}\label{eq:kz-arnold-bar-suspension-action}
\widetilde\rho_{\cA,V,n}(x)
=(s^{-1})^{\otimes n}\rho_{\cA,V,n}(x)s^{\otimes n}
\end{equation}
be the suspension-conjugate action on \((s^{-1}V)^{\otimes n}\).
Here $\mathfrak t_n$ is the infinitesimal braid Lie algebra generated
by $t_{ij}=t_{ji}$ subject to
\[
[t_{ij},t_{k\ell}]=0\quad
\text{for disjoint pairs},\qquad
[t_{ij},t_{ik}+t_{jk}]=0\quad
\text{for triples}.
\]
Then:
\begin{enumerate}[label=\textup{(\roman*)}]
\item On an affine genus-$0$ tangent screen the KZ pullback identity
acts on \(\cK^{\ord}_{n,V}(\cA)\) by
\[
\nabla^{B,0}_{\cA,V,n}
=d_\cA+d_{\mathrm{dR}}
-\sum_{i<j}\widetilde\rho_{\cA,V,n}(t_{ij})\,\eta_{ij}
=\KZ_{\cA,V,n}^{*}(\nabla_{\mathrm{Arnold}}).
\label{eq:kz-arnold-bar-pullback-identity}
\]
Its singular collision part is the simple-pole/current-window summand
of the ordered bar differential. This is the precise meaning of
$d_B=\KZ^*(\nabla_{\mathrm{Arnold}})$; it is not a global
positive-genus identity.
\item Universally, flatness of
$\nabla_{\mathrm{Arnold}}=d-\sum_{i<j}t_{ij}\eta_{ij}$ is equivalent
to the Arnold three-term relation for the forms together with the
infinitesimal braid relations for the generators. After applying
$\rho_{\cA,V,n}$, the same statement holds modulo
$\ker\rho_{\cA,V,n}$.
\item At genus~$1$ the affine form \(\eta_{ij}=d\log(z_i-z_j)\) is
replaced by the normalized elliptic propagator
\(\eta^{(1)}_{ij}\), and the KZ connection is replaced by the
Bernard--Felder KZB connection
\[
\nabla^{B,1}_{\cA,V,n}
=d-\hbar\sum_{i<j}
\widetilde\rho_{\cA,V,n}(t_{ij})\,
\wp_1(z_{ij},\tau)\,dz_{ij}
-\frac{\hbar}{2\pi i}\sum_iD_i\,d\tau
\]
of Definition~\ref{def:elliptic-kzb-connection}. The cyclic Arnold sum
acquires the period correction
\(\mathcal P_{ijk}(\tau)\) of
Proposition~\ref{prop:elliptic-arnold-relations}, and flatness is the
KZB statement of Proposition~\ref{prop:elliptic-kzb-flatness-jacobi}.
\item For genus \(g\ge2\), the propagator is built from Fay's prime
form \(E(p_i,p_j)\) of
Theorem~\ref{thm:prime-form-properties}; on a fixed fiber the local
singular term is \(d\log E(p_i,p_j)\), while the strict relative
differential is the period/Gauss--Manin corrected operator
\(\Dg{g}=\dfib+d_{\mathrm{per}}^{(g)}
+\nabla_{\mathrm{KS}}^{\mathrm{GM}}\) of
Remark~\ref{rem:period-integrals-bar}. Thus a formula involving only
\(\nabla_{\mathrm{Arnold}}\) is confined to genus~\(0\) or to the
associated graded tangent screen.
\item The local collision operator is coordinate independent: if
$w=f(z)$ is an \'{e}tale coordinate change and \(f'(z_j)\ne0\), then
\begin{equation}\label{eq:kz-coordinate-regular-difference}
d\log(f(z_i)-f(z_j))-d\log(z_i-z_j)
=d\log f'(z_j)
+d\log\!\left(
1+\frac{f''(z_j)}{2f'(z_j)}(z_i-z_j)+O((z_i-z_j)^2)
\right),
\end{equation}
a form regular along \(D_{ij}\). Hence the Poincar\'e residue along
\(D_{ij}\) is unchanged. If the chiral algebra carries a stress tensor
with central charge \(c\), a projective coordinate change is lifted
through the Virasoro action and the defect from tensorial
transformation is the Schwarzian cocycle
\begin{equation}\label{eq:kz-virasoro-coordinate-cocycle}
\widetilde T(f(z))
=(f'(z))^2T(f(z))-\frac{c}{12}\{f;z\},
\qquad
\{f;z\}=\frac{f'''}{f'}-\frac32\left(\frac{f''}{f'}\right)^2 .
\end{equation}
If no stress tensor is part of the structure, only the $\cD$-module
coordinate naturality is asserted.
\item For an \'{e}tale map of curves $u\colon Y\to X$, pullback
commutes with ordered FM collision divisors, logarithmic residues, and
the ordered bar differential:
$u^*\mathbb B^{\ord}_X(\cA)\simeq
\mathbb B^{\ord}_Y(u^*\cA)$. For a smooth base change of curve
families the same statement holds after de~Rham direct image under the
usual holonomic finiteness and proper-support hypotheses.
\end{enumerate}
\end{theorem}

\begin{proof}
The first assertion is the definition of the KZ structure morphism on
the finite current window: the universal generator $t_{ij}$ acts by the
Casimir tensor in the $i,j$ factors, while the universal form is
$\eta_{ij}=d\log(z_i-z_j)$. This gives the displayed pulled-back
connection. Conjugation by \(s^{\otimes n}\) transports this action to
the desuspended coefficient factor of \(\cK^{\ord}_{n,V}(\cA)\). The
residue part of the bar differential is the same collision operator
after applying the form residue; the full bar still contains the
internal differential and all OPE modes recorded in
Proposition~\ref{prop:pole-decomposition}.

For flatness, compute
\[
\nabla_{\mathrm{Arnold}}^2
=-\sum_{i<j}t_{ij}\,d\eta_{ij}
+\frac12\sum [t_{ij},t_{k\ell}]\,\eta_{ij}\wedge\eta_{k\ell}.
\]
Since $d\eta_{ij}=0$, disjoint-pair curvature terms are exactly
$[t_{ij},t_{k\ell}]$, and triple terms are the product of the
infinitesimal braid commutator with the Arnold cyclic sum. Thus the
universal connection is flat precisely when those two families of
relations hold. A representation can kill extra commutators; hence the
fixed-window statement is the corresponding modulo-kernel assertion.

The genus split is Proposition~\ref{prop:arnold-higher-genus},
Proposition~\ref{prop:elliptic-arnold-relations},
Definition~\ref{def:elliptic-kzb-connection},
Proposition~\ref{prop:elliptic-kzb-flatness-jacobi},
Theorem~\ref{thm:prime-form-properties}, and
Remark~\ref{rem:period-integrals-bar}. For the coordinate statement,
Taylor expansion gives
\[
f(z_i)-f(z_j)
=f'(z_j)(z_i-z_j)
\left(1+\frac{f''(z_j)}{2f'(z_j)}(z_i-z_j)
  +O((z_i-z_j)^2)\right).
\]
Taking \(d\log\) gives
\eqref{eq:kz-coordinate-regular-difference}. The final two summands in
that display are regular along \(D_{ij}\), so the logarithmic residue
is the same as the residue of \(d\log(z_i-z_j)\). The Virasoro
coordinate defect is exactly the stress-tensor transformation law
\eqref{eq:schwarzian-transformation} of
Proposition~\ref{prop:schwarzian-central-charge}, written in the same
sign convention in \eqref{eq:kz-virasoro-coordinate-cocycle}. The
\'{e}tale and smooth-base-change statements follow from base change for
logarithmic forms and holonomic $\cD$-module de~Rham direct image.
\end{proof}

\begin{proposition}[Coordinate independence of the collision-residue
bar differential; \ClaimStatusProvedHere]
\label{prop:bar-residue-coordinate-independence}
\index{bar differential!coordinate independence of residue part}
\index{Poincare residue!coordinate independence}
Type signature: \textup{(}Open quadrant, ordered FM/logarithmic
collision-residue presentation, Beilinson level~$2$, hypothesis
package: local affine/formal collision screen, \'{e}tale coordinate
change, logarithmic forms with simple pole along \(D_{ij}\), and the
fixed residue-orientation convention of
Appendix~\ref{app:signs}\textup{)}.
Let \(z\) and \(w=f(z)\) be two \'{e}tale coordinates near a collision
divisor \(D_{ij}\), with \(f'(z_j)\ne0\).  Put \(u=z_i-z_j\) and
\[
v=w_i-w_j=f(z_i)-f(z_j).
\]
Then \(v=u\,g\), where \(g\) is invertible along \(D_{ij}\), and
\[
d\log v=d\log u+d\log g .
\]
The second summand is regular along \(D_{ij}\).  Therefore, for every
logarithmic form \(\alpha\) whose pole along \(D_{ij}\) is at most
simple,
\[
\operatorname{Res}_{D_{ij}}\!\bigl(\alpha\wedge d\log(w_i-w_j)\bigr)
=
\operatorname{Res}_{D_{ij}}\!\bigl(\alpha\wedge d\log(z_i-z_j)\bigr),
\]
with the same determinant-line sign convention.  Consequently the
collision-residue summand \(d_{\mathrm{res}}\) of the ordered bar
differential is independent of the local affine/formal coordinate used
to represent the logarithmic normal form \(\eta_{ij}\).

This proposition concerns the residue part of the bar differential
only.  The transformation law for the full KZ/projective connection is
stronger: when a conformal stress tensor \(T\) with central charge \(c\)
is part of the structure, projective coordinate change contributes the
Schwarzian/Virasoro cocycle of
\eqref{eq:kz-virasoro-coordinate-cocycle}; without that stress-tensor
datum no Virasoro cocycle is asserted.
\end{proposition}

\begin{proof}
Taylor expansion at \(z_j\) gives
\[
f(z_i)-f(z_j)
=
(z_i-z_j)
\left(
 f'(z_j)+\frac{f''(z_j)}{2}(z_i-z_j)+O((z_i-z_j)^2)
\right).
\]
Thus \(v=u\,g\), with
\[
g=f'(z_j)+\frac{f''(z_j)}{2}u+O(u^2),
\qquad
g|_{D_{ij}}=f'(z_j)\ne0 .
\]
Since \(g\) is invertible on the collision screen, \(d\log g\) has no
logarithmic pole along \(u=0\).  In the local decomposition of a
logarithmic form,
\[
\omega=d\log u\wedge \omega_0+\omega_1,
\]
the Poincare residue is \(\omega_0|_{u=0}\); adding a regular
one-form to \(d\log u\) does not change this coefficient.  Hence the
two displayed residues agree, including the determinant-line sign
fixed by the order in which \(d\log u\) is moved to the residue slot.
Applying this equality to the coefficient forms in the ordered bar
construction proves coordinate independence of \(d_{\mathrm{res}}\).
The last paragraph is precisely the separation between residue
naturality and the stress-tensor transformation law already recorded
in Theorem~\ref{thm:kz-arnold-pullback-genus-package}.
\end{proof}

\begin{corollary}[\'{E}tale and holonomic base-change naturality of
the ordered bar differential; \ClaimStatusProvedHere]
\label{cor:ordered-bar-differential-base-change}
\index{bar differential!etale naturality@\'{e}tale naturality}
\index{bar differential!smooth base change}
\index{smooth base change!ordered bar differential}
Type signature: \textup{(}Open quadrant, ordered
Fulton--MacPherson/logarithmic \(\cD\)-module presentation,
Beilinson level~$2$, hypothesis package: augmented complete
finite-window \(\Eone\)-chiral algebra; for \'{e}tale pullback, the
usual pullback of chiral \(\cD\)-modules; for smooth base change, a
smooth proper family of curves, relative FM stages, holonomic finite
bar windows, and proper-support de~Rham direct-image base
change\textup{)}.
Let \(u\colon Y\to X\) be an \'{e}tale map of smooth curves and let
\(u_n\colon\FM_n(Y)\to\FM_n(X)\) be the induced map on ordered
Fulton--MacPherson compactifications. For every augmented finite-window
\(\Eone\)-chiral algebra \(\cA\) on \(X\) there is a natural
identification
\[
\Phi_{u,n}\colon
u_n^*\mathbb B^{\ord}_{X,n}(\cA)
\xrightarrow{\ \sim\ }
\mathbb B^{\ord}_{Y,n}(u^*\cA)
\]
under which the three summands of the ordered bar differential commute:
\[
\Phi_{u,n}\,u_n^*d_\cA=d_{u^*\cA}\,\Phi_{u,n},\qquad
\Phi_{u,n}\,u_n^*d_{\mathrm{dR}}
=d_{\mathrm{dR}}\,\Phi_{u,n},
\]
and
\[
\Phi_{u,n}\,u_n^*d_{\mathrm{res}}^X
=d_{\mathrm{res}}^Y\,\Phi_{u,n}.
\]
Equivalently,
\[
\Phi_{u,n}\,u_n^*d_B^X=d_B^Y\,\Phi_{u,n}.
\]

Now let
\[
\begin{tikzcd}
\mathcal X' \arrow[r,"g"] \arrow[d,"\pi'"] &
\mathcal X \arrow[d,"\pi"]\\
S' \arrow[r,"h"] & S
\end{tikzcd}
\]
be a cartesian smooth base change of smooth proper curve families.
Assume the arity-\(n\) relative ordered bar windows are holonomic
\(\cD\)-complexes on \(\FM_n(\mathcal X/S)\) and have proper support
over \(S\). Then the standard holonomic \(\cD\)-module base-change
isomorphism for de~Rham direct image identifies
\[
h^*\,R(\pi_n)_+
\bigl(\mathbb B^{\ord}_{\mathcal X/S,n}(\cA),d_B\bigr)
\;\cong\;
R(\pi'_n)_+
\bigl(\mathbb B^{\ord}_{\mathcal X'/S',n}(g^*\cA),d_B\bigr),
\]
and this isomorphism intertwines \(d_\cA\), \(d_{\mathrm{dR}}\),
\(d_{\mathrm{res}}\), and hence \(d_B\). Without the holonomic
finite-window and proper-support hypotheses this corollary asserts only
the aritywise pullback identity before relative de~Rham pushforward.
\end{corollary}

\begin{proof}
For an \'{e}tale map \(u\), the formal neighbourhood of a diagonal in
\(Y^n\) is the pullback of the formal neighbourhood of the corresponding
diagonal in \(X^n\). Thus
\[
u_n^{-1}(D^X_{ij})=D^Y_{ij}
\]
and the relative normal parameter on \(Y\) is obtained from the one on
\(X\) by an \'{e}tale coordinate change. Proposition~\ref{prop:bar-residue-coordinate-independence}
therefore gives the residue identity
\[
u_n^*\Res_{D^X_{ij}}=\Res_{D^Y_{ij}}\,u_n^*
\]
on logarithmic forms. The chiral product is a morphism of
\(\cD\)-modules, so \'{e}tale pullback carries
\[
\mu_{\cA}\colon j_*j^*(\cA\boxtimes\cA)\to\Delta_!\cA
\]
to the chiral product \(\mu_{u^*\cA}\). Hence the OPE-mode projections
and the residue operation commute with pullback, proving the displayed
identity for \(d_{\mathrm{res}}\). The identities for \(d_\cA\) and
\(d_{\mathrm{dR}}\) are the ordinary functoriality of the internal
differential and of logarithmic de~Rham forms. Summing the three
identities gives naturality of \(d_B\).

For the family statement, the same argument applies on the relative FM
spaces \(\FM_n(\mathcal X/S)\) before pushforward. Under the stated
holonomic finite-window and proper-support hypotheses, holonomic
\(\cD\)-module de~Rham direct image satisfies smooth/proper base change
\textup{(}Hotta--Takeuchi--Tanisaki~\cite[\S1.7]{HTT08}\textup{)}.
The three components \(d_\cA\), \(d_{\mathrm{dR}}\), and
\(d_{\mathrm{res}}\) are morphisms of the corresponding relative
holonomic complexes, and the base-change isomorphism is functorial for
such morphisms. It therefore carries the relative complex with
differential \(d_B\) to the pulled-back relative complex with the
differential of \(g^*\cA\). This proves the asserted pushed-forward
base-change identity and also explains why the holonomic/proper-support
hypotheses are not optional after de~Rham direct image.
\end{proof}

\begin{remark}[The ordered bar carrier, not
$\mathrm{SC}^{\mathrm{ch,top}}$]
\label{rem:bar-not-sc-coalgebra}
%
\index{Swiss-cheese!emerges on derived centre, not on bar}%
The bar object $B^{\mathrm{ord}}(\cA) = T^c(s^{-1}\bar\cA)$ is a
single ordered coassociative coalgebra with a single degree-$+1$
differential $d_{\mathrm{bar}}$ and the deconcatenation coproduct. Under
the finite locally free \(\chirAss\)-self-duality package of
Corollary~\ref{cor:ordered-bar-chirass-dual-carrier}, this is
equivalently a coalgebra over the chiral cooperad
\((\chirAss)^{!,c}\). It is \emph{not} a Swiss-cheese coalgebra:
the $\mathrm{SC}^{\mathrm{ch,top}}$ structure emerges on
the derived-centre pair
$(C^\bullet_{\mathrm{ch}}(\cA,\cA),\,\cA)$ via the chiral Hochschild
cochain complex (Theorem~H), not on $B^{\mathrm{ord}}(\cA)$ itself.
A statement such as ``the bar complex is an
$\mathrm{SC}^{\mathrm{ch,top}}$-coalgebra'' or ``the bar differential is
the closed-color piece of the Swiss-cheese structure'' is false; the
bar holds the $E_1$ data, the derived centre holds the
$\mathrm{SC}^{\mathrm{ch,top}}$ data
(\S\ref{ssec:climax-platonic-equation-g2} of
Chapter~\ref{chap:climax-platonic} for the $E_1$ pullback identity).
No assertion in this chapter uses an open--closed lift to prove a
property of $B^{\mathrm{ord}}(\cA)$ or $\barB_X(\cA)$.
\end{remark}

\begin{remark}[Curved $A_\infty$ and the genus-$g$ fiberwise curvature]
\label{rem:curved-fiberwise-not-bar-d-squared}
\index{curved A-infinity!fiberwise curvature versus bar differential}%
The bar differential satisfies $d_{\mathrm{bar}}^2 = 0$ unconditionally
in this chapter (Theorem~\ref{thm:bar-nilpotency-complete}). At higher
genus a separate fiberwise differential $d_{\mathrm{fib}}$ on the
Hodge-bundle direction acquires curvature
$d_{\mathrm{fib}}^{\,2} = \kappa(\cA) \cdot \omega_g$
(Chapter~\ref{chap:higher-genus}), and at the curved $A_\infty$
level the relation $\mu_1^2(a) = [\mu_0, a]$ supersedes the strict
$\mu_1^2 = 0$. The two nilpotence statements
do not conflict: the genus-$0$ bar differential and the curved
$A_\infty$ differential live on different sides of the bar
construction. The genus-$g$ fiberwise curvature reads off the
$\kappa$-conductor of $\cA$, which by identity~\textup{(G3)}
$\kappa(\cA) = -\,c_{\mathrm{ghost}}(\BRST(\cA))$ of
Chapter~\ref{chap:universal-conductor} is independent of the choice of
quasi-free BRST resolution.
\end{remark}

\begin{example}[Heisenberg degree-\texorpdfstring{$1$}{1} bar differential]\label{ex:heisenberg-d-deg1}
\label{ex:free-boson-d-deg1}
The Heisenberg algebra $\mathcal{H}_k$ with OPE $J(z)J(w) = k/(z-w)^2 + \text{regular}$ has bar differential
$d_{\mathrm{res}}(J \otimes J \otimes \eta_{12}) = J_{(1)}J = k \cdot \mathbf{1}$
at degree~$1$: the full chiral product extracts the double-pole mode $J_{(1)}J = k$. Since $J_{(0)}J = 0$ (no simple pole), only the curvature component $d_{\mathrm{curvature}}$ contributes (Proposition~\ref{prop:pole-decomposition}). The level~$k$ is visible at genus~$0$ through the bar differential; at genus~$g \geq 1$, it acquires a topological partner $\kappa(\cH_k)\cdot\omega_g$ from the Hodge bundle (Chapter~\ref{chap:higher-genus}). See Chapter~\ref{ch:heisenberg-frame}, \S\ref{sec:frame-bar-deg1} for the full computation.
\end{example}

% Orientation bundle: see Definition~\ref{def:orientation} below.

\subsection{Sign conventions}\label{sec:sign-conventions}

\begin{convention}[Enhanced sign system; \ClaimStatusDefinitional]\label{conv:orientations-enhanced}
The bar complex requires three types of sign conventions:

\emph{Type 1: Koszul signs (algebraic).}

The standard Koszul sign rule (Appendix~\ref{app:signs}, \S\ref{app:sign-conventions}; see also \S\ref{sec:sign-dictionary-complete}) applies with total degree $|\phi| + k$ for a field $\phi \in \mathcal{A}$ tensored with a form in $\Omega^k$.

\emph{Type 2: Orientation signs (geometric).}

Configuration spaces and their boundary divisors carry orientations:

\begin{enumerate}
\item \emph{Configuration space orientation:} $\overline{C}_{n+1}(\Sigma_g)$ is oriented via the complex structure:
\[\text{or}(\overline{C}_{n+1}) = dz_1 \wedge d\bar{z}_1 \wedge \cdots \wedge dz_n \wedge d\bar{z}_n\]
(after modding out by automorphisms; see \S\ref{sec:sign-conventions} below)

\item \emph{Divisor orientation:} Each boundary divisor $D_{ij}$ is oriented by the \emph{outward normal} convention:
\[\text{or}(D_{ij}) = d\epsilon_{ij} \wedge \text{or}(\text{tangent to } D_{ij})\]
where $\epsilon_{ij} = |z_i - z_j|$ points outward (into the interior).

\item \emph{Codimension-2 strata:} At intersections $D_{ij} \cap D_{jk} = D_{ijk}$:
\[\text{or}(D_{ijk}) = d\epsilon_{ij} \wedge d\epsilon_{jk} \wedge \text{or}(\text{tangent})\]

The key identity (Lemma~\ref{lem:orientation}, item~(3)):
\[\text{or}(D_{ijk})|_{D_{ij}} = -\text{or}(D_{ijk})|_{D_{jk}}\]
This sign difference ensures Stokes' theorem holds with correct cancellations.

\item \emph{Residue orientation:} When computing $\text{Res}_{D_{ij}}$, we use:
\[\text{Res}_{D_{ij}}\left(\frac{d\epsilon_{ij}}{\epsilon_{ij}} \wedge \alpha\right) = (+1) \cdot \alpha|_{D_{ij}}\]
(no extra sign for residue extraction)
\end{enumerate}

\emph{Type 3: Operadic signs.}

The bar complex has an operadic structure (composition of operations). When composing two operations, the sign arises from three sources: \emph{grafting trees} (reordering edges when attaching one tree to another), \emph{shuffle signs} (permuting tensor factors to bring colliding fields together), and the \emph{Koszul sign} (moving differential forms past fields).

The formula (for operads): if we compose operations of degree $m$ and $n$ at the $i$-th input:
\[\text{sign} = (-1)^{\epsilon}\]
where:
\[\epsilon = \sum_{j=1}^{i-1} |p_j| \cdot |q|\]
($|p_j|$ are degrees of inputs before position $i$, $|q|$ is degree of the composed operation)

\emph{Compatibility condition.}

These three types of signs must be compatible to ensure $d^2 = 0$. The key relations are:

\begin{enumerate}
\item \emph{Koszul-Orientation compatibility:}
\[\text{sign}_{\text{Koszul}}(\phi_i \leftrightarrow \phi_j) \cdot \text{sign}_{\text{orient}}(D_{ij} \leftrightarrow D_{ji}) = (-1)^{1}\]
(fields anticommute up to orientation sign)

\item \emph{Orientation-Residue compatibility:}
\[\text{Res}_{D_{ij}} \circ \text{Res}_{D_{jk}} + \text{Res}_{D_{jk}} \circ \text{Res}_{D_{ij}} = 0 \quad \text{(with correct signs)}\]
(residues anticommute at codimension-2 strata)

\item \emph{Koszul-Operadic compatibility:}
\[\text{sign}_{\text{Koszul}}(\text{reorder}) = \text{sign}_{\text{operadic}}(\text{compose})\]
(both give the same sign for the same operation)
\end{enumerate}

\end{convention}

\begin{lemma}[Sign compatibility; \ClaimStatusProvedHere]\label{lem:sign-compatibility}
The three types of signs (Koszul, orientation, operadic) are mutually compatible in the sense required for $d^2 = 0$.
\end{lemma}

\begin{proof}
We verify each compatibility relation:

\emph{Relation 1: Koszul-Orientation.}

Consider swapping two fields $\phi_i \otimes \phi_j \to \phi_j \otimes \phi_i$:
- Koszul sign: $(-1)^{|\phi_i| \cdot |\phi_j|}$
- This corresponds to swapping collision divisors $D_{ij} \leftrightarrow D_{ji}$
- Orientation sign: $\text{or}(D_{ji}) = -\text{or}(D_{ij})$ (from antisymmetry of differentials)

The product:
\[(-1)^{|\phi_i| \cdot |\phi_j|} \cdot (-1) = (-1)^{|\phi_i| \cdot |\phi_j| + 1}\]

For bosonic fields ($|\phi_i|, |\phi_j|$ even), this is $(-1)^{0+1} = -1$.
For fermionic fields ($|\phi_i|, |\phi_j|$ odd), this is $(-1)^{1+1} = +1$.

This yields the correct commutation and anticommutation relations for super-objects.

\emph{Relation 2: Orientation-Residue.}

At a codimension-2 stratum $D_{ijk} = D_{ij} \cap D_{jk}$:

Approach from $D_{ij}$ side:
\[\text{or}(D_{ijk})|_{D_{ij}} = d\epsilon_{jk} \wedge \text{or}(D_{ij})\]

Approach from $D_{jk}$ side:
\[\text{or}(D_{ijk})|_{D_{jk}} = d\epsilon_{ij} \wedge \text{or}(D_{jk})\]

By Lemma~\ref{lem:orientation}(3), these differ by a sign: $\text{or}(D_{ijk})|_{D_{ij}} = -\text{or}(D_{ijk})|_{D_{jk}}$.

Now compute double residue:
\begin{align*}
\text{Res}_{D_{ij}} \text{Res}_{D_{jk}}(\omega) + \text{Res}_{D_{jk}} \text{Res}_{D_{ij}}(\omega) &= \int_{D_{ijk}} \omega|_{D_{ijk}} \text{ (from }\text{or}(D_{ijk})|_{D_{ij}}\text{)} \\
&\quad + \int_{D_{ijk}} \omega|_{D_{ijk}} \text{ (from }\text{or}(D_{ijk})|_{D_{jk}}\text{)} \\
&= (+1) \int + (-1) \int = 0
\end{align*}

The orientations differ by exactly the sign needed for cancellation.

\emph{Relation 3: Koszul-Operadic.}

Consider composing two operations $\mu_1: V_1 \otimes V_2 \to W_1$ and $\mu_2: W_1 \otimes V_3 \to W_2$.

Koszul sign for moving $V_2$ past $W_1$:
\[(-1)^{|V_2| \cdot |W_1|}\]

Operadic sign for grafting:
\[(-1)^{\epsilon}\] where $\epsilon = |V_1| + |V_2|$ (degrees of inputs before the graft point)

These match when we account for the suspension in the bar construction ($W_1$ has degree shifted by 1). More precisely, the sign convention on the desuspension $s^{-1}\bar{\mathcal{A}}$ is \emph{chosen} so that the Koszul and operadic signs are compatible; we verify that this choice is consistent.

\emph{General case.}
Let $T$ be the rooted tree labeling the given bar element. Each of
the three sign sources defines a parity character on $\operatorname{Aut}(T)$:
Koszul signs come from permuting tensor factors, orientation signs
come from permuting configuration-space coordinates, and operadic
signs come from grafting odd inputs through the operad composition.

The key identity is checked on an adjacent transposition of two inputs
$\phi_i,\phi_j$. On the desuspension $s^{-1}\bar{\mathcal{A}}$, that
transposition contributes
\[
(-1)^{(|\phi_i|-1)(|\phi_j|-1)}
= (-1)^{|\phi_i||\phi_j|}
 (-1)^{|\phi_i|+|\phi_j|}
 (-1).
\]
The first factor is the Koszul sign, the second is the operadic sign,
and the last factor is the orientation sign from swapping the two
coordinates. Hence the combined action on
$s^{-1}\bar{\mathcal{A}}$ already contains all three sign sources.

Because every tree automorphism is generated by adjacent
transpositions, compatibility for generators implies
\[
\epsilon_{\mathrm{Koszul}} \cdot
\epsilon_{\mathrm{orient}} \cdot
\epsilon_{\mathrm{operad}} = +1
\]
for every rooted tree~$T$. The degree-$2$ computation above supplies
the base case, and induction on tree depth propagates the identity to
all graftings.
\end{proof}

\subsection{\texorpdfstring{Proof that $d^2 = 0$: nine-term verification}{Proof that d2 = 0: nine-term verification}}\label{sec:bar-nilpotency-nine-terms-complete}\label{sec:bar-nilpotency}

\begin{theorem}[Nilpotency of bar differential; \ClaimStatusProvedHere]\label{thm:bar-nilpotency-complete}
\index{bar differential!nilpotence}
\textup{[Regime: quadratic, genus-$0$
\textup{(}Convention~\textup{\ref{conv:regime-tags})}.]}

The differential $d = d_{\text{internal}} + d_{\text{residue}} + d_{\text{form}}$ on the bar complex satisfies:
\[d^2 = 0\]

More precisely, all nine cross-terms arising from $(d_1 + d_2 + d_3)^2$ cancel.
\end{theorem}

\begin{proof}
Write $d = d_1 + d_2 + d_3$ where $d_1 = d_{\text{int}}$ is the internal differential on $\mathcal{A}$, $d_2 = d_{\text{res}}$ extracts residues at collision divisors, and $d_3 = d_{\text{dR}}$ is the de Rham differential on forms. Expanding $(d_1 + d_2 + d_3)^2$, we obtain six terms (three diagonal and three anticommutators) which we verify separately.

\emph{Diagonal terms.} The identity $d_1^2 = 0$ holds because $d_\mathcal{A}^2 = 0$ on $\mathcal{A}$: diagonal terms vanish by hypothesis, and cross-terms (where $d_1$ hits two different tensor factors $\phi_i$ and $\phi_j$) cancel in pairs by the Koszul sign rule. The identity $d_3^2 = 0$ is the standard relation $d_{\text{dR}}^2 = 0$ for differential forms.

The identity $d_2^2 = 0$ is the most substantial. In the double sum
\[d_2^2 = \sum_{i<j}\sum_{k<\ell} (-1)^{\sigma_{ij} + \sigma_{k\ell}'}\, \text{Res}_{D_{k\ell}} \circ \text{Res}_{D_{ij}},\]
three cases arise according to how the index sets $\{i,j\}$ and $\{k,\ell\}$ overlap.

\emph{Case (i): Disjoint pairs} ($\{i,j\} \cap \{k,\ell\} = \emptyset$). The divisors $D_{ij}$ and $D_{k\ell}$ intersect transversally in $\partial\overline{C}_{n+1}(X)$. Since the divisors are disjoint (no shared indices), the residue operations themselves commute:
\[\text{Res}_{D_{ij}} \circ \text{Res}_{D_{k\ell}} = \text{Res}_{D_{k\ell}} \circ \text{Res}_{D_{ij}}.\]
The cancellation comes from the signed sum in $d_2$: the ordered pairs $(i,j),(k,\ell)$ and $(k,\ell),(i,j)$ contribute with opposite signs $(-1)^{\sigma_{ij}+\sigma_{k\ell}'}$ and $(-1)^{\sigma_{k\ell}+\sigma_{ij}'}$ (Lemma~\ref{lem:sign-compatibility}), and these signs differ by $(-1)$, so the two terms cancel.

\emph{Case (ii): One overlap} (say $j = k$).
This is the case that produced the Arnold relation in the
Heisenberg degree-$2$ model computation: three boundary strata meet
at a triple collision, and their contributions cancel cyclically.
We give the complete argument.

\smallskip
\emph{Step (a): Identify the codimension-$2$ locus.}
When $|\{i,j\}\cap\{k,\ell\}|=1$, say $j=k$, the
double residue $\text{Res}_{D_{j\ell}}\circ\text{Res}_{D_{ij}}$
approaches the locus $D_{ij\ell} \subset
\overline{C}_{n+1}(X)$ where all three points $z_i,z_j,z_\ell$
collide simultaneously. In the FM compactification, $D_{ij\ell}$
is a smooth codimension-$2$ stratum obtained by blowing up the
triple diagonal. Exactly three codimension-$1$ boundary
divisors contain $D_{ij\ell}$, corresponding to the three
ways of performing the triple collision in two stages:
\[
D_{ij}\cap D_{j\ell}, \qquad
D_{j\ell}\cap D_{i\ell}, \qquad
D_{i\ell}\cap D_{ij}.
\]
The first collapses $z_i\to z_j$ and then $z_j\to z_\ell$;
the second collapses $z_j\to z_\ell$ and then $z_i\to z_\ell$;
the third collapses $z_i\to z_\ell$ and then $z_j\to z_\ell$.

\smallskip
\emph{Step (b): Extract the three iterated-residue terms.}
Each codimension-$1$ stratum through $D_{ij\ell}$ contributes
one iterated residue to $d_2^2$. The chiral product
$\mu$ determines the algebra element produced at each stage:
\begin{enumerate}[label=\textup{(\arabic*)}]
\item \emph{First $D_{ij}$, then $D_{j\ell}$.}
 Collapsing $z_i\to z_j$ applies
 $\mu(\phi_i,\phi_j)$; collapsing the result with $z_\ell$
 applies $\mu(-,\phi_\ell)$. This gives:
 \[
 \text{Res}_{D_{j\ell}}\circ\text{Res}_{D_{ij}}
 \bigl[\mu\bigl(\mu(\phi_i,\phi_j),\,\phi_\ell\bigr)
 \otimes\omega\bigr].
 \]
\item \emph{First $D_{j\ell}$, then $D_{ij}$.}
 Collapsing $z_j\to z_\ell$ applies
 $\mu(\phi_j,\phi_\ell)$; collapsing with $z_i$ applies
 $\mu(\phi_i,-)$. This gives:
 \[
 \text{Res}_{D_{ij}}\circ\text{Res}_{D_{j\ell}}
 \bigl[\mu\bigl(\phi_i,\,\mu(\phi_j,\phi_\ell)\bigr)
 \otimes\omega\bigr].
 \]
\item \emph{First $D_{i\ell}$, then $D_{ij}$.}
 Collapsing $z_i\to z_\ell$ applies
 $\mu(\phi_i,\phi_\ell)$; collapsing the result with $z_j$
 applies $\mu(-,\phi_j)$. This gives:
 \[
 \text{Res}_{D_{ij}}\circ\text{Res}_{D_{i\ell}}
 \bigl[\mu\bigl(\mu(\phi_i,\phi_\ell),\,\phi_j\bigr)
 \otimes\omega\bigr].
 \]
\end{enumerate}
These are the only terms in $d_2^2$ that access the
codimension-$2$ locus $D_{ij\ell}$. (All other pairs in the
double sum either have disjoint indices or identical indices,
handled by Cases~(i) and~(iii).) The total contribution of
Case~(ii) at the triple $\{i,j,\ell\}$ is thus:
\begin{equation}
\sum_{\textup{cyclic}}
 \text{Res}_{D_{j\ell}}\circ\text{Res}_{D_{ij}}
 \bigl[\mu\bigl(\mu(\phi_i,\phi_j),\,\phi_\ell\bigr)
 \otimes\omega\bigr]
= 0,
\end{equation}
where the cyclic sum runs over the three terms~(1)--(3).

\smallskip
\emph{Step (c): The geometric origin: boundary of boundary.}
The relation $\partial^2=0$ on the cellular chain complex of
$\overline{C}_{n+1}(X)$ (Lemma~\ref{lem:orientation}) implies
that the signed sum of the codimension-$2$ faces arising from
any codimension-$2$ stratum vanishes. Applied to $D_{ij\ell}$:
the three codimension-$1$ strata $D_{ij}$, $D_{j\ell}$,
$D_{i\ell}$ each pair to give a codimension-$2$ face through
$D_{ij\ell}$, and the boundary-of-boundary relation
$\partial^2=0$ asserts that these three faces, with their
orientation signs, sum to zero. Since each face carries the
corresponding iterated residue of
$\mu(\mu(\phi_?,\phi_?),\phi_?)\otimes\omega$, the relation
$\partial^2=0$ translates directly into the displayed cyclic-sum
identity.

\smallskip
\emph{Step (d): Equivalence with the Borcherds identity.}
The displayed vanishing is precisely
axiom~(ii) of Definition~\ref{def:chiral-algebra}: the
associativity of the chiral product $\mu$, i.e., the
Borcherds identity. This identity couples
\emph{all} OPE pole orders simultaneously, unlike the
Jacobi identity, which governs only the simple-pole bracket
$\{-,-\} = (-)_{(0)}(-)$.
The Borcherds identity is both necessary and sufficient for the
vanishing: any meromorphic product $\mu$ satisfying
locality~(i) and producing $d_2^2=0$ must satisfy~(ii), and
conversely any chiral algebra (satisfying~(ii) by definition)
gives $d_2^2=0$ in Case~(ii), because the residue differential
selects exactly the triple-collision coefficient in that Borcherds
identity.

\smallskip
\emph{Step (e): Conclusion.}
By the Borcherds identity (axiom~(ii) of $\cA$), the cyclic
sum displayed above vanishes for every triple
$\{i,j,\ell\}$ and every choice of sections and forms.
Therefore $d_2^2 = 0$ on all Case~(ii) contributions.

\emph{Case (iii): Same pair} ($\{i,j\} = \{k,\ell\}$). The
Poincar\'e residue removes the unique normal logarithmic factor
$d\log(z_i-z_j)$ from the form part and restricts to the divisor
$D_{ij}$. After that restriction there is no second normal direction
for the same divisor, hence
$\text{Res}_{D_{ij}}^2 = 0$.

\emph{Mixed terms.} The anticommutator $\{d_1, d_2\} = 0$
follows because the chiral product~$\mu$ is a morphism
of complexes: as a map of $\cD_X$-modules,
$\mu \colon j_* j^*(\cA \boxtimes \cA) \to \Delta_!(\cA)$
satisfies
\[
d_\cA \circ \mu = \mu \circ (d_\cA \boxtimes \id
 + \id \boxtimes d_\cA),
\]
the Leibniz rule (this is an axiom of a \emph{differential
graded} chiral algebra: the chiral product intertwines the
internal differential). Since $d_{\mathrm{residue}}$ is
built from~$\mu$ via $\Res_{D_{ij}}$, the commutativity
$d_{\mathrm{internal}} \circ \Res_{D_{ij}}
= \Res_{D_{ij}} \circ d_{\mathrm{internal}}$ follows:
$\Res_{D_{ij}}$ is a $\cD_X$-module operation, and
$d_\cA$ is $\cD_X$-linear.

For $\{d_1, d_3\} = 0$: the operators $d_1$ and $d_3$ act on the algebra and form tensor factors respectively. The Koszul sign accompanying $d_3$ includes a factor $(-1)^{\sum_i |\phi_i|}$ from passing the de Rham differential past the algebra elements. After $d_1$ increases one algebra degree by $1$, the sign acquires an extra factor of $(-1)$. Thus $d_3 \circ d_1 = -d_1 \circ d_3$, and the anticommutator vanishes.

For $\{d_2, d_3\} = 0$: this follows from the residue exact sequence for logarithmic forms. On a normal crossing divisor $D = \bigcup D_{ij} \subset \overline{C}_{n+1}(X)$, the residue map fits into the exact sequence $0 \to \Omega^k \to \Omega^k(\log D) \xrightarrow{\text{Res}} \bigoplus_{D_{ij}} \Omega^{k-1}_{D_{ij}}(\log D') \to 0$, and the de Rham differential is compatible with the residue map: $\text{Res}_{D_{ij}} \circ d_{\text{dR}} = d_{\text{dR}}|_{D_{ij}} \circ \text{Res}_{D_{ij}}$. Since $d_2$ is a signed sum of residues and $d_3 = d_{\text{dR}}$, they anticommute (with signs from Convention~\ref{conv:orientations-enhanced}).
\end{proof}

\begin{remark}[Sign-level verification]\label{rem:sign-bridge}
For a complete sign-by-sign verification of $d^2 = 0$ at low degrees
(including the explicit computation for degree-3 elements $[a|b|c]$ in a
graded associative algebra), see Verification~\ref{ver:d2-degree3} and
Proposition~\ref{prop:master-sign} in Appendix~\ref{app:signs}.
\end{remark}

\begin{theorem}[Bar sign theorem; \ClaimStatusProvedHere]
\label{thm:bar-sign-coherence}
\index{bar differential!sign theorem}
\index{bar signs!determinant-line normalization}
\index{Fulton--MacPherson!nested boundary sign}
Type signature: \textup{(}Open quadrant, ordered logarithmic
Fulton--MacPherson presentation, Beilinson level~$2$, hypothesis
package: desuspended ordered bar convention of
Convention~\ref{conv:A-infinity-2-ordered-bar-signs}, enhanced
orientation convention of
Convention~\ref{conv:orientations-enhanced}, and determinant-line
normalization of the canonical sign source in Appendix~\ref{app:signs},
Convention~\ref{conv:canonical-sign-source-app-b}\textup{)}.
All displayed bar-sign conventions in this chapter agree with the
determinant-line signs of Theorem~\ref{thm:det-bar-cobar-signs} and the
master sign formula of Proposition~\ref{prop:master-sign}. In
particular:
\begin{enumerate}[label=\textup{(\roman*)}]
\item the local length-$2$ and length-$3$ formulas of
Computation~\ref{comp:bar-signs-2-1} agree with the geometric residue
signs of Convention~\ref{conv:orientations-enhanced};
\item for $[a|b|c]$ the product part satisfies
\[
d_{\mathrm{res}}^2[a|b|c]
=(-1)^{|b|}\bigl([(ab)c]-[a(bc)]\bigr)=0;
\]
\item for $[a|b|c|d]$, the nested boundary through the triple
collision of the first three inputs gives
\[
d_{\mathrm{res}}^2[a|b|c|d]\big|_{D_{12}\cap D_{123}}
=(-1)^{|b|}\bigl([(ab)c|d]-[a(bc)|d]\bigr)=0,
\]
with the same formula after multiplying by the prefix sign for any
other consecutive triple;
\item the general $n$ case follows by induction on the collision tree:
disjoint codimension-two faces cancel by residue anticommutation, and
nested codimension-two faces reduce to the length-$3$ associativity
calculation with the determinant-line prefix sign.
\end{enumerate}
\end{theorem}

\begin{proof}
Lemma~\ref{lem:LV-sign-comparison} identifies the geometric signs with
the operadic signs through the determinant-line isomorphism, while
Theorem~\ref{thm:det-bar-cobar-signs} and
Proposition~\ref{prop:master-sign} give the canonical desuspended
formula. The length-$3$ computation is
Verification~\ref{ver:d2-degree3}: applying the product differential
twice gives
$(-1)^{|b|}((ab)c-a(bc))$.

For the displayed length-$4$ nested face, first use
\[
d_{\mathrm{res}}[a|b|c|d]
=(-1)^{|a|}[ab|c|d]
 +(-1)^{|a|+|b|-1}[a|bc|d]
 +(-1)^{|a|+|b|+|c|}[a|b|cd].
\]
The component supported on $D_{12}\cap D_{123}$ is obtained from the
first two summands only. Applying $d_{\mathrm{res}}$ again to the first
two slots gives
\[
(-1)^{|a|}(-1)^{|a|+|b|}[(ab)c|d]
 +(-1)^{|a|+|b|-1}(-1)^{|a|}[a(bc)|d],
\]
which is the displayed
$(-1)^{|b|}([(ab)c|d]-[a(bc)|d])$. Associativity kills it. A different
consecutive triple only inserts the determinant-line prefix sign
$(-1)^{\sum_{q<i}(|a_q|-1)}$ before the same calculation.

Every codimension-two FM face is either disjoint or nested. The
disjoint case is residue anticommutation from
Convention~\ref{conv:orientations-enhanced}; the nested case contracts
one internal edge of a collision tree and reduces to the length-$3$
calculation. Induction on the number of internal edges gives the
general $n$ statement.
\end{proof}

\begin{remark}[Boundary mechanism]\label{rem:d-squared-synthesis}\label{rem:geometric-miracle}
The three mechanisms entering the proof of $d^2 = 0$,
the Borcherds identity (algebra), Stokes' theorem (topology), and
residue calculus (analysis), all arise from the boundary
stratification of $\overline{C}_n(X)$. For the Heisenberg algebra,
this reduced to the Arnold relation
(\S\ref{sec:frame-bar-deg2}); in general, the factorization structure
on $\operatorname{Ran}(X)$ forces the compatibility
(Beilinson--Drinfeld~\cite{BD04};
see \S\ref{sec:selection-principle} for the selection principle).
\end{remark}

\begin{convention}[Borcherds mode signs; \ClaimStatusDefinitional]
\label{conv:borcherds-mode-signs}
\index{Borcherds identity!mode signs}
\index{OPE modes!Borcherds convention}
For homogeneous fields \(a,b,c\) we use the super Borcherds convention
\begin{equation}\label{eq:borcherds-mode-signs}
\begin{aligned}
&\sum_{j\ge0}\binom{m}{j}
  (a_{(n+j)}b)_{(m+k-j)}c  \\
&\qquad =
\sum_{j\ge0}(-1)^j\binom{n}{j}
a_{(m+n-j)}\bigl(b_{(k+j)}c\bigr)
-(-1)^n(-1)^{|a||b|}
\sum_{j\ge0}(-1)^j\binom{n}{j}
b_{(n+k-j)}\bigl(a_{(m+j)}c\bigr),
\end{aligned}
\end{equation}
with \(|a|\) denoting parity before desuspension.  After desuspension
the additional bar signs are exactly those of
Convention~\ref{conv:A-infinity-2-ordered-bar-signs}; no second
Borcherds sign convention is used elsewhere in this chapter.
\end{convention}

\begin{proposition}[Pole decomposition of the bar differential; \ClaimStatusProvedHere]\label{prop:pole-decomposition}
\index{bar differential!pole decomposition}
\index{bar differential!bracket component}
\index{bar differential!curvature component}
Use the mode-sign convention of
Convention~\ref{conv:borcherds-mode-signs}.
The residue differential decomposes by OPE mode:
\begin{equation}\label{eq:pole-decomposition}
d_{\mathrm{res}}
=
\sum_{m\ge0} d_{(m)},\qquad
d_{(m)}\text{ applies }a_{(m)}b.
\end{equation}
The summand $d_{(0)}=:d_{\mathrm{bracket}}$ is the simple-pole
bracket part. On quadratic current-type sectors, the only
non-bracket pole mode is the scalar double-pole pairing
$a_{(1)}b=\kappa(a,b)|0\rangle$; in that restricted regime we write
$d_{\mathrm{curvature}}:=d_{(1)}$ and
$d_{\mathrm{res}}=d_{\mathrm{bracket}}+d_{\mathrm{curvature}}$.
For Virasoro, W-algebras, and other higher-pole chiral algebras, the
terms $d_{(m)}$ with $m\ge2$ are part of the residue differential and
must not be discarded.

This decomposition has the following properties:
\begin{enumerate}
\item The full signed sum satisfies $d_{\mathrm{res}}^2=0$ by the
Borcherds identity together with the Arnold boundary relation.
\item A single mode summand need not square to zero. In particular,
the simple-pole Jacobi identity controls only the pure
$d_{(0)}d_{(0)}$ part on the affine Lie screen; it does not cancel
mixed pole-order terms in a general chiral algebra.
\item On a quadratic current-type sector with scalar curvature
$d_{(1)}$, the scalarity of $a_{(1)}b$ gives
$d_{\mathrm{curvature}}^2=0$ and
\[
d_{\mathrm{bracket}}^2
=-\{d_{\mathrm{bracket}},d_{\mathrm{curvature}}\}.
\]
\end{enumerate}

The Lie-bracket Jacobi identity is therefore insufficient for
nilpotency of the chiral bar differential. The full Borcherds identity
for the chiral product~$\mu$, involving all pole orders
simultaneously, is required.
\end{proposition}

\begin{proof}
For a chiral algebra $\cA$ with
\[
a(z)b(w)\sim\sum_{m\ge0} a_{(m)}b\,(z-w)^{-m-1},
\]
let $\operatorname{pr}_m\mu$ denote the $m$-th OPE-mode projection of
the chiral product. If the logarithmic form factor is locally
$d\log(z_i-z_j)\wedge\alpha+\beta$, then
\[
\operatorname{Res}_{D_{ij}}^{\mathrm{form}}(\omega)=\alpha|_{D_{ij}}.
\]
The collision differential is the sum of the maps
\[
d_{(m)}
(\phi_0\otimes\cdots\otimes\phi_p\otimes\omega)
=
\sum_{i<j}(-1)^{\sigma_{ij}}\,
(\phi_i)_{(m)}\phi_j\otimes
\widehat{\phi_i,\phi_j}\otimes
\operatorname{Res}_{D_{ij}}^{\mathrm{form}}(\omega),
\]
with the omitted factors ordered by Convention~\ref{conv:set-notation}.
The Poincar\'e residue is applied only to the logarithmic form factor;
the OPE pole order is selected by $\operatorname{pr}_m\mu$. Thus no
term ever takes a residue of
$(z_i-z_j)^{-m-1}d\log(z_i-z_j)$.

For $m=0$ this is the bracket component
$d_{\mathrm{bracket}}$, applying
$\phi_i{}_{(0)}\phi_j$. For an affine current algebra,
$m=1$ gives the scalar level term
$(J^a)_{(1)}J^b=k\,\delta^{ab}|0\rangle$; for the Heisenberg algebra
it gives $J_{(1)}J=k$. These are the current-type curvature terms.
For Virasoro and W-type algebras, higher $m$ also occur and remain in
the sum~\eqref{eq:pole-decomposition}.

The square of the total residue differential is the one-overlap part
of Theorem~\ref{thm:bar-nilpotency-complete}: at a triple collision,
the three iterated products
\[
\mu(\mu(\phi_i,\phi_j),\phi_\ell),\qquad
\mu(\phi_i,\mu(\phi_j,\phi_\ell)),\qquad
\mu(\mu(\phi_i,\phi_\ell),\phi_j)
\]
are weighted by the three logarithmic two-forms in Arnold's relation.
The Borcherds identity cancels the OPE-mode sum carried by these
three products; Arnold cancels the configuration-form coefficients.
Hence the full $d_{\mathrm{res}}$ squares to zero.

\emph{Quadratic current regime.}
If the only non-bracket pole mode is a scalar $d_{(1)}$, then applying
$d_{(1)}$ replaces the colliding pair by a vacuum scalar. A second
positive mode involving that output vanishes because
$\mathbf{1}_{(m)}a=0$ for $m\ge0$. Thus
$d_{\mathrm{curvature}}^2=0$ in this regime. Expanding the total identity
$d_{\mathrm{res}}^2=0$ gives
\begin{equation}\label{eq:borcherds-pole-mix}
d_{\mathrm{bracket}}^2
+\{d_{\mathrm{bracket}},d_{\mathrm{curvature}}\}
+d_{\mathrm{curvature}}^2=0,
\end{equation}
and the scalarity just proved gives the displayed compensation formula.
\end{proof}

\begin{corollary}[Ordered Arnold--Borcherds residue cancellation;
\ClaimStatusProvedHere]
\label{cor:ordered-arnold-borcherds-residue-cancellation}
\index{bar differential!Arnold--Borcherds cancellation}
\index{residue differential!de Rham compatibility}
\index{OPE modes!higher residues}
Type signature: \textup{(}Open quadrant, ordered logarithmic
Fulton--MacPherson presentation, Beilinson level~$2$, hypothesis
package: augmented dg chiral algebra, the mode signs of
Convention~\ref{conv:borcherds-mode-signs}, and the bar sign package
of Theorem~\ref{thm:bar-sign-coherence}\textup{)}.
In the sign convention used in
Proposition~\ref{prop:ordered-bar-local-differential-identities}, the
mixed de~Rham/residue term and the pure residue square vanish:
\[
d_{\mathrm{dR}}d_{\mathrm{res}}
+d_{\mathrm{res}}d_{\mathrm{dR}}=0,\qquad
d_{\mathrm{res}}^2=0.
\]
The first identity is the total-complex form of the logarithmic
residue-chain-map identity
\[
\Res_{D_{ij}}\circ d_{\mathrm{dR}}
=d_{\mathrm{dR},D_{ij}}\circ\Res_{D_{ij}},
\]
with the anticommutation sign supplied by
\eqref{eq:ordered-bar-ddR}.  For the second identity, the ordered
three-point Arnold relation is, for \(i<j<k\),
\begin{equation}\label{eq:ordered-arnold-three-sign}
\eta_{ij}\wedge\eta_{jk}
-\eta_{ij}\wedge\eta_{ik}
+\eta_{jk}\wedge\eta_{ik}=0,
\end{equation}
which is the cyclic convention of Theorem~\ref{thm:arnold-three}
written with increasing wedge order and
\(\eta_{ab}=\eta_{ba}\). The algebraic coefficient multiplying
\eqref{eq:ordered-arnold-three-sign} is exactly the super Borcherds
identity \eqref{eq:borcherds-mode-signs}.

After projecting to the simple-pole summand \(d_{(0)}\), this becomes
the super-Jacobi identity
\[
(a_{(0)}b)_{(0)}c-a_{(0)}(b_{(0)}c)
+(-1)^{|a||b|}b_{(0)}(a_{(0)}c)=0,
\]
equivalently the Arnold/Jacobi screen of
Theorem~\ref{thm:arnold-jacobi}. For \(m\ge1\), the same residue
operation sees the higher OPE mode \(a_{(m)}b\): the form residue
removes the single logarithmic normal factor, while the projection
\(\operatorname{pr}_m\mu\) selects pole order \(m+1\). Thus second and
higher poles enter \(d_{\mathrm{res}}\) through Borcherds mode
projections, not through iterated form residues.
\end{corollary}

\begin{proof}
The de~Rham statement is the standard residue exact sequence for
normal-crossings logarithmic forms, inserted into the totalization
signs of \eqref{eq:ordered-bar-ddR} and
\eqref{eq:ordered-bar-dres}. This is the mixed
\(\{d_{\mathrm{dR}},d_{\mathrm{res}}\}\) calculation in
Theorem~\ref{thm:bar-nilpotency-complete}.

For \(d_{\mathrm{res}}^2\), disjoint collision divisors cancel because
Poincar\'e residues anticommute. Repeating the same divisor gives zero
because the first residue removes its unique logarithmic normal
factor. The only remaining case is a triple collision. Rewriting the
cyclic Arnold relation of Theorem~\ref{thm:arnold-three} in increasing
wedge order gives \eqref{eq:ordered-arnold-three-sign}; the ordered
bar signs are those of Theorem~\ref{thm:bar-sign-coherence}. The
coefficient of these three logarithmic two-forms is precisely the
Borcherds identity in the mode convention
\eqref{eq:borcherds-mode-signs}, so the triple-overlap contribution
vanishes. Setting \(m=n=k=0\) in
\eqref{eq:borcherds-mode-signs} gives the displayed simple-pole
Jacobi identity. The assertion about higher residues is
Theorem~\ref{thm:bd-ope-residue-full-poles}, equivalently the
modewise decomposition of Proposition~\ref{prop:pole-decomposition}.
\end{proof}

\begin{remark}[Bicomplex obstruction]\label{rem:bicomplex-obstruction}
\index{bar complex!bicomplex structure}
Proposition~\ref{prop:pole-decomposition} shows that the pole-order
summands of $d_{\mathrm{res}}=\sum_{m\ge0}d_{(m)}$ do \emph{not} form a
bicomplex. On the quadratic current regime this reduces to the familiar
two-term obstruction:
\[
d_{\mathrm{bracket}}^2
=-\{d_{\mathrm{bracket}},d_{\mathrm{curvature}}\},
\]
with $d_{\mathrm{curvature}}^2=0$ only because the current-type
curvature is scalar. For Virasoro and W-type algebras the higher-pole
summands $d_{(m)}$, $m\ge2$, are part of the same Borcherds
cancellation and cannot be suppressed.

For \emph{free fields} (Heisenberg, free fermion), the OPE has only a double pole ($a_{(0)}b = 0$), so $d_{\mathrm{bracket}} = 0$ and $d_{\mathrm{res}} = d_{\mathrm{curvature}}$, which squares to zero automatically. This explains the simpler algebraic structure of free-field bar complexes (Corollary~\ref{cor:subexp-free-field}).

For Kac--Moody algebras, the mixed term is controlled by the level-$k$
double-pole OPE coefficient $k \cdot (\,,\,)$, whose scalar projection
is the conductor
\[
\kappa(\widehat{\mathfrak{g}}_k)
= \dim(\mathfrak{g})\cdot(k+h^\vee)/(2h^\vee)
\]
(Chapter~\ref{chap:deformation-theory}).
\end{remark}

\begin{proposition}[Operator-valued collision residue and scalar trace;
\ClaimStatusProvedHere]
\label{prop:operator-valued-collision-residue-trace}
\index{collision residue!operator-valued}
\index{r-matrix@$r$-matrix!operator-valued collision residue}
\index{scalar trace!degree-two collision residue}
Work in the Open quadrant, ordered logarithmic bar presentation,
Beilinson level~$2$. Assume the finite current trace-window package
$\mathsf H_{\mathrm{cur,tr}}$: a finite locally free homogeneous
subquotient $V\subset\bar\cA$ closed under the current OPE modes below,
an even non-degenerate invariant pairing $(\,,\,)_V$ on~$V$, and a
quadratic current OPE
\begin{equation}\label{eq:finite-current-window-ope}
a(z)b(w)
\;\sim\;
\frac{\ell_V(a,b)\mathbf{1}}{(z-w)^2}
\;+\;
\frac{[a,b](w)}{z-w}
\;+\; \textup{regular}
\qquad (a,b\in V),
\end{equation}
where $[a,b]=a_{(0)}b$ is a graded Lie bracket and
\[
\ell_V([a,b],c)=\ell_V(a,[b,c])
\]
with the Koszul sign convention fixed by
Convention~\ref{conv:orientations-enhanced}. Let
$\{e_\alpha\}$ be a homogeneous basis of~$V$ and
$\{e^\alpha\}$ the $(\,,\,)_V$-dual basis. The degree-two collision
residue is the tensor-valued one-form
\begin{equation}\label{eq:operator-valued-collision-r}
r^{\mathrm{coll}}_{\cA,V}(z)
\;:=\;
\Omega_{\ell,V}\,\frac{dz}{z},
\qquad
\Omega_{\ell,V}
\;:=\;
\sum_{\alpha,\beta}\ell_V(e_\alpha,e_\beta)\,
e^\alpha\otimes e^\beta
\;\in\; V\otimes V\subset\bar\cA\otimes\bar\cA .
\end{equation}
Thus $r^{\mathrm{coll}}_{\cA,V}(z)$ is a class in
$\bar\cA\otimes\bar\cA\otimes
\Omega^1_{\Conf^{\mathrm{ord}}_2}(*D_{12})$, not a scalar.
Rank-one formulas such as $k\,dz/z$ are matrix coefficients of
\eqref{eq:operator-valued-collision-r}.

The logarithmic form $\eta_{12}=d\log(z_1-z_2)$ sees the OPE modes in
two different ways. In the bar differential, Proposition~\ref{prop:pole-decomposition}
uses every singular mode: $a_{(m)}b$ contributes to~$d_{(m)}$. In the
polar connection kernel, one $d\log$ has already been absorbed: the
simple-pole bracket mode $m=0$ becomes regular, the double-pole mode
$m=1$ gives the simple pole $dz/z$ in
\eqref{eq:operator-valued-collision-r}, and a pole of order $m+1\ge3$
would give $z^{-m}dz$. The current-window hypothesis
\eqref{eq:finite-current-window-ope} is precisely the regime in which
only the degree-two polar tensor remains.

Under the same hypotheses, $r^{\mathrm{coll}}_{\cA,V}$ satisfies the
classical Yang--Baxter equation on
$\Conf^{\mathrm{ord}}_3(\bC)$, with brackets taken in the graded tensor
algebra and the one-form factors wedged:
\begin{equation}\label{eq:current-window-cybe}
[r^{12}(z_{12}),r^{13}(z_{13})]
+[r^{12}(z_{12}),r^{23}(z_{23})]
+[r^{13}(z_{13}),r^{23}(z_{23})]
=0 .
\end{equation}
The scalar attached to the same degree-two window is the normalized
trace of the homotopy-coinvariant operator:
\begin{equation}\label{eq:kappa-degree-two-trace}
\kappa_{\mathrm{tr},V}(\cA)
\;:=\;
\frac{1}{\operatorname{rk}V}\,
\operatorname{contr}_{(\,,\,)_V}
\!\left(\operatorname{Coeff}_{dz/z}
q_2(K^{\mathrm{coll}}_{\cA,V})\right)
\;=\;
\frac{1}{\operatorname{rk}V}
\sum_\alpha \ell_V(e_\alpha,e^\alpha).
\end{equation}
When the manuscript's family normalization declares this finite window
to be the conductor window, $\kappa(\cA)=\kappa_{\mathrm{tr},V}(\cA)$
in that normalization. Before this trace and normalization, there is no
scalar $\kappa(\cA)$ in the collision kernel.
\end{proposition}

\begin{proof}
The bilinear form $\ell_V$ determines the tensor
$\Omega_{\ell,V}$ by the non-degenerate pairing $(\,,\,)_V$; changing a
homogeneous basis changes the two dual bases by inverse matrices with
the Koszul signs already included in the dual-pairing convention, so
\eqref{eq:operator-valued-collision-r} is basis-independent.

In the local relative coordinate $z=z_1-z_2$, the OPE mode
$a_{(m)}b$ is the coefficient of $z^{-m-1}$. The ordered bar connection
kernel has one logarithmic collision form $d\log z=dz/z$. Passing from
the OPE density to the polar connection kernel removes one pole order:
$z^{-m-1}$ becomes $z^{-m}dz$. Hence the bracket mode $m=0$ is regular,
the double-pole mode $m=1$ is the simple-pole one-form
$z^{-1}dz$, and higher modes would produce higher poles. This is the
connection-kernel version of the mode decomposition in
Proposition~\ref{prop:pole-decomposition}; it is not a second residue
of $z^{-m-1}d\log z$.

Let $x=z_{12}$ and $y=z_{23}$, so $z_{13}=x+y$. Invariance of
$\ell_V$ and the graded Jacobi identity for $[\, ,\,]$ give the graded
infinitesimal braid identities
\[
[\Omega_{\ell,V}^{12},
 \Omega_{\ell,V}^{13}+\Omega_{\ell,V}^{23}]=0,
\qquad
[\Omega_{\ell,V}^{13},
 \Omega_{\ell,V}^{12}+\Omega_{\ell,V}^{23}]=0 .
\]
Writing
$A=[\Omega_{\ell,V}^{12},\Omega_{\ell,V}^{13}]$,
$B=[\Omega_{\ell,V}^{12},\Omega_{\ell,V}^{23}]$, and
$C=[\Omega_{\ell,V}^{13},\Omega_{\ell,V}^{23}]$, these identities give
$B=-A$ and $C=A$. The left side of
\eqref{eq:current-window-cybe} is therefore
\[
A\left(\frac{1}{x(x+y)}-\frac{1}{xy}
        +\frac{1}{(x+y)y}\right)\,dx\wedge dy,
\]
and the parenthesized rational function is zero. This proves the CYBE
on the finite current trace-window surface.

Formula~\eqref{eq:kappa-degree-two-trace} is the contraction of the
$V\otimes V$ coefficient of $q_2(K^{\mathrm{coll}}_{\cA,V})$ against
the same trace pairing.
The contraction of a tensor defined from a bilinear form is invariant
under change of homogeneous dual basis. For a rank-one Heisenberg
window $V=\bC J$ with $(J,J)_V=1$ and $\ell_V(J,J)=k$, the tensor
residue is $k\,J\otimes J\,dz/z$ and the normalized trace is~$k$.
\end{proof}

\begin{corollary}[Averaging preserves the scalar conductor;
\ClaimStatusConditional]
\label{cor:averaging-preserves-scalar-kappa}
\index{averaging map!scalar conductor}
\index{scalar conductor!averaging invariance}
Work under the finite current trace-window hypotheses of
Proposition~\ref{prop:operator-valued-collision-residue-trace} and
assume the family normalization identifies
$\kappa(\cA)=\kappa_{\mathrm{tr},V}(\cA)$ on that window. Then the
ordered-to-symmetric Reynolds projection preserves this scalar:
\[
\operatorname{contr}_{(\,,\,)_V}
\operatorname{Coeff}_{dz/z}
\operatorname{Av}_2
\bigl(r^{\mathrm{coll}}_{\cA,V}(z)\bigr)
=
\operatorname{contr}_{(\,,\,)_V}
\operatorname{Coeff}_{dz/z}
r^{\mathrm{coll}}_{\cA,V}(z).
\]
Equivalently, averaging kills the nontrivial $\Sigma_2$-isotypic
collision data but leaves the trivial scalar trace, hence leaves the
declared scalar $\kappa(\cA)$ unchanged.
\end{corollary}

\begin{proof}
The contraction against the invariant pairing is a
$\Sigma_2$-equivariant trace functional on $V\otimes V$. It vanishes on
the antisymmetric summand
$\ker(\operatorname{Av}_2)$ described in
Proposition~\ref{prop:ordered-bar-low-arity-averaging-kernels} and is
the identity on the trivial summand. Therefore applying
$\operatorname{Av}_2$ before taking the coefficient of $dz/z$ and the
trace contraction gives the same scalar as tracing first. The
hypothesis that $\kappa(\cA)$ is this finite-window trace is exactly the
normalization needed to turn the invariant trace into the manuscript's
scalar conductor.
\end{proof}

\begin{corollary}[Functoriality; \ClaimStatusConditional]\label{cor:bar-functorial}
Since $d^2 = 0$, the bar complex
$(\bar{B}^{\bullet}(\cA), d)$ is a chain complex for every
chiral algebra~$\cA$.
(The tensor-valued collision residue
$r(z)=\operatorname{Res}^{\mathrm{coll}}_{0,2}(\Theta_\cA)$ is the
universal twisting morphism; on a finite current trace-window it is
typed by Proposition~\ref{prop:operator-valued-collision-residue-trace}.
See Theorem~\ref{thm:gz26-commuting-differentials} for its genus-$0$
realization as commuting Hamiltonians.) Morphisms of chiral algebras induce
chain maps (apply~$f$ to algebra factors, leave forms
unchanged); the resulting functor
$\bar{B}^{\mathrm{geom}} \colon
\mathsf{ChirAlg}_X^{\mathrm{aug}} \to
\mathsf{dgCoalg}_X$
is stated and proved in full generality in
Theorem~\ref{thm:bar-functorial-complete}.
\end{corollary}

\begin{proof}
The first sentence is the square-zero statement already proved for the
total bar differential. If
$f \colon \cA \to \cA'$ is a morphism of chiral algebras, applying~$f$
to each algebra factor commutes with the internal differential, with
the collision residues, and with the de~Rham differential, so
$\bar B^{\mathrm{geom}}(f)$ is a chain map. Compatibility with the
coproduct, identities, and composition is checked on the same tensor
factors, yielding the stated functor.
\end{proof}

\subsection{Stokes' theorem on configuration spaces}

\begin{theorem}[Stokes' theorem on configuration spaces; \ClaimStatusProvedHere]\label{thm:stokes-config}
For the Fulton--MacPherson compactification $\overline{C}_{n+1}(\Sigma_g)$ with boundary divisor $D = \bigcup_{i<j} D_{ij}$:

For any $\omega \in \Omega^k(\overline{C}_{n+1}(\Sigma_g))$ (a smooth $k$-form):
\[\int_{\overline{C}_{n+1}(\Sigma_g)} d_{\text{dR}}(\omega) = \sum_{i<j} \epsilon_{ij} \int_{D_{ij}} \omega|_{D_{ij}}\]
where $\epsilon_{ij} = \pm 1$ is the orientation sign.

For logarithmic forms $\omega \in \Omega^k(\log D)$:
\[\int_{\overline{C}_{n+1}} d_{\text{dR}}(\omega) = \sum_{i<j} \epsilon_{ij} \int_{D_{ij}} \text{Res}_{D_{ij}}(\omega)\]
\end{theorem}

\begin{proof}
The configuration space $\overline{C}_{n+1}(\Sigma_g)$ is a \emph{manifold with corners}. The boundary consists of multiple smooth divisors $D_{ij}$ meeting transversely along higher codimension strata.

Stokes' theorem for manifolds with corners (Theorem of Melrose, Mazzeo, et al.) states:
\[\int_M d\omega = \sum_{\text{faces } F} \epsilon_F \int_F \omega|_F\]
where faces are the codimension-1 boundary components.

\emph{Step 1: Identify faces}

The faces of $\overline{C}_{n+1}(\Sigma_g)$ are precisely the divisors $D_{ij}$ for $i < j$.

Codimension: Each $D_{ij}$ has codimension 1 in $\overline{C}_{n+1}$:
\[\dim_{\mathbb{C}} D_{ij} = \dim_{\mathbb{C}} \overline{C}_{n+1} - 1 = (n+1) - 1 = n\]

\emph{Step 2: Orientation of faces}

Each face $D_{ij}$ inherits an orientation from the \emph{outward normal} convention (Convention \ref{conv:orientations-enhanced}):
\[\text{or}(D_{ij}) = d\epsilon_{ij} \wedge \text{or}_{\text{tangent}}\]
where $\epsilon_{ij} = |z_i - z_j|$ increases towards the interior.

The sign $\epsilon_{ij}$ in Stokes' theorem is:
\[\epsilon_{ij} = +1 \quad \text{if } \text{or}(D_{ij}) = \text{outward normal orientation}\]
\[\epsilon_{ij} = -1 \quad \text{if opposite}\]

With our conventions: $\epsilon_{ij} = +1$ for all $i < j$.

\emph{Step 3: Corners}

The divisors $D_{ij}$ and $D_{k\ell}$ (for distinct pairs) intersect along codimension-2 strata:
\[D_{ij} \cap D_{k\ell} = D_{ijk\ell}\]

At these corners, we must verify that contributions from different faces cancel appropriately.

Consider the corner $D_{ijk} = D_{ij} \cap D_{jk}$ (where three points collide). Approaching from different faces:

From $D_{ij}$:
\[\text{contribution} = \int_{D_{ijk}} \omega|_{D_{ij}}|_{D_{ijk}} \cdot \epsilon_{jk|D_{ij}}\]

From $D_{jk}$:
\[\text{contribution} = \int_{D_{ijk}} \omega|_{D_{jk}}|_{D_{ijk}} \cdot \epsilon_{ij|D_{jk}}\]

By Lemma~\ref{lem:orientation} (orientation consistency), these have opposite signs:
\[\epsilon_{jk|D_{ij}} = -\epsilon_{ij|D_{jk}}\]

So the corner contributions cancel.

\emph{Step 4: Apply Stokes' theorem}

With corners handled correctly:
\[\int_{\overline{C}_{n+1}} d_{\text{dR}}(\omega) = \sum_{i<j} \int_{D_{ij}} \omega|_{D_{ij}}\]

For logarithmic forms, $\omega|_{D_{ij}}$ is not well-defined (it has a pole), but $\text{Res}_{D_{ij}}(\omega)$ is:
\[\int_{\overline{C}_{n+1}} d_{\text{dR}}(\omega) = \sum_{i<j} \int_{D_{ij}} \text{Res}_{D_{ij}}(\omega)\]
\end{proof}

\begin{example}[Stokes for three points]\label{ex:stokes-three-points}
Consider $\overline{C}_3(\mathbb{C})$ (three points on the complex plane, compactified).

\emph{Boundary.} $D = D_{12} \cup D_{23} \cup D_{13}$ (three divisors)

\emph{2-form.} $\omega = \eta_{12} \wedge \eta_{23}$ (logarithmic 2-form)

\emph{Differential.}
\begin{align*}
d_{\text{dR}}(\eta_{12} \wedge \eta_{23}) &= d(\eta_{12}) \wedge \eta_{23} - \eta_{12} \wedge d(\eta_{23}) \\
&= 0
\end{align*}
(since $d(\eta_{ij}) = 0$ for logarithmic 1-forms)

\emph{Stokes.}
\[\int_{\overline{C}_3} d_{\text{dR}}(\omega) = 0 = \int_{D_{12}} \text{Res}_{D_{12}}(\omega) + \int_{D_{23}} \text{Res}_{D_{23}}(\omega) + \int_{D_{13}} \text{Res}_{D_{13}}(\omega)\]

\emph{Residues.}
\begin{itemize}
\item $\text{Res}_{D_{12}}(\eta_{12} \wedge \eta_{23}) = \eta_{23}|_{D_{12}}$.
\item $\text{Res}_{D_{23}}(\eta_{12} \wedge \eta_{23}) = -\eta_{12}|_{D_{23}}$ (sign from wedge order).
\item $\text{Res}_{D_{13}}(\eta_{12} \wedge \eta_{23}) = 0$ (no pole along $D_{13}$).
\end{itemize}

Therefore
\[0 = \int_{D_{12}} \eta_{23} - \int_{D_{23}} \eta_{12} + 0\]

This is the boundary-shadow of the Arnold relation: the two displayed
residues are the restrictions of the three-term identity
$\eta_{12}\wedge\eta_{23}
+\eta_{23}\wedge\eta_{31}
+\eta_{31}\wedge\eta_{12}=0$ to the corresponding faces.
\end{example}

\begin{corollary}[Residues anticommute at corners; \ClaimStatusProvedHere]\label{cor:residues-anticommute}
For transverse divisors $D_{ij}$ and $D_{k\ell}$ meeting at a codimension-2 stratum:
\[\text{Res}_{D_{ij}} \text{Res}_{D_{k\ell}} + \text{Res}_{D_{k\ell}} \text{Res}_{D_{ij}} = 0\]
(up to sign)
\end{corollary}

\begin{proof}
Since $D_{ij}$ and $D_{k\ell}$ meet transversally at the codimension-2 stratum $D_{ij} \cap D_{k\ell}$, we may choose local coordinates $(u,v, \ldots)$ near the intersection with $D_{ij} = \{u = 0\}$ and $D_{k\ell} = \{v = 0\}$. For a logarithmic form $\omega$ with at most simple poles along both divisors, the iterated residue $\mathrm{Res}_{v=0}\,\mathrm{Res}_{u=0}\,\omega$ extracts the coefficient of $d\log u \wedge d\log v$, while $\mathrm{Res}_{u=0}\,\mathrm{Res}_{v=0}\,\omega$ extracts the coefficient of $d\log v \wedge d\log u = -d\log u \wedge d\log v$. Hence $\mathrm{Res}_{D_{ij}} \circ \mathrm{Res}_{D_{k\ell}} = -\mathrm{Res}_{D_{k\ell}} \circ \mathrm{Res}_{D_{ij}}$.
\end{proof}

\subsection{Arnold relations: proofs from three perspectives}

\begin{convention}[Set ordering and position notation]\label{conv:set-ordering-arnold}
Throughout this manuscript, we adopt the following conventions for ordered sets:

\begin{enumerate}
\item \emph{Natural Ordering:} For any finite subset $S \subseteq \mathbb{N}$,
we always use the ordering inherited from $\mathbb{N}$:
\[S = \{k_1, k_2, \ldots, k_m\} \quad \text{where} \quad k_1 < k_2 < \cdots < k_m\]

\item \emph{Position Function:} For $k \in S$, we denote by $|k|_S$ (or simply $|k|$
when $S$ is clear from context) the \emph{position} of $k$ in this ordering:
\[k = k_{|k|} \quad \iff \quad |k| = i \text{ where } k \text{ is the } i\text{-th smallest element of } S\]

\item \emph{Sign Convention:} Signs arising from reordering are computed via the
Koszul rule. Moving an element $k$ past position $|k|$ introduces sign $(-1)^{|k|-1}$.

\item \emph{Multi-indices:} For multi-index sets (e.g., in partitions), we use
lexicographic ordering.
\end{enumerate}

\emph{Example.} For $S = \{2, 5, 7\}$:
\begin{itemize}
\item $|2|_S = 1$ (first position)
\item $|5|_S = 2$ (second position)
\item $|7|_S = 3$ (third position)
\end{itemize}

In Arnold relations, the notation $(-1)^{|k|}$ means $(-1)^{|k|_S}$ where $S$ is the
index set of the collision divisor under consideration.
\end{convention}

\begin{theorem}[Arnold relations; \ClaimStatusProvedHere]\label{thm:arnold-three}
\index{Arnold relations|textbf}
We encountered these relations concretely in the Heisenberg degree-$2$
bar model, where they appeared as the identity forcing $d^2 = 0$.
In full generality:

For distinct indices $i, j, k \in \{1, \ldots, n\}$, the logarithmic 1-forms $\eta_{ij} = d\log(z_i - z_j)$ satisfy:

\emph{Formulation 1 (basic).}
\[\eta_{ij} \wedge \eta_{jk} + \eta_{jk} \wedge \eta_{ki} + \eta_{ki} \wedge \eta_{ij} = 0\]

\emph{Formulation 2 (Orlik--Solomon).}
The algebra $H^*(C_n(\mathbb{C}); \mathbb{Q})$ is the quotient of the exterior algebra $\bigwedge\langle \eta_{ij} \mid 1 \le i < j \le n \rangle$ by the ideal generated by the three-term Arnold relations of Formulation~1. All higher-degree relations in the Orlik--Solomon algebra follow from these by repeated application; there is no independent ``general Arnold relation'' beyond the three-term version.

\emph{Formulation 3 (logarithmic FM model).}
On the Fulton--MacPherson compactification, the logarithmic de~Rham
complex $\Omega^\bullet_{\overline{C}_n(X)}(\log D)$ restricts to the
open configuration space $C_n(X)$ and computes its cohomology. On
the type-$A$ affine chart $X=\mathbb{A}^1$ this is the
Orlik--Solomon algebra of Formulation~2. For a general curve~$X$,
additional classes pulled from $H^\bullet(X)$ may occur; the Arnold
three-term relation is the universal diagonal-collision relation,
not a presentation of the whole compactified cohomology ring.
\end{theorem}

\begin{proof}
See Theorem~\ref{thm:arnold-relations} for the partial-fractions proof.
The topological, combinatorial, and operadic perspectives give
parallel proofs of the same three-term relation and introduce no
additional hypothesis.
\end{proof}

\begin{remark}[Arnold relation and the Jacobi screen]\label{rem:arnold-jacobi-geometric}
On the simple-pole affine screen, the geometric content of
$\dzero^2=0$ at degree~$3$ is the Jacobi identity: the three boundary
strata of $\overline{C}_4(\mathbb{P}^1)$ that contribute to
$\dzero^2$ carry the three terms of the chiral bracket Jacobi sum.
Concretely, the Arnold relation
$\eta_{12} \wedge \eta_{23} + \eta_{23} \wedge \eta_{31}
+ \eta_{31} \wedge \eta_{12} = 0$ supplies the geometric coefficient of
\[
{f^{ab}}_d\,{f^{dc}}_e + {f^{bc}}_d\,{f^{da}}_e + {f^{ca}}_d\,{f^{db}}_e = 0,
\]
the Jacobi identity for the simple-pole structure constants. This is
the content of Theorem~\ref{thm:arnold-jacobi}: Arnold is equivalent
to Jacobi after projecting the chiral product to its simple-pole
bracket. For the full bar differential, higher OPE modes remain and
the required algebraic identity is Borcherds, as in
Theorem~\ref{thm:bar-nilpotency-complete}. See
Computation~\ref{comp:geom-alg-comparison-deg3} for the
term-by-term affine $\widehat{\mathfrak{sl}}_{2,k}$ verification.
\end{remark}

\begin{corollary}[Cohomology of configuration spaces {\cite{Arnold69}}; \ClaimStatusProvedElsewhere]\label{cor:cohomology-config}
The cohomology ring $H^*(C_n(\mathbb{C}); \mathbb{Q})$ of the open configuration space is:
\[H^*(C_n(\mathbb{C})) \cong \mathbb{Q}[\eta_{ij} : 1 \leq i < j \leq n] / \mathcal{I}_{\text{Arnold}}\]
where $\mathcal{I}_{\text{Arnold}}$ is the ideal generated by Arnold relations.
\end{corollary}

\begin{proof}
By Arnol'd~\cite{Arnold69}, the 1-forms $\eta_{ij} = d\log(z_i - z_j)$ generate $H^*(C_n(\mathbb{C}); \mathbb{Q})$ and satisfy the Arnold relations $\eta_{ij} \wedge \eta_{jk} + \eta_{jk} \wedge \eta_{ki} + \eta_{ki} \wedge \eta_{ij} = 0$. That these generate \emph{all} relations follows from Cohen's computation of $H^*(C_n(\mathbb{C}))$ as the cohomology of the Arnol'd--Brieskorn braid arrangement complement~\cite{Coh76}. Appendix~\ref{app:arnold-relations} records the three proofs of the
Arnold relations (topological, combinatorial, operadic) used
elsewhere in the manuscript.
\end{proof}

\subsection{Low-degree explicit computations}

\subsubsection{Degree 0: the vacuum}

\begin{computation}[Degree 0; \ClaimStatusProvedHere]\label{comp:deg0}
In degree 0:
\[\bar{B}^0(\mathcal{A}) = \Gamma\left(\overline{C}_1(\Sigma_g), \mathcal{A} \otimes \Omega^0(\log D)\right)\]

But $\overline{C}_1(\Sigma_g) = \Sigma_g$ (one marked point, so no collision divisor), and $\Omega^0(\log D) = \mathcal{O}_{\Sigma_g}$ (functions).

Therefore
\[\bar{B}^0(\mathcal{A}) = \Gamma(\Sigma_g, \mathcal{A}) = H^0(\Sigma_g, \mathcal{A})\]

This is the space of global sections of the chiral algebra.

\emph{Physical interpretation.} This is the vacuum sector: states with no operator insertions.

\emph{Differential.} $d: \bar{B}^0 \to \bar{B}^{-1}$. But there is no $\bar{B}^{-1}$ (negative degree), so $d|_{\bar{B}^0} = 0$.
\end{computation}

\begin{proof}
The definition of the reduced geometric bar complex at bar degree~$0$
uses configurations of one marked point, so the configuration space is
\(\overline{C}_1(\Sigma_g)=\Sigma_g\) and there is no boundary divisor.
Hence the logarithmic form factor is \(\Omega^0(\log D)=\mathcal{O}_{\Sigma_g}\),
which gives the displayed identification
\(\bar B^0(\mathcal A)=\Gamma(\Sigma_g,\mathcal A)\).
The bar differential lowers bar degree by one, so its restriction to
\(\bar B^0\) is zero because \(\bar B^{-1}=0\).
\end{proof}

\subsubsection{Degree 1: two-point functions}

\begin{computation}[Degree 1; \ClaimStatusProvedHere]\label{comp:deg1-general}
In degree 1:
\[\bar{B}^1(\mathcal{A}) = \Gamma\left(\overline{C}_2(\Sigma_g), j_*j^*(\mathcal{A} \boxtimes \mathcal{A}) \otimes \Omega^1(\log D_{12})\right)\]

\emph{Configuration space.} $\overline{C}_2(\Sigma_g)$ parametrizes two points on $\Sigma_g$.
- At genus 0: After modding out $\text{PSL}_2$, $\overline{C}_2(\mathbb{P}^1) \cong \mathbb{P}^1$
- At genus $g \geq 1$: $\overline{C}_2(\Sigma_g)$ is more complex (includes period matrix data)

\emph{Logarithmic 1-forms.} $\Omega^1(\log D_{12})$ is 1-dimensional, spanned by:
\[\eta_{12} = \frac{dz_1 - dz_2}{z_1 - z_2} = d\log(z_1 - z_2)\]

\emph{Basis.} After choosing generators \(\{\phi_i\}\) for the algebra factor,
the degree-$1$ term is spanned by:
\[\{\phi_i(z_1) \otimes \phi_j(z_2) \otimes \eta_{12} : 1 \le i,j \le N\}\]

If $\mathcal{A}$ has $N$ generators, then:
there are \(N^2\) basic tensors of this form.

\emph{Differential.} $d: \bar{B}^1 \to \bar{B}^0$
\[d(\phi_i \otimes \phi_j \otimes \eta_{12}) = \text{Res}_{D_{12}}[\mu(\phi_i, \phi_j) \otimes \eta_{12}]\]

where $\mu$ is the chiral product (OPE).

If the OPE is:
\[\phi_i(z)\phi_j(w) = \sum_k \frac{C_{ij}^k}{(z-w)^{\Delta_k}} \phi_k(w) + \text{regular}\]

then the bar differential acts by extracting the full chiral
product~$\mu$ (Convention~\ref{conv:product-vs-bracket}):
\[d_{\mathrm{res}}(\phi_i \otimes \phi_j \otimes \eta_{12})
 = \sum_{m \geq 0} (\phi_i)_{(m)}\phi_j.\]
The mode decomposition of Proposition~\ref{prop:pole-decomposition}
separates the contributions:
\begin{itemize}
\item pole order $1$: $d_{(0)}=d_{\mathrm{bracket}}$ applies
 $(\phi_i)_{(0)}\phi_j$;
\item pole order $2$: $d_{(1)}$ applies
 $(\phi_i)_{(1)}\phi_j$, the scalar curvature pairing in
 current-type examples;
\item pole order $m+1\ge3$: $d_{(m)}$ records the higher OPE modes;
\item regular terms give no bar-residue contribution.
\end{itemize}
The OPE pole is not multiplied by the propagator
$\eta_{12}=d\epsilon/\epsilon$ to form a second residue. The
logarithmic form records the collision direction; the Poincar\'e
residue removes that form factor, while the chiral product supplies
the OPE coefficient.
\end{computation}

\begin{proof}
By definition,
\[
\bar B^1(\mathcal A)
=
\Gamma\!\left(\overline C_2(\Sigma_g),
j_*j^*(\mathcal A\boxtimes\mathcal A)\otimes\Omega^1(\log D_{12})\right).
\]
Choosing generators \(\{\phi_i\}\) for the algebra factor gives the
displayed spanning tensors, and the logarithmic degree-$1$ factor is
generated locally by the standard form
\(\eta_{12}=d\log(z_1-z_2)\).
The bar differential in degree~$1$ is the residue map along the collision
divisor \(D_{12}\), so applying it to a generator tensor produces the
chiral product \(\mu(\phi_i,\phi_j)\), which is exactly the displayed
formula for \(d_{\mathrm{res}}\). The final warning is the usual
separation between the OPE pole and the logarithmic propagator:
\(\eta_{12}\) supplies the normal logarithmic form, the Poincar\'e
residue removes that form, and the singular part of
\(\mu(\phi_i,\phi_j)\) is selected by the OPE-mode projection.
\end{proof}

\begin{remark}[Iterated residues as $A_\infty$ operations]
\label{rem:iterated-residues-ainfty}
\index{Ainfty@$A_\infty$!from FM boundary strata}
\index{Gaiotto--Kulp--Wu!higher operations}
The iterated residues on FM boundary strata are the $A_\infty$
operations~$m_k$. The $k$-th boundary stratum of
$\overline{\operatorname{Conf}}_{k+1}(\mathbb{C})$ contributes to~$m_k$
via the HPL transfer
(Proposition~\ref{prop:ainfty-formality-implies-koszul}).
The Arnold relations (the sum of iterated residues at a
codimension-$2$ stratum vanishes) are the Stasheff $A_\infty$
identity $\sum_{i+j=n+1} m_i \circ m_j = 0$.
This is the geometric content of $d^2 = 0$: the Borcherds identity
at the algebraic level and the codimension-$2$ boundary cancellation
at the geometric level produce the same quadratic relations that
Gaiotto--Kulp--Wu~\cite{GKW2025} derive from Wess--Zumino
consistency of higher operations in the holomorphic-topological
framework.
\end{remark}

\begin{example}[Heisenberg at degree 1]\label{ex:heisenberg-deg1-complete}
The Heisenberg OPE $J(z)J(w) \sim k/(z{-}w)^2$ has no simple pole
($J_{(0)}J = 0$); the entire bar differential at degree~$1$ comes
from the double-pole extraction $J_{(1)}J = k$:
\[
d_{\mathrm{res}}(J \otimes J \otimes \eta_{12}) = k \cdot 1
\quad\in\; \bar{B}^0 = \mathbb{C}.
\]
The level~$k$ is visible at genus~$0$ through
$d_{\mathrm{curvature}}$
(Proposition~\ref{prop:pole-decomposition});
at genus~$\geq 1$ it acquires a topological component
$m_0 = \kappa \cdot \omega_g$
(Chapter~\ref{chap:higher-genus}).
\end{example}

\begin{example}[\texorpdfstring{$\beta\gamma$}{beta-gamma} system at degree 1]\label{ex:betagamma-deg1}
For the $\beta\gamma$ system with generators $\beta(z), \gamma(z)$ and OPE:
\[\beta(z)\gamma(w) = \frac{1}{z-w} + \text{regular}, \quad \beta(z)\beta(w) = 0, \quad \gamma(z)\gamma(w) = 0\]

\emph{Bar degree 1.}
\[\bar{B}^1(\mathcal{FG}) = \text{span}\{\beta \otimes \beta \otimes \eta, \beta \otimes \gamma \otimes \eta, \gamma \otimes \beta \otimes \eta, \gamma \otimes \gamma \otimes \eta\}\]

\emph{Differential.} Only the $\beta \otimes \gamma$ term contributes:
\begin{align*}
d(\beta \otimes \gamma \otimes \eta_{12}) &= \text{Res}_{D_{12}}\bigl[\mu(\beta, \gamma) \otimes \eta_{12}\bigr] \\
&= \bigl[\text{Res}_{\epsilon=0}\,\eta_{12}\bigr] \cdot \mu(\beta, \gamma)\big|_{z_1=z_2}
\end{align*}
Setting $\epsilon = z_1 - z_2$, the propagator is $\eta_{12} = d\epsilon/\epsilon$, which has $\text{Res}_{\epsilon=0}[d\epsilon/\epsilon] = 1$. The chiral product $\mu(\beta, \gamma)$ extracts the simple-pole OPE coefficient $\beta(z)\gamma(w) \sim 1/(z-w)$, giving $\mathbb{1}$ (the unit element). Hence $d(\beta \otimes \gamma \otimes \eta_{12}) = \mathbb{1}$.

(Note: the residue is of the \emph{logarithmic} 1-form $d\epsilon/\epsilon$, not of $d\epsilon/\epsilon^2$; one must not multiply the OPE pole by the propagator pole.)

Similarly: $d(\gamma \otimes \beta \otimes \eta_{12}) = -\mathbb{1}$ (the sign arises from the OPE $\gamma(z)\beta(w) \sim -1/(z-w)$, which carries a sign relative to $\beta(z)\gamma(w) \sim +1/(z-w)$ from the ordering convention for the chiral product, not from any fermionic anticommutativity; the $\beta\gamma$ system is bosonic).

\emph{Cohomology.} $H^1(\bar{B}^{\bullet}(\mathcal{FG})) = \text{span}\{\beta \otimes \beta, \gamma \otimes \gamma\}$ (2-dimensional).
\end{example}

\begin{remark}[Particle interpretation of low bar cohomology]
\label{rem:v1-bar-cohomology-particle-scattering}
For the reduced ordered bar complex
\[
B^{\mathrm{ord}}(\cA) = T^c(s^{-1}\bar\cA),
\qquad
\bar\cA = \ker(\varepsilon),
\]
the low-degree generators-and-relations truncation admits a
particle/scattering interpretation on the quadratic-dual side:
\[
\dim H^1(B^{\mathrm{ord}}(\cA))
= \text{number of single-particle species},
\]
\[
\dim H^2(B^{\mathrm{ord}}(\cA))
= \text{number of independent scattering channels},
\]
where the second count records the independent quadratic relations.
For the standard families:
\begin{itemize}
\item Heisenberg: $\dim H^1 = 1$, corresponding to one boson, and
 $\dim H^2 = 0$, reflecting the absence of quadratic interaction
 channels in the free theory.
\item $\widehat{\mathfrak{sl}}_{2,k}$: $\dim H^1 = 3 = \dim \mathfrak{sl}_2$
 and $\dim H^2 = 5$. The degree-$2$ count is $5$, not $6$.
\item $\widehat{\mathfrak{sl}}_{3,k}$: $\dim H^1 = 8 = \dim \mathfrak{sl}_3$
 and $\dim H^2 = 36$. This is the chiral bar value; the
 Chevalley--Eilenberg count gives $20$ and does not include the
 extra chiral Orlik--Solomon contribution.
\end{itemize}
\end{remark}

\begin{remark}[Fixed-curve configuration geometry across genera]
For a smooth curve $\Sigma_g$, chiral operator insertions are organised by the
genus expansion
\[
\langle \phi_1(z_1) \cdots \phi_n(z_n) \rangle = \sum_{g=0}^{\infty} \lambda^{2g-2} \int_{\text{field space}} \mathcal{D}\phi \, e^{-S[\phi]} \phi_1(z_1) \cdots \phi_n(z_n)
\]
The singularities as $z_i \to z_j$ encode the operator algebra structure at
that fixed genus. The configuration space
$C_n(\Sigma_g) = \Sigma_g^n \setminus \{\text{diagonals}\}$ parametrizes
non-colliding insertion points. Its Fulton--MacPherson compactification
$\overline{C}_n(\Sigma_g)$ adds only collision strata and tangent screens.
No curve degeneration occurs for a fixed $\Sigma_g$; stable-curve
degeneration enters only in the relative compactification over
$\overline{\mathcal M}_{g,n}$. Logarithmic forms
$d\log(z_i-z_j)$ carry the local OPE singularities, while the dependence on
the genus-$g$ complex structure enters through the prime form and period
matrix. The bar differential is the residue operator along collision
faces; modular corrections arise from the variation of these residue kernels
with the period matrix.
\end{remark}

\begin{definition}[Orientation line bundle]\label{def:orientation}
The \emph{orientation line bundle} $\text{or}_{p+1}^{(g)}$ on $\overline{C}_{p+1}(\Sigma_g)$ is
\[
\text{or}_{p+1}^{(g)} = \det(T\overline{C}_{p+1}(\Sigma_g)) \otimes \text{sgn}_{p+1} \otimes \mathcal{L}_g
\]
where:
\begin{itemize}
\item $\det(T\overline{C}_{p+1}(\Sigma_g))$ is the top exterior power of the tangent bundle
\item $\text{sgn}_{p+1}$ is the sign representation of $\mathfrak{S}_{p+1}$
\item $\mathcal{L}_g$ is the orientation line from the period-matrix
variation; on a fixed curve it is a constant line
\item The tensor product ensures that exchanging two points introduces a sign and modular covariance
\end{itemize}
With the Arnold--Orlik--Solomon relation on each collision screen, these
orientation lines give the Stokes signs in the cancellation proving
$d_{\mathrm{bar}}^2=0$.
\end{definition}

\subsection{Explicit low-degree terms}

\begin{example}[Unreduced bar complex in low degrees: orientation only]\label{ex:bar-unreduced-low-degrees}
The display in this example uses the full algebra~$\mathcal{A}$ rather
than the augmentation ideal~$\bar{\mathcal{A}}$. It is the
\emph{unreduced} bar complex. The canonical theorem-level complex is
the reduced ordered bar
$B^{\mathrm{ord}}(\mathcal{A}) = T^c(s^{-1}\bar{\mathcal{A}})$,
and its symmetric Ran quotient is $\barB_X(\mathcal{A})$.

In the unreduced presentation:
\begin{align}
\bar{B}^0_{\mathrm{unred}}(\mathcal{A}) &= \mathcal{A} \\
\bar{B}^1_{\mathrm{unred}}(\mathcal{A}) &= \Gamma(C_2(X), \mathcal{A} \boxtimes \mathcal{A} \otimes \eta_{12}) \\
\bar{B}^2_{\mathrm{unred}}(\mathcal{A}) &= \Gamma(\overline{C}_3(X), \mathcal{A}^{\boxtimes 3} \otimes \Omega^2_{\log}).
\end{align}

Bar degree~$n$ uses sections over $C_{n+1}(X)$, and the bar
differential lowers bar degree. The reduced ordered complex used from
\S\ref{subsec:bar-functoriality} onward replaces $\mathcal{A}$ by its
augmentation ideal
\[
\bar{\mathcal A}=\ker(\varepsilon\colon\mathcal A\to\omega_X)
\]
and starts at $B^{\mathrm{ord},0}=\mathbb C$. The augmentation separates scalar
vacuum expectation values from genuine generators. Thus
$T^c(s^{-1}\bar{\mathcal A})$ tensors only non-scalar modes, carries
the reduced coproduct, and is conilpotent. This conilpotent cofree
coalgebra is the input of the bar-cobar adjunction
\textup{(}Theorem~\textup{\ref{thm:bar-cobar-adjunction}}\textup{)}.
The reduced complex is the theorem-level object; the unreduced display
records the orientation convention for the first residue.
\begin{align}
d: B^{\mathrm{ord},1} &\to B^{\mathrm{ord},0} \\
a_1 \otimes a_2 \otimes \eta_{12} &\mapsto \text{Res}_{z_1=z_2}[a_1(z_1) \cdot a_2(z_2) \cdot \eta_{12}]
\end{align}
This residue extraction computes the OPE product by composing the
Poincar\'e residue of $\eta_{12}$ with the chiral mode projection.
\end{example}

\subsection{Functoriality of the bar construction}
\label{subsec:bar-functoriality}
\index{bar-cobar!functoriality}

\begin{theorem}[Bar construction is functorial; \ClaimStatusProvedHere]\label{thm:bar-functorial-complete}
\label{thm:bar-functorial}
The geometric bar construction defines a functor:
\[\bar{B}^{\text{geom}}: \mathsf{ChirAlg}_X^{\mathrm{aug}} \to \mathsf{dgCoalg}_X\]
from augmented chiral algebras to dg coalgebras, that is:
\begin{enumerate}
\item \emph{Well-defined on objects:} For each augmented chiral algebra $\mathcal{A}$, $\bar{B}^{\text{geom}}(\mathcal{A})$ is a differential graded coalgebra
\item \emph{Well-defined on morphisms:} For each morphism $f: \mathcal{A} \to \mathcal{B}$ of augmented chiral algebras, there is an induced coalgebra morphism $\bar{B}^{\text{geom}}(f): \bar{B}^{\text{geom}}(\mathcal{A}) \to \bar{B}^{\text{geom}}(\mathcal{B})$
\item \emph{Preserves identities:} $\bar{B}^{\text{geom}}(\text{id}_\mathcal{A}) = \text{id}_{\bar{B}^{\text{geom}}(\mathcal{A})}$
\item \emph{Preserves composition:} $\bar{B}^{\text{geom}}(g \circ f) = \bar{B}^{\text{geom}}(g) \circ \bar{B}^{\text{geom}}(f)$
\end{enumerate}
\end{theorem}

\begin{proof}
Well-definedness on objects is
Theorem~\ref{thm:bar-nilpotency-complete}.
For a morphism $f\colon \cA \to \cB$ of augmented chiral
algebras, define
$\bar{B}^{\mathrm{geom}}(f)(a_0 \otimes \cdots \otimes a_n
\otimes \omega) = f(a_0) \otimes \cdots \otimes f(a_n) \otimes
\omega$\label{thm:bar-induced-map}
(apply $f$ to algebra factors, leave forms unchanged).

\emph{Chain map property.}\label{thm:bar-induced-chain-map}
Since $f$ commutes with the internal differential ($f$ is a
$\mathcal{D}$-module map), with the residue differential ($f$
preserves the chiral product), and does not affect the de~Rham
component, $\bar{B}^{\mathrm{geom}}(f)$ commutes with
$d = d_{\mathrm{int}} + d_{\mathrm{res}} + d_{\mathrm{dR}}$.

\emph{Coalgebra morphism.}\label{lem:bar-induced-coalgebra}
The coproduct $\Delta$ is defined by restricting to collision
divisors; since $f$ acts on algebra factors and $\Delta$ acts
on the indexing set,
$\Delta \circ \bar{B}(f) = (\bar{B}(f) \otimes \bar{B}(f))
\circ \Delta$. The counit compatibility and preservation of
identities and composition are immediate from the definition.
\end{proof}

\begin{corollary}[Natural transformation property; \ClaimStatusConditional]\label{cor:bar-natural}
\label{cor:why-functoriality}
A commutative diagram of chiral algebras induces a commutative diagram of bar complexes: if $k \circ f = g \circ h$ in $\mathsf{ChirAlg}_X$, then $\bar{B}(k) \circ \bar{B}(f) = \bar{B}(g) \circ \bar{B}(h)$ in $\mathsf{dgCoalg}_X$. This is immediate from functoriality.
\end{corollary}

\begin{proof}
By Corollary~\ref{cor:bar-functorial}, the bar construction is a
functor. Applying that functor to the equality
$k \circ f = g \circ h$ gives
$\bar{B}(k \circ f) = \bar{B}(g \circ h)$, and functoriality expands
this as
$\bar{B}(k)\circ\bar{B}(f) = \bar{B}(g)\circ\bar{B}(h)$.
\end{proof}

\begin{proposition}[Based comparison of bar models; \ClaimStatusProvedHere]
\label{prop:model-independence}
\index{model independence}
Let $\cA$ be a chiral algebra on~$X$.
Let $B_1$ and $B_2$ be admissible dg factorisation-coalgebra models,
equipped with specified equivalences
\[
 u_i\colon B_i\xrightarrow{\;\sim\;}\mathcal B_X(\cA),
 \qquad i=1,2,
\]
to a fixed bar object.  Then the mapping space between
$(B_1,u_1)$ and $(B_2,u_2)$ in the slice
$\infty$-category
$\mathsf{dgCoalg}_{X,/\mathcal B_X(\cA)}$ is contractible.  Hence
cohomology and every functorial filtered invariant of the two based
models are identified through a contractible space of maps over
$\mathcal B_X(\cA)$.

After forgetting the maps~$u_i$, the equivalence component carries the
natural action of
$\operatorname{Aut}(\mathcal B_X(\cA))$.  Thus an unbased comparison
is a choice in the corresponding automorphism torsor.
\end{proposition}

\begin{proof}
In the slice $\infty$-category over $\mathcal B_X(\cA)$, each
$(B_i,u_i)$ is equivalent to the terminal object
$(\mathcal B_X(\cA),\mathrm{id})$.  The mapping space from any object
to a terminal object is contractible, and equivalence with the target
preserves mapping spaces.  This proves the based assertion.

For the unbased assertion, choose one equivalence
$f\colon B_1\xrightarrow{\sim}B_2$.  Composition identifies every
other equivalence in the same component with
$\alpha\circ f$ for a homotopy automorphism~$\alpha$ of~$B_2$.
Transport along~$u_2$ identifies this automorphism space with
$\operatorname{Aut}(\mathcal B_X(\cA))$.
\end{proof}

\subsection{Coalgebra structure}

\begin{theorem}[Bar coalgebra; \ClaimStatusConditional]\label{thm:bar-coalgebra}
The bar complex carries a natural coalgebra structure:
\[\Delta: \bar{B}^{\text{geom}}(\mathcal{A}) \to \bar{B}^{\text{geom}}(\mathcal{A}) \otimes \bar{B}^{\text{geom}}(\mathcal{A})\]
induced by the diagonal map $X \to X \times X$.
This is the coshuffle coproduct on the symmetric bar $\mathrm{Sym}^c(s^{-1}\bar{\mathcal{A}})$, summing over all $2^n$ unordered bipartitions (Theorem~\textup{\ref{thm:three-bar-complexes}}\textup{(ii)}). The ordered bar $T^c(s^{-1}\bar{\mathcal{A}})$ carries a distinct deconcatenation coproduct with $n+1$ terms (Theorem~\textup{\ref{thm:three-bar-complexes}}\textup{(i)}); the coshuffle is its $\Sigma_n$-coinvariant image.
\end{theorem}

\begin{proof}
Define $\Delta$ on tensors by summing over ordered bipartitions of the input
index set and restricting the logarithmic form to the corresponding boundary
factorization stratum, as in the explicit formula used in
Theorem~\ref{thm:coassociativity-complete}. This defines a morphism
$\bar{B}^{\mathrm{geom}}(\mathcal{A})\to
\bar{B}^{\mathrm{geom}}(\mathcal{A})\otimes\bar{B}^{\mathrm{geom}}(\mathcal{A})$.

Coassociativity of this coproduct is proved in
Theorem~\ref{thm:coassociativity-complete}; counit compatibility is proved in
Theorem~\ref{thm:counit-axioms}. Hence the stated coproduct gives a coalgebra
structure on $\bar{B}^{\mathrm{geom}}(\mathcal{A})$.
\end{proof}

\subsection{Coalgebra axioms for the bar coalgebra}
\label{subsec:coalgebra-axioms-complete}

\begin{theorem}[Coassociativity; \ClaimStatusProvedHere]\label{thm:coassociativity-complete}
\index{coassociativity}
The coproduct $\Delta: \bar{B}_n(\mathcal{A}) \to \bigoplus_{p+q=n} \bar{B}_p(\mathcal{A}) \otimes \bar{B}_q(\mathcal{A})$ is coassociative:
\[(\Delta \otimes \text{id}) \circ \Delta = (\text{id} \otimes \Delta) \circ \Delta\]
\end{theorem}

\begin{proof}

\smallskip\noindent\emph{Step 1: Definition of coproduct.}

Recall that the coproduct on $\bar{B}_n(\mathcal{A})$ is defined by: for $(a_0 \otimes \cdots \otimes a_n \otimes \omega) \in \bar{B}_n(\mathcal{A})$,
\[\Delta(a_0 \otimes \cdots \otimes a_n \otimes \omega) = \sum_{I \sqcup J = [0,n]} (a_I \otimes \omega_I) \otimes (a_J \otimes \omega_J)\]
where:
\begin{itemize}
\item $I, J \subseteq [0,n]$ partition the index set
\item $a_I = a_{i_0} \otimes \cdots \otimes a_{i_p}$ for $I = \{i_0, \ldots, i_p\}$
\item $\omega_I$ is the restriction of $\omega$ to the configuration space $\overline{C}_{|I|}(X)$
\end{itemize}

\smallskip\noindent\emph{Step 2: Left side, $(\Delta \otimes \mathrm{id}) \circ \Delta$.}

Apply $\Delta$ first:
\[\Delta(a_0 \otimes \cdots \otimes a_n \otimes \omega) = \sum_{I \sqcup J = [0,n]} (a_I \otimes \omega_I) \otimes (a_J \otimes \omega_J)\]

Now apply $\Delta \otimes \text{id}$ to each term:
\begin{align*}
&(\Delta \otimes \text{id})\left((a_I \otimes \omega_I) \otimes (a_J \otimes \omega_J)\right) \\
&= \Delta(a_I \otimes \omega_I) \otimes (a_J \otimes \omega_J) \\
&= \left(\sum_{I' \sqcup I'' = I} (a_{I'} \otimes \omega_{I'}) \otimes (a_{I''} \otimes \omega_{I''})\right) \otimes (a_J \otimes \omega_J) \\
&= \sum_{I' \sqcup I'' = I} (a_{I'} \otimes \omega_{I'}) \otimes (a_{I''} \otimes \omega_{I''}) \otimes (a_J \otimes \omega_J)
\end{align*}

Summing over all partitions $I \sqcup J = [0,n]$:
\[(\Delta \otimes \text{id}) \circ \Delta = \sum_{I' \sqcup I'' \sqcup J = [0,n]} (a_{I'} \otimes \omega_{I'}) \otimes (a_{I''} \otimes \omega_{I''}) \otimes (a_J \otimes \omega_J)\]

\smallskip\noindent\emph{Step 3: Right side, $(\mathrm{id} \otimes \Delta) \circ \Delta$.}

Apply $\Delta$ first (same as before):
\[\Delta(a_0 \otimes \cdots \otimes a_n \otimes \omega) = \sum_{I \sqcup J = [0,n]} (a_I \otimes \omega_I) \otimes (a_J \otimes \omega_J)\]

Now apply $\text{id} \otimes \Delta$:
\begin{align*}
&(\text{id} \otimes \Delta)\left((a_I \otimes \omega_I) \otimes (a_J \otimes \omega_J)\right) \\
&= (a_I \otimes \omega_I) \otimes \Delta(a_J \otimes \omega_J) \\
&= (a_I \otimes \omega_I) \otimes \left(\sum_{J' \sqcup J'' = J} (a_{J'} \otimes \omega_{J'}) \otimes (a_{J''} \otimes \omega_{J''})\right) \\
&= \sum_{J' \sqcup J'' = J} (a_I \otimes \omega_I) \otimes (a_{J'} \otimes \omega_{J'}) \otimes (a_{J''} \otimes \omega_{J''})
\end{align*}

Summing over all partitions $I \sqcup J = [0,n]$:
\[(\text{id} \otimes \Delta) \circ \Delta = \sum_{I \sqcup J' \sqcup J'' = [0,n]} (a_I \otimes \omega_I) \otimes (a_{J'} \otimes \omega_{J'}) \otimes (a_{J''} \otimes \omega_{J''})\]

\smallskip\noindent\emph{Step 4: Comparison.}

\emph{LHS.} Sum over ordered partitions $(I', I'', J)$ with $I' \sqcup I'' \sqcup J = [0,n]$

\emph{RHS.} Sum over ordered partitions $(I, J', J'')$ with $I \sqcup J' \sqcup J'' = [0,n]$

These are the same set of ordered triples, in different notation.

Relabeling $I' \to K_1$, $I'' \to K_2$, $J \to K_3$ on LHS and $I \to K_1$, $J' \to K_2$, $J'' \to K_3$ on RHS:

Both sides equal:
\[\sum_{K_1 \sqcup K_2 \sqcup K_3 = [0,n]} (a_{K_1} \otimes \omega_{K_1}) \otimes (a_{K_2} \otimes \omega_{K_2}) \otimes (a_{K_3} \otimes \omega_{K_3})\]

Therefore: $(\Delta \otimes \text{id}) \circ \Delta = (\text{id} \otimes \Delta) \circ \Delta$.
\end{proof}

\begin{theorem}[Counit axioms; \ClaimStatusProvedHere]\label{thm:counit-axioms}
The counit $\epsilon: \bar{B}_n(\mathcal{A}) \to \mathbb{C}$ satisfies:
\begin{enumerate}
\item \emph{Left counit:} $(\epsilon \otimes \text{id}) \circ \Delta = \text{id}$
\item \emph{Right counit:} $(\text{id} \otimes \epsilon) \circ \Delta = \text{id}$
\end{enumerate}
\end{theorem}

\begin{proof}

The counit $\epsilon$ is:
\[\epsilon(a_0 \otimes \cdots \otimes a_n \otimes \omega) =
\begin{cases}
\langle a_0, 1 \rangle & n = 0 \\
0 & n > 0
\end{cases}\]

\smallskip\noindent\emph{Left counit axiom.}

For $(a_0 \otimes \cdots \otimes a_n \otimes \omega) \in \bar{B}_n(\mathcal{A})$:
\begin{align*}
&(\epsilon \otimes \text{id})(\Delta(a_0 \otimes \cdots \otimes a_n \otimes \omega)) \\
&= (\epsilon \otimes \text{id})\left(\sum_{I \sqcup J = [0,n]} (a_I \otimes \omega_I) \otimes (a_J \otimes \omega_J)\right) \\
&= \sum_{I \sqcup J = [0,n]} \epsilon(a_I \otimes \omega_I) \cdot (a_J \otimes \omega_J)
\end{align*}

Now $\epsilon(a_I \otimes \omega_I) = 0$ unless $|I| = 0$ (i.e., $I = \emptyset$).

When $I = \emptyset$, we have $J = [0,n]$ and:
\[\epsilon(1 \otimes \omega_{\emptyset}) \cdot (a_0 \otimes \cdots \otimes a_n \otimes \omega) = 1 \cdot (a_0 \otimes \cdots \otimes a_n \otimes \omega)\]

Therefore:
\[(\epsilon \otimes \text{id}) \circ \Delta = \text{id}\]

\smallskip\noindent\emph{Right counit axiom.}

Similarly:
\begin{align*}
&(\text{id} \otimes \epsilon)(\Delta(a_0 \otimes \cdots \otimes a_n \otimes \omega)) \\
&= \sum_{I \sqcup J = [0,n]} (a_I \otimes \omega_I) \cdot \epsilon(a_J \otimes \omega_J)
\end{align*}

$\epsilon(a_J \otimes \omega_J) = 0$ unless $|J| = 0$ (i.e., $J = \emptyset$).

When $J = \emptyset$, we have $I = [0,n]$ and:
\[(a_0 \otimes \cdots \otimes a_n \otimes \omega) \cdot \epsilon(1 \otimes \omega_{\emptyset}) = (a_0 \otimes \cdots \otimes a_n \otimes \omega) \cdot 1\]

Therefore:
\[(\text{id} \otimes \epsilon) \circ \Delta = \text{id}\]
\end{proof}

\begin{corollary}[Bar complex is DG-coalgebra; \ClaimStatusConditional]\label{cor:bar-is-dgcoalg}
The bar complex $\bar{B}^{\text{geom}}(\mathcal{A})$ with:
\begin{itemize}
\item Differential $d = d_{\text{int}} + d_{\text{res}} + d_{\text{dR}}$ (satisfying $d^2 = 0$)
\item Coproduct $\Delta$ (coassociative)
\item Counit $\epsilon$ (satisfying counit axioms)
\end{itemize}
is a differential graded coalgebra.
\end{corollary}
\begin{proof}
This assembles Theorem~\ref{thm:bar-nilpotency-complete} ($d^2 = 0$), Theorem~\ref{thm:coassociativity-complete} (coassociativity of $\Delta$), and Theorem~\ref{thm:counit-axioms} (counit axioms). The compatibility $d \circ \Delta = \Delta \circ d$ is Theorem~\ref{thm:bar-coalgebra}.
\end{proof}

\begin{remark}[Geometric meaning of coassociativity]\label{rem:coassoc-geometric}
Coassociativity asserts that the order of splitting is immaterial: $(\Delta \otimes \text{id}) \circ \Delta$ (split left first) equals $(\text{id} \otimes \Delta) \circ \Delta$ (split right first). Geometrically, corners of $\overline{C}_n(X)$ where multiple divisors intersect admit approach from different directions, yielding consistent boundary data.
\end{remark}

\begin{theorem}[Differential is coderivation; \ClaimStatusProvedHere]\label{thm:diff-is-coderivation}
\textup{[Regime: genus-$0$
\textup{(}Convention~\textup{\ref{conv:regime-tags})};
at genus~$g \geq 1$ the curved differential~$\dfib$ fails the
coderivation property, see
\textup{Chapter~\ref{chap:higher-genus}}.]}

The differential $d$ on $\bar{B}(\mathcal{A})$ is a coderivation:
\[\Delta \circ d = (d \otimes \text{id} + \text{id} \otimes d) \circ \Delta\]
\end{theorem}

\begin{proof}
The coderivation identity $\Delta \circ d = (d \otimes \mathrm{id} + \mathrm{id} \otimes d) \circ \Delta$ is verified for each component of $d = d_{\mathrm{int}} + d_{\mathrm{res}} + d_{\mathrm{dR}}$ separately.

For $d_{\mathrm{int}}$: the internal differential acts on each tensor factor $\phi_i$ independently, and the coshuffle coproduct $\Delta(a_0 \otimes \cdots \otimes a_n) = \sum_{I \sqcup J = [0,n]} \pm\, (a_I) \otimes (a_J)$ (sum over all $2^{n+1}$ ordered bipartitions, as defined in Theorem~\ref{thm:coassociativity-complete}) splits the tensor factors. The Leibniz rule $\Delta \circ d_{\mathrm{int}} = (d_{\mathrm{int}} \otimes \mathrm{id} + \mathrm{id} \otimes d_{\mathrm{int}}) \circ \Delta$ follows because $d_{\mathrm{int}}$ acts within each factor and commutes with the splitting.

For $d_{\mathrm{res}}$: the residue at $D_{ij}$ (collision of points $i$ and $j$) produces a term in bar degree $n-1$. The coproduct restricts to boundary strata: $\Delta|_{D_{ij}}$ splits the remaining points into two groups, and the residue extraction commutes with this splitting because the divisor $D_{ij}$ lies entirely within the factor containing both $i$ and $j$.

For $d_{\mathrm{dR}}$: the exterior derivative satisfies $d_{\mathrm{dR}}(\omega_I \wedge \omega_J) = d_{\mathrm{dR}}\omega_I \wedge \omega_J + (-1)^{|\omega_I|}\omega_I \wedge d_{\mathrm{dR}}\omega_J$, which is the Leibniz rule with respect to the K\"unneth decomposition $\Omega^*(C_{n+1}(X)) \cong \bigoplus_{I \sqcup J} \Omega^*(C_{|I|}(X)) \otimes \Omega^*(C_{|J|}(X))$ used by $\Delta$.

This is the standard coderivation property for bar differentials; see~\cite{LV12}, Proposition~2.2.1 for the algebraic version.
\end{proof}

\begin{definition}[Genus-graded geometric bar complex]\label{def:geom-bar}
For a chiral algebra $\mathcal{A}$ on a Riemann surface $\Sigma_g$ of genus $g$, the \emph{genus-graded geometric bar complex} is the bigraded complex:
\[
\bar{B}^{(g)}_{p,q}(\mathcal{A}) = \Gamma\left(\overline{C}_{p+1}(\Sigma_g), j_*j^*\mathcal{A}^{\boxtimes(p+1)} \otimes \Omega^q_{\overline{C}_{p+1}(\Sigma_g)}(\log D^{(g)}) \otimes \text{or}_{p+1}^{(g)}\right)
\]
where:
\begin{itemize}
\item $\overline{C}_{p+1}(\Sigma_g)$ is the Fulton--MacPherson compactification at genus $g$
\item $D^{(g)} = \overline{C}_{p+1}(\Sigma_g) \setminus C_{p+1}(\Sigma_g)$ is the boundary divisor with genus-dependent stratification
\item $j: C_{p+1}(\Sigma_g) \hookrightarrow \overline{C}_{p+1}(\Sigma_g)$ is the open inclusion
\item $\Omega^q_{\overline{C}_{p+1}(\Sigma_g)}(\log D^{(g)})$ includes logarithmic forms and period integrals
\item $\text{or}_{p+1}^{(g)}$ is the genus-graded orientation bundle
\end{itemize}

The total bar complex is:
\[\bar{B}(\mathcal{A}) = \bigoplus_{g=0}^{\infty} \bar{B}^{(g)}(\mathcal{A})\]
\end{definition}

\begin{remark}[Orientation bundle]
The orientation bundle $\text{or}_{p+1}^{(g)}$ is necessary because configuration spaces are not naturally
oriented at each genus. It is the determinant line of $T_{C_{p+1}(\Sigma_g)}$ with genus-dependent corrections, ensuring consistent signs across all face maps. At genus~$0$, this guarantees $\dzero^2 = 0$ (Theorem~\ref{thm:bar-nilpotency-complete}); at genus~$g \geq 1$, the fiberwise differential satisfies $\dfib^{\,2} = \kappa(\cA)\cdot\omega_g$ and the total corrected differential $\Dg{g}^2 = 0$ is established in Chapter~\ref{chap:higher-genus}.
\end{remark}

\subsection{The differential: detailed components}

The three-component decomposition $d = d_{\text{internal}} + d_{\text{residue}} + d_{\text{form}}$
of Definition~\ref{def:bar-differential-complete} is now expanded at genus~$g$.

\subsubsection{Internal differential}

\begin{definition}[Internal differential]
For $\alpha = \alpha_1 \otimes \cdots \otimes \alpha_{n+1} \otimes \omega \otimes \theta \in
\bar{B}^{n,q}_{\text{geom}}(\mathcal{A})$ where $\theta \in \text{or}_{n+1}$:
\[
d_{\text{int}}(\alpha) = \sum_{i=1}^{n+1} (-1)^{|\alpha_1| + \cdots + |\alpha_{i-1}|}
\alpha_1 \otimes \cdots \otimes d_{\mathcal{A}}(\alpha_i) \otimes \cdots \otimes \alpha_{n+1} \otimes \omega \otimes \theta
\]
where $d_{\mathcal{A}}$ is the internal differential on $\mathcal{A}$ (if present) and $|\alpha_i|$ denotes
the cohomological degree.
\end{definition}

\subsubsection{Factorization differential}

\begin{definition}[Factorization differential]\label{def:diff-fact}
 The factorization differential encodes the chiral algebra structure:
 \[
 d_{\text{fact}}
 =
 \sum_{1 \leq i < j \leq n+1}
 (-1)^{\sigma(i,j)}
 \sum_{m\ge0}
 \bigl((-)_{(m)}(-)\bigr)_{ij}
 \otimes
 \operatorname{Res}^{\mathrm{form}}_{D_{ij}}
 \]
 where the sign is:
 \[\sigma(i,j) = i + j + \left(\sum_{k=1}^{i-1} |\alpha_k|\right) \cdot |\eta_{ij}|\]
 (Here $|\alpha_k|$ denotes the cohomological degree of the $k$-th element and $|\eta_{ij}| = 1$ for the logarithmic 1-form.)

 The operator $\operatorname{Res}^{\mathrm{form}}_{D_{ij}}$ is the
 Poincar\'e residue on the logarithmic form factor. The OPE mode
 $(-)_{(m)}(-)$ is supplied by the chiral product. Thus
 $d_{\text{fact}}$ does not take the analytic residue of an OPE pole
 multiplied by $\eta_{ij}$; it composes the form residue with the
 chiral mode projection.

 This accounts for:
 \begin{itemize}
 \item Koszul sign from moving the normal logarithmic direction past the fields $\alpha_k$
 \item Orientation of the divisor $D_{ij}$
 \item Parity of the permutation after collision
 \end{itemize}
 \end{definition}

 \begin{lemma}[Orientation convention; \ClaimStatusProvedHere]\label{lem:orientation}
 Fix orientations on boundary divisors by:
 \begin{enumerate}
 \item For $D_{ij}$ where $z_i = z_j$:
 \[\text{or}_{D_{ij}} = dz_1 \wedge \cdots \wedge \widehat{dz_i} \wedge \cdots \wedge dz_{n+1}\]
 (omit $dz_i$, keep others including $dz_j$)

 \item For codimension-2 strata $D_{ijk} = D_{ij} \cap D_{jk}$:
 \[\text{or}_{D_{ijk}} = \text{or}_{D_{ij}} \wedge \text{or}_{D_{jk}}\]

 \item This implies the relation:
 \[\text{or}_{D_{ijk}} = -\text{or}_{D_{ik}} \wedge \text{or}_{D_{jk}} = \text{or}_{D_{jk}} \wedge \text{or}_{D_{ik}}\]
 \end{enumerate}

 These choices ensure $\partial^2 = 0$ for the boundary operator on $\overline{C}_{n+1}(X)$.
 \end{lemma}

 \begin{proof}
 The FM compactification $\overline{C}_{n+1}(X)$ is constructed by iterated real blow-up along diagonals (\cite[\S2]{FM94}). After blow-up, each codimension-1 boundary stratum $D_S$ (indexed by subsets $S \subset [n+1]$ with $|S| \geq 2$) is a smooth divisor, and distinct divisors $D_S, D_T$ meet transversally when $S \subset T$ or $T \subset S$ (nested), and are disjoint otherwise. A codimension-2 stratum therefore arises as $D_S \cap D_T$ where $S \subsetneq T$ (a two-step degeneration), and the two orderings of the iterated boundary (first collapse $S$, then collapse $T/S$, versus directly collapsing $T$ and then refining to $S$) yield opposite orientations by the antisymmetry of the normal bundle orientations (Corollary~\ref{cor:residues-anticommute}).
 \end{proof}

 \begin{remark}[Significance of signs]
 The sign conventions are forced by $d^2 = 0$, following Kontsevich's principle: ``signs should be determined by geometry, not combinatorics.''
 \end{remark}

 \begin{lemma}[Residue properties; \ClaimStatusProvedHere]
 \label{lem:residue-properties}
 The residue operation satisfies:
 \begin{enumerate}
 \item $\text{Res}_{D_{ij}}^2 = 0$ for the same boundary normal direction
 \item For disjoint pairs: $\text{Res}_{D_{ij}} \circ \text{Res}_{D_{k\ell}} = -\text{Res}_{D_{k\ell}} \circ \text{Res}_{D_{ij}}$
 \item For overlapping pairs with $j = k$: contributions combine by the Arnold relation on forms and the Borcherds identity on OPE modes
 \end{enumerate}
 \end{lemma}

 \begin{proof}
 Part~(1): A section $\omega \in \Omega^*_{\log}(\overline{C}_{n+1}(X))$ has at most a simple pole along $D_{ij}$. In a local coordinate $u$ with $D_{ij} = \{u = 0\}$, write
 $\omega = f(u)\,d\log u + g(u)\,(\text{terms regular in }u)$
 where $f(0) = \mathrm{Res}_{D_{ij}}\omega$. The residue $\mathrm{Res}_{D_{ij}}\omega = f(0)$ is a section on $D_{ij}$ with no pole along $D_{ij}$, so $\mathrm{Res}_{D_{ij}}^2\omega = \mathrm{Res}_{D_{ij}}(f(0)) = 0$.

 Part~(2): This is Corollary~\ref{cor:residues-anticommute}: for disjoint index pairs, the divisors $D_{ij}$ and $D_{k\ell}$ are transverse, and the anticommutativity follows from the antisymmetry of $d\log u \wedge d\log v$.

 Part~(3): For overlapping pairs with $j = k$, three codimension-1 strata $D_{ij}$, $D_{jk}$, and $D_{ik}$ meet at the triple collision stratum $D_{ijk}$. The form part of the iterated residues:
 \[
 \mathrm{Res}_{D_{jk}} \circ \mathrm{Res}_{D_{ij}} + \mathrm{Res}_{D_{ij}} \circ \mathrm{Res}_{D_{jk}} + \mathrm{Res}_{D_{ik}} \circ \mathrm{Res}_{D_{ij}}
 \]
 is the codimension-two Arnold cancellation. After applying the
 corresponding OPE modes, the algebraic coefficients cancel by the
 Borcherds identity, as verified in Case~(ii) of
 Theorem~\ref{thm:bar-nilpotency-complete}.
 \end{proof}

% Well-definedness and anticommutativity of residues: see
% Lemma~\ref{lem:residue-properties} and Corollary~\ref{cor:residues-anticommute}.

\subsubsection{Configuration differential}

\begin{definition}[Configuration differential]
 The configuration differential is the de Rham differential on forms:
 \[d_{\text{config}} = d_{\text{config}}^{\text{dR}} + d_{\text{config}}^{\text{Lie*}}\]
 where the de~Rham component $d_{\text{config}}^{\text{dR}} = \text{id}_{\mathcal{A}^{\boxtimes(n+1)}} \otimes d_{\text{dR}} \otimes \text{id}_{\text{or}}$ acts on the differential forms, and the Lie$^*$ component $d_{\text{config}}^{\text{Lie*}} = \sum_{I \subset [n+1]} (-1)^{\epsilon(I)} d_{\text{Lie}}^{(I)} \otimes \text{id}_{\Omega^*}$ acts via the Lie$^*$ algebra structure (when present).

 For general chiral algebras without Lie* structure, $d_{\text{config}}^{\text{Lie*}} = 0$.
 \end{definition}

\begin{remark}[Geometric meaning]
 The configuration differential captures how $\mathcal{A}$ varies over configuration space: $d_{\text{dR}}$ measures variation of insertion points, $d_{\text{Lie*}}$ encodes infinitesimal symmetries. This parallels the Cartan model for equivariant cohomology.
 \end{remark}

\subsection{Explicit residue computations}\label{sec:residue-calculus}

\begin{lemma}[Geometric--operadic sign comparison; \ClaimStatusProvedHere]
\label{lem:LV-sign-comparison}
\label{rem:LV-signs}
Our sign conventions for the bar construction follow the geometric approach, which differs slightly from the operadic conventions in Loday--Vallette~\cite{LV12}.

The differences are as follows.
\begin{enumerate}
\item \emph{Koszul sign rule}: We use the \emph{geometric} Koszul rule where moving a differential form of degree $p$ past an operator of degree $q$ introduces $(-1)^{pq}$.

\item \emph{Residue orientation}: Our residues include an orientation factor from the normal bundle to collision divisors. This introduces signs when collision divisors intersect.

\item \emph{Suspension}: Loday--Vallette use operadic suspension $s: V \to sV$ with $|s| = 1$. We work with geometric forms directly, so suspension is implicit in the degree shift of $\Omega^n(\log D)$.
\end{enumerate}

\emph{Translation between conventions.}
\begin{center}
\begin{tabular}{|>{\raggedright\arraybackslash}p{0.43\textwidth}|>{\raggedright\arraybackslash}p{0.43\textwidth}|}
\hline
\textbf{Loday--Vallette (Operadic)} & \textbf{Ours (Geometric)} \\
\hline
$d_{op}(s a_1 \otimes \cdots \otimes s a_n)$ & $d_{geom}(a_1 \otimes \cdots \otimes a_n \otimes \omega_n)$ \\
Sign: $(-1)^{|a_1| + \cdots + |a_{i-1}|}$ &
Internal sign:
$(-1)^{\sum_{q<i}(|a_q|-1)}$; residue sign:
$(-1)^{\sum_{q<i}(|a_q|-1)+|a_i|}$ plus the FM
boundary-orientation parity of
Proposition~\ref{prop:ordered-bar-local-differential-identities}. \\
Suspension degree $|s a_i| = |a_i| + 1$ & Form degree $|\omega| = n$ \\
\hline
\end{tabular}
\end{center}

The two conventions agree under the suspension isomorphism
\[
\iota_n\colon
s\bar\cA^{\otimes n}
\xrightarrow{\;\sim\;}
\bar\cA^{\otimes n}\otimes
\det H^0\!\bigl(\FM_n(X),\Omega^1_{\FM_n(X)}(\log D)\bigr),
\]
up to the single scalar fixed by orienting the one-point stratum. This
scalar is absorbed into the normalization of the cyclic pairing.

The nine-term proof of $d^2 = 0$ (Theorem~\ref{thm:bar-nilpotency-complete}) uses geometric signs throughout. Translating to operadic conventions via the dictionary above preserves $d^2 = 0$.
\end{lemma}

\begin{proof}
Both sign systems are multiplicative characters of the groupoid of
ordered collision trees. It is therefore enough to compare them on two
generators: an adjacent transposition of two inputs and a single
codimension-one residue face. For an adjacent transposition of
homogeneous inputs \(a_i,a_{i+1}\), the Loday--Vallette suspended word
contributes
\[
(-1)^{(|a_i|+1)(|a_{i+1}|+1)}
=
(-1)^{|a_i||a_{i+1}|}
(-1)^{|a_i|+|a_{i+1}|}
(-1).
\]
Under \(\iota_n\), the first factor is the ordinary Koszul sign of the
fields, the middle factor is the suspension/form-degree shift, and the
last factor is the orientation sign for interchanging the two
configuration coordinates. This is exactly the enhanced orientation
convention of Convention~\ref{conv:orientations-enhanced} together with
the desuspension convention of
Convention~\ref{conv:A-infinity-2-ordered-bar-signs}.

For a residue face \(D_{ij}\), the operadic cocomposition sign is the
sign obtained by moving the suspended collision block through the
preceding suspended factors. Under \(\iota_n\) this is the sign obtained
by moving the local affine/formal representative \(d\log(z_i-z_j)\) of
the logarithmic normal form through the preceding coefficient factors
before applying the Poincare residue.
The outward-normal convention supplies the same normal-orientation
character. Thus \(\iota_n d_{\bar B}\iota_n^{-1}\) is the geometric
operator \(d_{\mathrm{dR}}+d_{\mathrm{res}}\) with the signs used in
this chapter. The only remaining freedom is the orientation of the
one-point stratum, a global scalar; fixing it fixes all tree signs
because every tree is generated by adjacent transpositions and residue
faces. Hence the geometric and operadic bar differentials are
conjugate under \(\iota_n\), and nilpotence is preserved.
\end{proof}

\begin{theorem}[Ordered, symmetric, and Harrison bars; \ClaimStatusConditional]\label{thm:three-bar-complexes}
\index{bar complex!three variants|textbf}
\index{coshuffle coproduct|textbf}
\index{deconcatenation coproduct|textbf}
\index{Eulerian idempotent!bar complex decomposition|textbf}
\index{PBW isomorphism!coalgebra level}
\index{R-matrix!descent from $T^c$ to $\mathrm{Sym}^c$}
Work in characteristic zero.
Let\/ $\cA$ be an augmented chiral algebra on\/~$X$ with
augmentation ideal\/ $\bar\cA = \ker(\varepsilon)$.
There is a chain of comparison maps
\begin{equation}\label{eq:three-bar-inclusions}
T^c(s^{-1}\bar\cA)
\xrightarrow{\ \pi_\Sigma\ }
\mathrm{Sym}^c(s^{-1}\bar\cA)
\xrightarrow{\ e^{(1)}\ }
\mathrm{Lie}^c(s^{-1}\bar\cA)
\end{equation}
whose three terms are:
\begin{enumerate}[label=\textup{(\roman*)}]
\item \textup{(Ordered tensor bar.)}
 $T^c(s^{-1}\bar\cA)$:
 the cofree conilpotent coassociative
 \textup{(}non-cocommutative\textup{)} coalgebra on
 $s^{-1}\bar\cA$.
 In degree~$n$ the component is
 $(s^{-1}\bar\cA)^{\otimes n}$,
 with the deconcatenation coproduct
 \[
 \Delta(a_1 \otimes \cdots \otimes a_n)
 \;=\;
 \sum_{i=0}^{n}\,
 (a_1 \otimes \cdots \otimes a_i)
 \otimes
 (a_{i+1} \otimes \cdots \otimes a_n),
 \]
 a sum over the\/ $n+1$ consecutive cuts.
 Under the finite locally free quadratic package of
 Definition~\ref{def:chirass-dual-cooperad} and
 Proposition~\ref{prop:chirAss-self-dual}, this ordered tensor
 cooperad is the chiral cooperad \((\chirAss)^{!,c}\); without that
 package it remains the proved level-$2$ ordered tensor carrier
 carrying the $R$-matrix tower.
\item \textup{(Symmetric/coshuffle bar.)}
 $\mathrm{Sym}^c(s^{-1}\bar\cA)$:
 the cofree cocommutative coassociative conilpotent
 coalgebra on $s^{-1}\bar\cA$.
 In degree~$n$ the component is
 $\bigl((s^{-1}\bar\cA)^{\otimes n}\bigr)_{S_n}$,
 with the coshuffle coproduct
 \[
 \Delta(a_1 \cdots a_n)
 \;=\;
 \sum_{I \sqcup J = [n]} \pm\, (a_I) \otimes (a_J),
 \]
 a sum over all\/ $2^n$ unordered bipartitions.
\item \textup{(Harrison/coLie bar.)}
 $\mathrm{Lie}^c(s^{-1}\bar\cA)$:
 the weight-$1$ Eulerian summand of the symmetric bar, equivalently
 the cofree conilpotent coLie coalgebra on $s^{-1}\bar\cA$.
 In degree~$n$ it carries the Lie representation
 of dimension $(n{-}1)!$, and its coproduct is the coLie cobracket.
 This is the Koszul dual cooperad of\/ $\chirCom$.
\end{enumerate}
The following structural results hold.
\begin{enumerate}[label=\textup{(\alph*)}]
\item \textup{(Coinvariant projection.)}
 The map
 $\pi_\Sigma\colon T^c(s^{-1}\bar\cA)\twoheadrightarrow
 \mathrm{Sym}^c(s^{-1}\bar\cA)$ is the degreewise
 $\Sigma_n$-coinvariant quotient. It is a morphism of graded
 coalgebras from deconcatenation to coshuffle, and it is a chain map
 whenever the ordered bar differential is $\Sigma_n$-equivariant.
 Its kernel is the non-symmetric ordered sector.
\item \textup{(Averaging idempotent.)}
 The characteristic-zero averaging section
 $j_\Sigma\colon\mathrm{Sym}^c(s^{-1}\bar\cA)\to
 T^c(s^{-1}\bar\cA)$ satisfies
 $\pi_\Sigma j_\Sigma=\id$ and makes
 $j_\Sigma\pi_\Sigma$ the symmetrizing idempotent on $T^c$.
 This splitting is a statement about graded coalgebras; it is not a
 quasi-isomorphism statement.
\item \textup{(PBW decomposition.)}
 The Eulerian idempotent decomposition gives a
 direct-sum splitting of graded coalgebras
 \begin{equation}\label{eq:eulerian-bar-decomposition}
 \mathrm{Sym}^c(s^{-1}\bar\cA)
 \;\cong\;
 \bigoplus_{r \ge 1}
 \bigl(\mathrm{Lie}^c(s^{-1}\bar\cA)\bigr)^{\odot r}
 \end{equation}
 as $S_n$-modules, where\/ $\odot$ denotes the
 symmetric product of cooperads.
 The first summand\/ $(r = 1)$ is the Harrison
 complex; the higher summands are its symmetric
 decomposables.
\item \textup{(Acyclicity criterion.)}
 The projection $\pi_\Sigma$ is a quasi-isomorphism exactly on those
 surfaces where the non-symmetric ordered kernel is acyclic. This
 acyclicity is the PBW/Koszul averaging hypothesis used for the
 symmetric Theorem~A/H statements; it is false as an unrestricted
 assertion about ordered chiral bars.
\item \textup{(R-matrix descent.)}
 The $R$-matrix tower on\/ $T^c$ descends to the
 shadow obstruction tower on\/ $\mathrm{Sym}^c$
 via\/ $\pi_\Sigma$, and the Harrison
 subcomplex\/ $\mathrm{Lie}^c$ detects the primitives.
 In the notation of\/
 \S\textup{\ref{subsec:e1-as-primitive}}: at degree\/
 $r$, $\pi_\Sigma(r_r) = S_r$ is the
 $E_\infty$-shadow at degree\/~$r$
 \textup{(}$\kappa$ at\/ $r = 2$;
 the cubic shadow\/ $\mathfrak{C}$ at\/ $r = 3$\textup{)}.
\end{enumerate}
\end{theorem}

\begin{proof}
The three coalgebras are the cofree constructions on
$s^{-1}\bar\cA$ for the cooperads
$\mathrm{Ass}^c$, $\mathrm{Com}^c$, and $\mathrm{Lie}^c$
(Loday--Vallette \cite[Theorem~1.3.6,
Proposition~8.6.9]{LV12}).

\emph{Part\/ \textup{(a)}.}
The projection $\pi_\Sigma$ is the
coinvariant quotient in each arity:
\[
\pi_\Sigma(a_1\otimes\cdots\otimes a_n)=a_1\cdots a_n.
\]
Applying $\pi_\Sigma$ to deconcatenation and then using the
$S_n$-coinvariant universal property gives the coshuffle formula:
every consecutive cut in each ordered representative becomes one of
the unordered bipartitions of $[n]$, with the Koszul sign inherited
from the representative. If the differential on $T^c$ is
$S_n$-equivariant, then the formula defining it commutes with the
coinvariant quotient, so
$d_{\mathrm{Sym}}\pi_\Sigma=\pi_\Sigma d_T$.

\emph{Part\/ \textup{(b)}.}
In characteristic zero the averaging section is
\[
j_\Sigma(a_1\cdots a_n)
=\frac{1}{n!}\sum_{\sigma\in S_n}
\epsilon(\sigma)\,
a_{\sigma(1)}\otimes\cdots\otimes a_{\sigma(n)}.
\]
The equality $\pi_\Sigma j_\Sigma=\id$ is immediate, and
$j_\Sigma\pi_\Sigma$ is the standard symmetrizing idempotent. For
$n\ge2$ its kernel contains, for homogeneous $a,b$,
\[
a\otimes b-(-1)^{(|a|-1)(|b|-1)}b\otimes a,
\]
so ordered information is genuinely lost by $\pi_\Sigma$.

\emph{Part\/ \textup{(c)}.}
The Eulerian idempotents
$e^{(r)}_n \in \mathbb{Q}[S_n]$
($r = 1, \ldots, n$) form a complete set of
orthogonal idempotents that decompose the regular
representation into weight spaces for the
descent algebra. The weight-$1$ idempotent
$e^{(1)}_n$ projects onto the Lie representations
in~$\mathbb{Q}[S_n]$; its image is precisely the
space of primitive elements of the coshuffle
coalgebra, i.e.\ the Harrison component
$\mathrm{Lie}^c(s^{-1}\bar\cA)$ in degree~$n$.
The direct-sum decomposition
$\mathrm{Sym}^c \cong \bigoplus_{r \ge 1}
(\mathrm{Lie}^c)^{\odot r}$ follows from the
PBW theorem for Lie coalgebras in characteristic zero
(Reutenauer \cite[Chapter~9]{Reutenauer93};
Loday--Vallette \cite[\S1.3.4]{LV12}):
the symmetric coalgebra on a Lie coalgebra is the
universal enveloping cocommutative coalgebra, and the
grading by Eulerian weight supplies the splitting.

\emph{Part\/ \textup{(d)}.}
The quasi-isomorphism assertion is precisely the vanishing of
$H^*(\ker\pi_\Sigma)$. On the PBW/Koszul averaging surface this is the
standard acyclicity input: the non-symmetric shuffle sector is
contracted by the PBW homotopy, so the quotient computes the same
Ext object. Without that hypothesis the displayed antisymmetric class
above can survive, and no formal property of cofree coalgebras kills
it. Thus the unrestricted quasi-isomorphism claim is false; the true
statement is the stated acyclicity criterion.

\emph{Part\/ \textup{(e)}.}
The $R$-matrix $r_r \in \operatorname{End}(V^{\otimes r})
\otimes \mathcal{O}(*\Delta)$ is the degree-$r$
genus-$0$ component of the $E_1$ MC element
$\Theta_\cA^{E_1}$ on the tensor bar $T^c$
(Proposition~\ref{prop:e1-shadow-r-matrix}).
Under the coinvariant projection, the $S_r$-coinvariant of
$r_r$ is the degree-$r$ shadow $S_r$ on the
symmetric bar $\mathrm{Sym}^c$: this is the content
of the coinvariant shadow theorem
(Theorem~\ref{thm:e1-coinvariant-shadow}).
At $r = 2$ the shadow $S_2$ equals $\kappa(\cA)$ for
abelian and scalar families, and $\kappa(\cA) -
\dim(\mathfrak{g})/2$ for non-abelian affine
Kac--Moody (the Sugawara shift;
Theorem~\ref{thm:e1-coinvariant-shadow} for the
precise identification in each family).
The Harrison subcomplex $\mathrm{Lie}^c$ detects
the primitives of $\mathrm{Sym}^c$ under the
coshuffle coproduct; these are the weight-$1$
Eulerian components, which carry the essential
algebraic content (higher Eulerian weights are
symmetric decomposables determined by the primitives).
\end{proof}

\begin{remark}[Primacy direction]\label{rem:primacy-direction}
\index{bar complex!primacy direction}
Information flows in one direction through
\eqref{eq:three-bar-inclusions}: from the ordered tensor bar to the
symmetric coinvariant shadow, then to the Harrison primitive summand.
The projection
$\pi_\Sigma\colon T^c \twoheadrightarrow \mathrm{Sym}^c$ is always
surjective, and it is a quasi-isomorphism only under the acyclicity
criterion of Theorem~\ref{thm:three-bar-complexes}(d). Its kernel
contains the $R$-matrix's non-coinvariant Fourier modes. The ordered
bar is strictly richer; Theorems~A--D use the symmetric bar exactly
where their statements are $\Sigma_n$-invariant.
\end{remark}

\begin{remark}[Topological factorization on the $\bR$-direction]
\label{rem:deconcatenation-topological}
\index{deconcatenation coproduct!topological factorization}%
\index{Swiss-cheese!topological factorization on $\bR$}%
The $n+1$ terms of the deconcatenation coproduct on~$T^c$
correspond to the $n+1$ ways of cutting
$\Conf_n^{<}(\bR) = \{t_1 < \cdots < t_n\}$ into two
consecutive segments $\{t_1,\ldots,t_p\} \sqcup
\{t_{p+1},\ldots,t_n\}$. This is topological factorization
on the~$\bR$-direction of~$\bC \times \bR$ in the
Costello--Gwilliam framework~\cite[\S5.4]{CG17}: the bar
differential is holomorphic factorization on~$\FM_k(\bC)$
(Remark~\ref{rem:costello-gwilliam-factorization}); the bar
coproduct is topological factorization on~$\Conf_n^{<}(\bR)$.
Together they exhibit the two geometric directions that later
feed the Swiss-cheese story. On the bar complex itself they still
form only the single-coloured $\Eone$ dg coalgebra
$T^c(s^{-1}\bar{\cA})$; the actual
$\mathrm{SC}^{\mathrm{ch,top}}$ datum appears only on the
derived-center pair
$(C^\bullet_{\mathrm{ch}}(\cA,\cA),\,\cA)$.
\end{remark}

\begin{remark}[Conventions for this chapter]\label{rem:three-bar-conventions}
\index{bar complex!convention in this chapter}
In this chapter, the unadorned notation
$\bar{B}^{\mathrm{geom}}(\mathcal{A})$ denotes the symmetric Ran bar
$\mathrm{Sym}^c(s^{-1}\bar{\mathcal{A}})$
(Theorem~\ref{thm:three-bar-complexes}(ii)); its coproduct is the
coshuffle on $\mathrm{Sym}^c$, matching the decomposition of
$\overline{C}_{n+1}(X)$ into unordered collision strata. The notation
$B_X^{\mathrm{ord}}(\cA)$ always denotes the ordered tensor bar with
deconcatenation. Statements involving $R$-matrices, Yangian data, or
ordered convolution are made before applying $\pi_\Sigma$; statements
about Theorem~A, Verdier duality, and Ran factorization use the
symmetric quotient.
\end{remark}

\begin{theorem}[Geometric bar \texorpdfstring{$=$}{=} operadic bar; \ClaimStatusProvedHere]
\label{thm:geometric-equals-operadic-bar}
\index{geometric equals operadic bar}
\index{bar construction!algebraic}
\index{twisting morphism!chiral}
Let $\mathcal{P}^{\mathrm{ch}}$ be a chiral operad on $X$ (either
$\chirCom$ or $\chirAss$) and let $\mathcal{A}$ be an augmented
$\mathcal{P}^{\mathrm{ch}}$-algebra. Then the geometric bar complex
$\bar{B}_{\mathrm{geom}}(\mathcal{A})$
(Definition~\ref{def:geometric-bar}) is naturally quasi-isomorphic
to the operadic bar construction
$B_{\mathcal{P}^{\mathrm{ch}}}(\mathcal{A})$ in the sense of
\cite[\S6.5]{LV12}:
\[
\bar{B}_{\mathrm{geom}}(\mathcal{A})
\;\xrightarrow{\;\sim\;}\;
B_{\mathcal{P}^{\mathrm{ch}}}(\mathcal{A}).
\]
The quasi-isomorphism is natural in $\mathcal{A}$ and compatible
with the coalgebra structures on both sides.
\end{theorem}

\begin{proof}
The operadic bar construction
$B_{\mathcal{P}}(\mathcal{A}) =
(\mathcal{P}^{!,c} \circ \bar{\mathcal{A}}, d_{\bar{B}})$
is a cofree $\mathcal{P}^!$-coalgebra on the augmentation ideal
$\bar{\mathcal{A}} = \ker(\varepsilon)$, with differential $d_{\bar{B}}$
encoding the algebra structure of $\mathcal{A}$ via the twisting
morphism $\kappa: \mathcal{P}^{!,c} \to \mathcal{P}$
\cite[Theorem~6.5.7]{LV12}.

\emph{Step 1} (Geometric realization of $\mathcal{P}^{!,c}$).
For $\mathcal{P} = \operatorname{Com}$ (the $\Einf$-chiral case),
$\mathcal{P}^{!,c} = \operatorname{Lie}^c$, and the cooperad
$\operatorname{Lie}^c(n)$ is computed by the logarithmic de~Rham
model of the open configuration space $C_n(X)$ on the
Fulton--MacPherson compactification
$\overline{C}_n(X)$. On the affine type-$A$ chart, the Arnold
relations in the Orlik--Solomon algebra give the Lie cooperad.
For $\mathcal{P} = \operatorname{Ass}$,
$\mathcal{P}^{!,c} = \operatorname{Ass}^c$ (the associative
cooperad), realized by ordered configurations.

\emph{Step 2} (Identification of differentials).
The operadic differential $d_{\bar{B}}$ decomposes into $d_1$ (the cooperadic cocomposition $\mathcal{P}^{!,c}(n)
\xrightarrow{\Delta} \mathcal{P}^{!,c}(k)
\otimes \mathcal{P}^{!,c}(n_1) \otimes \cdots
\otimes \mathcal{P}^{!,c}(n_k)$) and $d_2$ (the twisting morphism applied to pairs of consecutive inputs).
Under the geometric realization, $d_1$ corresponds to the de~Rham
differential $d_{\mathrm{dR}}$ on $\Omega^*(\log D)$ (which detects
the codimension-one boundary strata of $\overline{C}_n(X)$), and
$d_2$ corresponds to the residue differential $d_{\mathrm{res}}$
(which extracts OPE data at collision divisors). The sign
  comparison of Lemma~\ref{lem:LV-sign-comparison} shows these agree up to
the suspension isomorphism
$s\bar{\mathcal{A}} \cong \bar{\mathcal{A}} \otimes
\Omega^1(\log D)|_{\text{codim-1}}$.

\emph{Step 3} (Quasi-isomorphism).
The comparison map sends a geometric bar element
$a_1 \otimes \cdots \otimes a_n \otimes \omega
\in \bar{B}_{\mathrm{geom}}^n(\mathcal{A})$
to the operadic bar element
$[sa_1 | \cdots | sa_n] \in B_{\mathcal{P}}(\mathcal{A})_n$
by evaluating the logarithmic form $\omega$ at the fundamental class
of the fiber. This evaluation is a chain map by Stokes's theorem
(boundary terms match the residue differential by Step~2).
It is a quasi-isomorphism by the logarithmic comparison theorem:
the complex $\Omega^*_{\overline{C}_{n+1}(X)}(\log D)$ computes
$H^*(C_{n+1}(X); \mathbb{C})$ (the cohomology of the open configuration
space), and the operadic bar differential encodes precisely the
combinatorics of the boundary divisor stratification of
$\overline{C}_{n+1}(X)$; both are controlled by the operad structure
on $\{C_{n+1}(X)\}_{n \geq 0}$
(cf.\ \cite[Theorem~3.7.4]{BD04}, \cite[Theorem~5.1]{Get95}).
\end{proof}

\begin{remark}[Lurie comparison]
Lurie's bar construction for an augmented algebra is the
geometric realization of the two-sided simplicial bar object
\[
\mathbf{1} \leftarrow A \leftarrow A^{\otimes 2}
\leftarrow A^{\otimes 3} \leftarrow \cdots,
\]
equivalently the relative tensor product
$\mathbf{1}\otimes_A \mathbf{1}$
\textup{(}Construction~4.4.2.7, Theorem~4.4.2.8, and
\S5.2.1 of \cite{HA}\textup{)}.
Our symmetric chiral bar $\barB_X(\cA)$ is the curve-valued
factorization analogue of that reduced augmented bar resolution:
ordinary tensor powers are replaced by configuration-space
factorization data on $\Ran(X)$, while the ordered tensor coalgebra
$B_X^{\mathrm{ord}}(\cA)=T^c(s^{-1}\bar\cA)$ remains the primitive
$\Eone$ lift.

It is therefore inaccurate to identify $\barB_X(\cA)$ itself with the
\v{C}ech nerve of the augmentation
$\varepsilon\colon \cA \to \omega_X$.
The \v{C}ech complex that appears in this chapter is the Weiss-cover
descent complex used to compute factorization homology
\textup{(}Theorem~\textup{\ref{thm:bar-NAP-homology}}\textup{)}.
That is a cover-theoretic resolution of the same homological output,
not the definition of the bar coalgebra.
In the connected group-object case treated in \cite[\S5.2.6]{HA},
the same simplicial pattern becomes a delooping nerve.
That extra identification is not part of the present chiral bar
formalism.
\end{remark}

\begin{theorem}[Arbitrary-mode ordered residue formula;
\ClaimStatusProvedHere]\label{thm:residue-formula}
\index{residue!iterated}
\index{OPE modes!arbitrary-mode bar residue}
Type signature: \textup{(}Open quadrant, ordered logarithmic
Fulton--MacPherson presentation, Beilinson level~$2$, hypothesis
package: finite OPE window in a dg chiral algebra, ordered
boundary-orientation convention
\ref{conv:orientations-enhanced}, and bar-sign convention
\ref{thm:bar-sign-coherence}\textup{)}.
Let \(V\subset\bar\cA\) be a finite homogeneous window closed under
the finitely many singular products \(a_{(m)}b\) occurring in the
window. For homogeneous \(a_1,\ldots,a_N\in V\), write
\([a_1|\cdots|a_N]=s^{-1}a_1\otimes\cdots\otimes s^{-1}a_N\).
Let \(D_I\), \(I=\{i,j\}\) with \(i<j\), be a pairwise collision
divisor and write
a logarithmic form near \(D_I\) as
\[
\omega=\eta_I\wedge\alpha+\beta,\qquad
\eta_I=d\log(z_i-z_j),
\]
with \(\alpha,\beta\) logarithmic forms having no normal
\(d\log(z_i-z_j)\)-factor. The \(m\)-th OPE-mode contribution of the
ordered residue differential along \(D_I\) is
\begin{equation}\label{eq:ordered-residue-arbitrary-mode}
d_{(m),I}\bigl([a_1|\cdots|a_N]\otimes\omega\bigr)
=
(-1)^{\epsilon_B(I;a_\bullet,\omega)}
[a_1|\cdots|a_i{}_{(m)}a_j|\cdots|
\widehat{a_j}|\cdots|a_N]\otimes\alpha|_{D_I}.
\end{equation}
The parity \(\epsilon_B(I;a_\bullet,\omega)\) is the unique parity
fixed by Theorem~\ref{thm:bar-sign-coherence}; for the adjacent
ordered face \(I=\{i,i+1\}\) and \(\omega=\eta_{i,i+1}\wedge\alpha\),
it is
\[
\epsilon_B(I;a_\bullet,\omega)
\equiv
\sum_{q<i}(|a_q|-1)+|a_i|
\pmod 2,
\]
the exponent in \eqref{eq:ordered-bar-dres}. For non-adjacent FM
faces, the same formula is multiplied by the block-permutation and
boundary-orientation parity of
Convention~\ref{conv:orientations-enhanced}.
If the normal factor \(\eta_I\) is absent, then
\(\operatorname{Res}^{\mathrm{form}}_{D_I}(\omega)=0\), hence the
summand is zero.

The cases \(m=0\), \(m=1\), and \(m\ge2\) are respectively the
simple-pole bracket, double-pole, and higher-pole contributions. The
Poincar\'e residue removes only the logarithmic form factor
\(\eta_I\); it never imposes \(m=0\). The pole order is selected
solely by the mode projection of the Beilinson--Drinfeld chiral
operation:
\[
\operatorname{pr}_m\mu_{\mathrm{BD}}(a_i,a_j)
=a_i{}_{(m)}a_j .
\]
\end{theorem}

\begin{proof}
Near \(D_I\), put \(t=z_i-z_j\). Every logarithmic form has the
normal decomposition
\[
\omega=d\log t\wedge\alpha+\beta,\qquad
\operatorname{Res}^{\mathrm{form}}_{D_I}(\omega)=\alpha|_{D_I}.
\]
The BD chiral product decomposes on the chosen OPE window as the
finite sum of mode projections
\[
\mu_{\mathrm{BD}}=\sum_{m\ge0}\operatorname{pr}_m\mu_{\mathrm{BD}},
\qquad
\operatorname{pr}_m\mu_{\mathrm{BD}}(a_i,a_j)=a_i{}_{(m)}a_j.
\]
Thus the boundary operation obtained by first taking the logarithmic
form residue and then applying the \(m\)-th mode projection is exactly
\eqref{eq:ordered-residue-arbitrary-mode}. The sign is not a new
convention: for adjacent collisions it is the exponent of
\eqref{eq:ordered-bar-dres}; for non-adjacent FM faces it is the same
desuspended bar sign together with the block and boundary-orientation
parity fixed in Convention~\ref{conv:orientations-enhanced}. Since
the decomposition of \(\mu_{\mathrm{BD}}\) is over all singular modes
in the finite window, no higher pole can be lost by the form-residue
operation.
\end{proof}

\begin{theorem}[BD chiral operation and full OPE residue;
\ClaimStatusProvedHere]
\label{thm:bd-ope-residue-full-poles}
\index{Beilinson--Drinfeld!chiral operation and bar residue}
\index{OPE residue!full pole retention}
\index{bar differential!arbitrary OPE mode}
Type signature: \textup{(}Open quadrant, symmetric
Beilinson--Drinfeld/Ran presentation compared with the ordered
logarithmic Fulton--MacPherson presentation, Beilinson level~$2$,
hypothesis package: chiral algebra locality, ordered-to-symmetric
comparison only after a separate descent datum, and the sign package
of Theorem~\ref{thm:bar-sign-coherence}\textup{)}.
On the formal neighbourhood of a pairwise diagonal, the OPE residue
used in the ordered logarithmic bar is the Beilinson--Drinfeld chiral
operation
\[
\mu_{\mathrm{BD}}\colon
j_*j^*(\cA\boxtimes\cA)\longrightarrow \Delta_!\cA
\]
written on the Fulton--MacPherson boundary. In a local coordinate
normal to the diagonal,
\[
a(z)b(w)\sim
\sum_{m\ge0}(a_{(m)}b)(w)(z-w)^{-m-1}
\quad\Longleftrightarrow\quad
\operatorname{pr}_m\mu_{\mathrm{BD}}(a,b)=a_{(m)}b .
\]
This theorem is the ordered boundary statement before taking
coinvariants. No singular OPE mode is discarded: for arbitrary tensor
length and arbitrary pole order, the contribution of the $m$-th mode
\(a_{(m)}b\) is exactly the summand
\eqref{eq:ordered-residue-arbitrary-mode} of
Theorem~\ref{thm:residue-formula}. Higher poles are therefore part of
$d_{\mathrm{res}}$; only an additional current-window hypothesis may
project them to a scalar trace or remove them from a specific
comparison. This bar-residue statement is not the polar connection
kernel: the logarithmic form supplies the normal orientation and does
not multiply the OPE pole by a second propagator pole.
\end{theorem}

\begin{proof}
The chiral product in the sense of Beilinson--Drinfeld
\cite[\S\S3.3--3.4]{BD04} is the $\cD$-module morphism from the
extension across the complement of the diagonal to the diagonal
pushforward. Its local-coordinate expansion along the diagonal is the
singular OPE expansion, so the \(m\)-th normal coefficient is
\(\operatorname{pr}_m\mu_{\mathrm{BD}}(a,b)=a_{(m)}b\). The FM
boundary divisor $D_{ij}$ is the normal-crossings resolution of that
same diagonal collision. Applying the Poincar\'e residue to the
logarithmic normal form and then applying \(\mu_{\mathrm{BD}}\) is
therefore the boundary model for the BD operation. The ordered complex
keeps the labelled \(E_1\) lift before quotienting.

The arbitrary-mode formula is
\eqref{eq:ordered-residue-arbitrary-mode}. The form residue removes
only the logarithmic form factor $\eta_{ij}$.
The pole order is selected separately by the OPE-mode projection
$\operatorname{pr}_m\mu_{\mathrm{BD}}\), so the construction includes
every $m\ge0$ term. Proposition~\ref{prop:pole-decomposition} records
the same fact at the level of the differential
$d_{\mathrm{res}}=\sum_{m\ge0}d_{(m)}$. The signs in the displayed
formula are those fixed by Theorem~\ref{thm:bar-sign-coherence}.
\end{proof}

\begin{corollary}[Symmetric BD residue after descent;
\ClaimStatusConditional]
\label{cor:bd-ope-symmetric-ran-differential}
\index{Beilinson--Drinfeld!symmetric residue after descent}
\index{symmetric bar!conditional BD differential}
Type signature: \textup{(}Open quadrant, symmetric
Beilinson--Drinfeld/Ran presentation, Beilinson level~$2$, hypothesis
package: the ordered residue theorem
\ref{thm:bd-ope-residue-full-poles} plus the ordered-to-symmetric
descent datum of
Proposition~\ref{prop:symmetric-bar-descent-criterion}\textup{)}.
If the descent criterion of
Proposition~\ref{prop:symmetric-bar-descent-criterion} is satisfied,
then the symmetric quotient of the ordered residue differential is the
Cousin/factorisation differential on the Beilinson--Drinfeld Ran
coalgebra. Without that descent datum this quotient is only a graded
coinvariant shadow and is not asserted to be a dg factorisation
coalgebra.
\end{corollary}

\begin{proof}
Theorem~\ref{thm:bd-ope-residue-full-poles} identifies the ordered
Fulton--MacPherson boundary residue with the BD chiral operation and
retains all pole modes. Proposition~\ref{prop:symmetric-bar-descent-criterion}
is exactly the compatibility condition that makes the
$\Sigma_\bullet$-coinvariant differential, de~Rham differential, and
coalgebra structure descend to the unordered Ran quotient. Under that
condition the quotient of the ordered residue is therefore the
BD Ran-side Cousin differential. If the condition is absent, the same
coinvariant object has no asserted differential compatible with
factorisation, so the symmetric statement stops at the formal graded
quotient.
\end{proof}

\subsection{Uniqueness and functoriality}

\begin{theorem}[Uniqueness and functoriality; \ClaimStatusProvedElsewhere]\label{thm:bar-uniqueness-functoriality}
The geometric bar construction is the unique functor
\[\bar{B}_{geom}: \text{ChirAlg}_X \to \text{dgCoalg}\]
satisfying:
\begin{enumerate}
\item \emph{Locality:} For $j: U \hookrightarrow X$ open, $j^*\bar{B}_{geom}(\mathcal{A}) \cong \bar{B}_{geom}(j^*\mathcal{A})$
\item \emph{External product:} $\bar{B}_{geom}(\mathcal{A} \boxtimes \mathcal{B}) \cong \bar{B}_{geom}(\mathcal{A}) \boxtimes \bar{B}_{geom}(\mathcal{B})$
\item \emph{Normalization:} $\bar{B}_{geom}(\mathcal{O}_X) = \Omega^*(\overline{\mathcal{C}}_{*+1}(X))$
\end{enumerate}
up to unique natural isomorphism.

It defines a functor from chiral algebras to filtered conilpotent chiral coalgebras, with essential image the coalgebras having logarithmic coderivations supported on collision divisors.
\end{theorem}

Uniqueness follows from the universal property of the cofree conilpotent coalgebra and the locality axiom; functoriality is established in the preceding construction.

\begin{definition}[Conilpotent chiral coalgebra]\label{def:conilpotent-chiral-coalgebra}
A chiral coalgebra $C$ is \emph{filtered conilpotent} if the iterated comultiplication
$\Delta^{(n)} : C \to C^{\otimes(n+1)}$ satisfies: For each $c \in C$, there exists
$N$ such that $\Delta^{(n)}(c) = 0$ for all $n \geq N$. This ensures the cobar
construction $\Omega^{\text{ch}}(C)$ is well-defined without completion.
\end{definition}



\begin{proof}
\emph{Step 1: Existence.} We verify each axiom explicitly:
\begin{itemize}
\item \emph{Locality:} For $j: U \hookrightarrow X$ open,
we have $C_n(U) = j^{-1}(C_n(X))$.
The maximal extension $j_*j^*$ commutes with sections over
configuration spaces:
\begin{align*}
j^*\bar{B}_{\text{geom}}(A)
&= j^*\Gamma(\overline{C}_{n+1}(X), \cdots) \\
&= \Gamma(\overline{C}_{n+1}(U), \cdots)
= \bar{B}_{\text{geom}}(j^*A)
\end{align*}

\item \emph{External product:} For disjoint open subsets
$U_1, U_2 \subset X$, the identification
$C_n(U_1 \sqcup U_2) \cong
\bigsqcup_{k+l=n} C_k(U_1) \times C_l(U_2)$
induces a natural isomorphism
\[
\bar{B}_{\mathrm{geom}}(\mathcal{A}|_{U_1 \sqcup U_2})
\cong \bigoplus_{k+l=n}
\bar{B}_k(\mathcal{A}|_{U_1}) \otimes
\bar{B}_l(\mathcal{A}|_{U_2}),
\]
compatible with boundary stratifications.

\item \emph{Normalization:} For $A = \mathcal{O}_X$, there are no nontrivial OPEs, so
$d_{\text{fact}} = 0$, and we are left with just the de Rham complex on configuration spaces.
\end{itemize}

\emph{Step 2: Uniqueness.} Let $F, G$ be two such functors.

For the structure sheaf: By normalization,
\[F(\mathcal{O}_X) = G(\mathcal{O}_X) = \Omega^*(\overline{\mathcal{C}}_{*+1}(X))\]

For the free chiral algebra $\text{Free}^{\text{ch}}(V)$ on a vector bundle $V$, the locality and external product axioms determine:
\[F(\text{Free}^{\text{ch}}(V)) \cong \text{Sym}^*(s^{-1}V) \otimes \Omega^*(\overline{C}_{*+1}(X))\]
(here $s^{-1}V$ is the desuspension $|s^{-1}v| = |v| - 1$ of Convention~\ref{conv:bar-coalgebra-identity}, matching the desuspension in $T^c(s^{-1}\bar{\mathcal{A}})$), and similarly for $G$, giving a canonical isomorphism $\eta_V\colon F(\text{Free}^{\text{ch}}(V)) \xrightarrow{\sim} G(\text{Free}^{\text{ch}}(V))$.


\begin{align}
F(\text{Free}^{\text{ch}}(V)) &= F(V^{\otimes_{\text{ch}} \bullet})\\
&\cong F(V)^{\otimes \bullet} \quad \text{(external product)}\\
&\cong (s^{-1}V \otimes F(\mathcal{O}_X))^{\otimes \bullet} \quad \text{(locality)}\\
&\cong \text{Sym}^*(s^{-1}V) \otimes \Omega^*(\overline{\mathcal{C}}_{*+1}(X))
\end{align}

Similarly for $G$, giving a canonical isomorphism $\eta_{V}\colon F(\text{Free}^{\text{ch}}(V)) \xrightarrow{\sim} G(\text{Free}^{\text{ch}}(V))$.

For general $\mathcal{A} = \text{Free}^{\text{ch}}(V)/R$:
The relations $R$ determine boundaries via the same residue formulas in both $F(A)$ and $G(A)$:
\begin{itemize}
\item Each relation $r \in R$ maps to $d_{\text{fact}}(r)$ computed via residues
\item The residue formula is determined by the OPE structure
\item Locality ensures these agree on all affine charts
\end{itemize}

\emph{Step 3: Natural isomorphism.}
For morphism $\phi: \mathcal{A} \to \mathcal{B}$, the diagram
\[
\begin{tikzcd}
F(\mathcal{A}) \arrow[r, "\eta_\mathcal{A}"] \arrow[d, "F(\phi)"] & G(\mathcal{A}) \arrow[d, "G(\phi)"]\\
F(\mathcal{B}) \arrow[r, "\eta_\mathcal{B}"] & G(\mathcal{B})
\end{tikzcd}
\]
commutes by construction of $\eta$ using universal properties.

\emph{Verification that relations map to boundaries}: Let $r \in R \subset \text{Free}^{\text{ch}}(V) \otimes \text{Free}^{\text{ch}}(V)$.
Under $F$, we have:
\[F(r) \in F(\text{Free}^{\text{ch}}(V) \otimes \text{Free}^{\text{ch}}(V)) = F(\text{Free}^{\text{ch}}(V))^{\otimes 2}\]
\[ = (V[1] \otimes \Omega^*(C_{*+1}(X)))^{\otimes 2}\]
The differential $d_F$ maps $r$ to the boundary because:
\[d_F(r) = d_{\text{fact}}(r) + d_{\text{config}}(r) + d_{\text{int}}(r)\]
where $d_{\text{fact}}$ implements the relation via residue extraction. Similarly for $G$.
The agreement $F(r) = G(r)$ in cohomology follows from the universal property
of free chiral algebras and the uniqueness of residue extraction.

\emph{Step 4: Uniqueness of isomorphism.}
Any other natural isomorphism $\eta': F \Rightarrow G$ must agree on $\mathcal{O}_X$ by normalization,
hence on free algebras by external product, hence on all algebras by locality.
\end{proof}

\begin{theorem}[Ordered bar as a complete conilpotent functor;
\ClaimStatusProvedHere]
\label{thm:ordered-bar-complete-conilpotent-functor}
\index{ordered bar!complete conilpotent functor}
\index{complete conilpotent dg coalgebra!ordered bar}
Type signature: \textup{(}Open quadrant, ordered logarithmic
Fulton--MacPherson/Ran presentation, Beilinson level~$2$, hypothesis
package: augmented complete finite-window $\Eone$-chiral algebra,
holonomic/excision descent on disjoint opens, finite output in every
fixed word-length window, and the sign package of
Theorem~\ref{thm:bar-sign-coherence}\textup{)}.
Let
\[
\widehat{\mathbb B}^{\ord}_X(\cA)
        :=\prod_{n\geq 0}\mathbb B^{\ord}_{X,n}(\cA)
\]
be the completed ordered bar carrier of
Definition~\ref{def:ordered-chiral-bar-term-log}. The assignment
\[
\cA\longmapsto
\bigl(\widehat{\mathbb B}_X^{\ord}(\cA),d_B,\Delta_{\mathrm{dec}},
\varepsilon_B,F^\bullet\bigr)
\]
extends to a functor from augmented $\Eone$-chiral algebras on $X$ to
complete conilpotent dg coalgebras on the ordered Ran space:
\[
B_X^{\ord}\colon
\mathsf{ChirAlg}^{\Eone,\mathrm{aug,comp}}_X
\longrightarrow
\mathsf{dgCoalg}^{\mathrm{conil,comp}}(\Ran^{\ord}(X)).
\]
More precisely:
\begin{enumerate}
\item the differential is
\[
d_B=d_\cA+d_{\mathrm{dR}}+d_{\mathrm{res}},
\]
satisfies \(d_B^2=0\), and is a coderivation for the
deconcatenation coproduct \(\Delta_{\mathrm{dec}}\);
\item \(\varepsilon_B\) is projection to word length~\(0\), the
coaugmentation is the unit word, and the reduced augmentation ideal
removes all positive-length vacuum degeneracies;
\item the filtration \(F^\bullet\) is the completed word-length
filtration, every finite word-length quotient is conilpotent under
the reduced deconcatenation coproduct, and the product completion is
therefore complete pro-conilpotent;
\item for a morphism \(f\colon\cA\to\cA'\), the induced coalgebra map is
the tensorwise map \(s^{-1}\bar f\) on augmentation-ideal factors and
the identity on the logarithmic Fulton--MacPherson form factor.
\end{enumerate}
This is an ordered statement. It asserts neither descent to the
symmetric Ran quotient nor a global KZ/Arnold comparison; those require
separate locality and genus hypotheses.
\end{theorem}

\begin{proof}
For an augmented complete finite-window \(\Eone\)-chiral algebra
\(\cA\), Definition~\ref{def:ordered-chiral-bar-term-log} gives the
aritywise logarithmic Fulton--MacPherson \(\cD\)-module
\(\mathbb B^{\ord}_{X,n}(\cA)\). Taking the product over word length
gives \(\widehat{\mathbb B}^{\ord}_X(\cA)\), and the finite-output
hypothesis ensures that the internal differential, de Rham
differential, and residue/OPE collision differential are continuous
for this completion.

Proposition~\ref{prop:ordered-bar-local-differential-identities}
computes the three summands
\[
d_\cA,\qquad d_{\mathrm{dR}},\qquad d_{\mathrm{res}}
\]
on ordered logarithmic charts and proves the identities
\[
d_\cA^2=d_{\mathrm{dR}}^2=d_{\mathrm{res}}^2=0,\qquad
d_\cA d_{\mathrm{dR}}+d_{\mathrm{dR}}d_\cA=
d_\cA d_{\mathrm{res}}+d_{\mathrm{res}}d_\cA=
d_{\mathrm{dR}}d_{\mathrm{res}}+d_{\mathrm{res}}d_{\mathrm{dR}}=0.
\]
Hence \(d_B=d_\cA+d_{\mathrm{dR}}+d_{\mathrm{res}}\) is square-zero.
The same proposition checks that \(d_B\) is a coderivation for
ordered deconcatenation; Theorem~\ref{thm:bar-sign-coherence} supplies
the global sign compatibility for passing between ordered charts.
Proposition~\ref{prop:ordered-bar-carrier-ran-factorisation} supplies
the factorisation restrictions on disjoint ordered opens, so the
coalgebra object lives on \(\Ran^{\ord}(X)\).

The counit is projection to the empty word and the coaugmentation is
inclusion of the empty word. Passing to the reduced augmentation ideal
removes positive-length words containing a vacuum factor by the
normalization in
Proposition~\ref{prop:ordered-bar-local-differential-identities}. The
word-length filtration is complete and is preserved by \(d_B\) and by
\(\Delta_{\mathrm{dec}}\) by
Convention~\ref{conv:bar-length-filtration}. On a word of length
\(n\), the \(N\)-fold reduced deconcatenation is zero for \(N>n\);
therefore every finite word-length quotient is conilpotent. The full
product object is the inverse limit of these conilpotent finite
quotients, which is the asserted complete pro-conilpotence.

Let \(f\colon\cA\to\cA'\) be a morphism in
\(\mathsf{ChirAlg}^{\Eone,\mathrm{aug,comp}}_X\). Since \(f\) preserves
augmentations, it restricts to \(\bar f\colon\bar\cA\to\bar\cA'\). The
map on word length \(n\) is
\[
(s^{-1}\bar f)^{\otimes n}\otimes
\mathrm{id}_{\Omega^\bullet_{\log}(\FM_X(n))},
\]
with the evident ordered Ran pullback on the carrier. Preservation of
the internal differential gives compatibility with \(d_\cA\);
preservation of the \(\Eone\)-chiral OPE products gives compatibility
with \(d_{\mathrm{res}}\); the identity on the logarithmic
Fulton--MacPherson form factor gives compatibility with
\(d_{\mathrm{dR}}\). Tensorwise application also commutes with
deconcatenation, the empty-word counit, the coaugmentation, and
\(F^\bullet\). Identities and composition are checked in each word
length, hence the construction is functorial.

No quotient by the symmetric group is used in the argument, and no
KZ/Arnold pullback identity is invoked. The theorem therefore remains
strictly ordered and finite-window complete.
\end{proof}

\subsection{Bar complex as chiral coalgebra}

\begin{theorem}[Bar complex is chiral; \ClaimStatusConditional]\label{thm:bar-chiral}
\index{factorization coalgebra|textbf}
The geometric bar complex $\bar{B}^{\text{ch}}(\mathcal{A})$ naturally carries the structure of a differential graded chiral coalgebra.
The coproduct is the coshuffle coproduct on $\mathrm{Sym}^c(s^{-1}\bar{\mathcal{A}})$; the ordered refinement on $T^c(s^{-1}\bar{\mathcal{A}})$ with deconcatenation coproduct (Theorem~\textup{\ref{thm:three-bar-complexes}}\textup{(i)}) carries strictly richer data, including the $R$-matrix tower.
\end{theorem}

\begin{proof}
We construct the chiral coalgebra structure explicitly:

\emph{Comultiplication.} The map $\Delta: \bar{B}^{\text{ch}}(\mathcal{A}) \to \bar{B}^{\text{ch}}(\mathcal{A}) \otimes \bar{B}^{\text{ch}}(\mathcal{A})$ is induced by:
\[
\Delta: \overline{C}_{n+1}(X) \to \bigcup_{I \sqcup J = [n+1]} \overline{C}_{|I|}(X) \times \overline{C}_{|J|}(X)
\]
where the union is over ordered partitions with $0 \in I$. Explicitly:
\[
\Delta(\phi_0 \otimes \cdots \otimes \phi_n \otimes \omega) = \sum_{I \sqcup J} \pm \left(\bigotimes_{i \in I} \phi_i \otimes \omega|_I\right) \otimes \left(\bigotimes_{j \in J} \phi_j \otimes \omega|_J\right)
\]

\emph{Counit.} $\epsilon: \bar{B}^{\text{ch}}(\mathcal{A}) \to \mathbb{C}$ is given by projection onto degree 0:
\[
\epsilon(\phi_0 \otimes \cdots \otimes \phi_n \otimes \omega) = \begin{cases}
\int_X \phi_0 & \text{if } n = 0 \\
0 & \text{if } n > 0
\end{cases}
\]

\emph{Coassociativity.} Follows from the associativity of configuration space stratifications
(the stratification $\overline{C}_{n_1+n_2}(X) \supset C_{n_1}(X) \times C_{n_2}(X)$ is associative in the sense that iterated boundary extraction $(I_1, I_2, I_3) \mapsto (I_1 \mid I_2 \mid I_3)$ is independent of the order of extraction, by the transversality of the strata):
\[
(\Delta \otimes \text{id}) \circ \Delta = (\text{id} \otimes \Delta) \circ \Delta
\]

\emph{Compatibility with differential.} The comultiplication is a chain map:
\[
\Delta \circ d = (d \otimes \text{id} + \text{id} \otimes d) \circ \Delta
\]
This follows from the compatibility of residues with the stratification of configuration spaces
(the residue of $d\omega$ along a boundary stratum equals the residue of $\omega$ on the induced stratum, by the compatibility of the Leray residue with the boundary map in the long exact sequence of the pair).
\end{proof}

\medskip\noindent
Three inputs have now produced one output. The augmentation ideal of an
augmented chiral algebra, the Fulton--MacPherson compactification of its
configuration spaces, and the logarithmic kernel
$\eta_{ij}=d\log(z_i-z_j)$ assemble into the dg factorization coalgebra
$\barB_X(\cA)$, with the Arnold relation enforcing $\dzero^2=0$ at
genus~$0$. The symmetric bar is the universal-conductor descent of the
primitive ordered bar
$B^{\mathrm{ord}}_X(\cA)=T^c(s^{-1}\bar\cA)$: ordinary
$\Sigma_n$-coinvariants in the untwisted characteristic-zero splitting,
and $L_R$-twisted derived completed coinvariants in the braided case.
On the PBW/Koszul $E_\infty$ surface this descent preserves the
cohomology used in Theorem~H; outside that surface it is only the
averaging map and need not preserve the ordered cohomology. The
cobar--bar adjunction
(Chapter~\ref{chap:cobar-construction}) inverts this construction on
the Koszul locus; the fiberwise curvature
$\dfib^{\,2}=\kappa(\cA)\cdot\omega_g$
(Chapter~\ref{chap:higher-genus}) records its higher-genus anomaly. The
characteristic datum $\kappa(\cA)$ is the residue of that failure at the
branch point.

\section{Ordered K3 convolution}
\label{sec:barconstr-supplement}

\begin{definition}[Hall-theoretic ordered bar differential; \ClaimStatusDefinitional]
\label{rem:barconstr-1}
The ordered bar construction on the K3 chiral bialgebra
$\mathbf H_{\Delta_5}$ takes the form
\[
B^{E_1}_{\mathrm{ch}}(\mathbf H_{\Delta_5}) \;=\; T^c\bigl(s^{-1}\overline{\mathbf H_{\Delta_5}}\bigr) \;\in\; \mathrm{Hall}(\mathrm{CoHA}_{K3\times E}),
\]
the tensor coalgebra on the desuspended augmentation-reduction
$\overline{\mathbf H_{\Delta_5}} = \ker(\epsilon:\mathbf H_{\Delta_5}\to\mathbb C)$,
taken in the Hall subcategory of the K3-elliptic Cohomological Hall
Algebra. In the Hall model the bar differential is written
$d_B = d_{\mathrm{Hall}} + d_{\mathrm{Sieg}} + \hbar\, d_{\mathrm{assoc}} + \hbar^2\, d_{\Phi_{10}^{\mathrm{un}}/\eta^{24}}$
as a decomposition by geometric source. The Hall piece
$d_{\mathrm{Hall}}$ distinguishes the Hall-theoretic bar from classical
bar-construction theory: in the classical case $d_{\mathrm{Hall}} = 0$
because the algebra lives in a symmetric monoidal category, but the Hall
category is \emph{only} monoidal (non-symmetric), and the Hall piece
is genuinely non-trivial. This is the source of the non-abelian chiral
structure of $\mathbf H_{\Delta_5}$. Primary-lit: Schiffmann--Vasserot
2012, Kontsevich--Soibelman 2011 ($\mathrm{CoHA}$), Davison 2017 (CY-3
extension).
\end{definition}

\begin{definition}[Ordered convolution dGLA $\mathfrak{C}_{\mathbf{H}_{\Delta_5}}$;
\ClaimStatusDefinitional]
\label{def:convolution-dgla-k3}
Let $A = \mathbf{H}_{\Delta_5}$ with product
$\mu\colon A\otimes A\to A$, augmentation
$\epsilon\colon A\to \mathbb{C}$, and reduced algebra
$\bar A = \ker\epsilon$. Let
$B^{\mathrm{ord}}(A) = T^c(s^{-1}\bar A)$ be the ordered bar coalgebra,
equipped with the reduced comultiplication
$\bar\Delta\colon B^{\mathrm{ord}}(A) \to B^{\mathrm{ord}}(A)\otimes
B^{\mathrm{ord}}(A)$,
$\bar\Delta(s^{-1}x_1\otimes\cdots\otimes s^{-1}x_n) =
\sum_{k=1}^{n-1} (s^{-1}x_1\otimes\cdots\otimes s^{-1}x_k)\otimes
(s^{-1}x_{k+1}\otimes\cdots\otimes s^{-1}x_n)$
(deconcatenation).
The \emph{ordered convolution dGLA} of $\mathbf{H}_{\Delta_5}$ is the
graded vector space
\[
 \mathfrak{C}_{\mathbf{H}_{\Delta_5}}
 \;:=\;
 \mathrm{Hom}_{\Bbbk}\!\bigl(B^{\mathrm{ord}}(\mathbf{H}_{\Delta_5}),\,
 \mathbf{H}_{\Delta_5}\bigr),
 \qquad
 \mathfrak{C}^p_{\mathbf{H}_{\Delta_5}}
 \;=\;
 \prod_{n\ge 0}
 \mathrm{Hom}^p\!\bigl(B^{\mathrm{ord}}_n(\mathbf{H}_{\Delta_5}),\,
 \mathbf{H}_{\Delta_5}\bigr),
\]
working in the weight-completed pro-category with respect to the
Mukai-lattice grading $\Lambda_{\mathrm{Muk}}$.
Let $d_B = d_{\mathrm{Hall}} + d_{\mathrm{Sieg}} + \hbar d_{\mathrm{assoc}}
+ \hbar^2 d_{\Phi_{10}^{\mathrm{un}}/\eta^{24}}$ be the bar differential of
Remark~\ref{rem:barconstr-1} and let $d_A$ be the BKM differential on
$\mathbf{H}_{\Delta_5}$ (Chevalley--Eilenberg differential of the
completed Lie bialgebra
$(\mathfrak{g}_{\Delta_5},\delta_{\mathrm{GN}})$ lifted, after a
chosen associator and pro-\(\hbar\)-adic topology are fixed, to the
specified Etingof--Kazhdan super-quasi-Hopf target; Borcherds 1988
\S 4).
\begin{enumerate}[label=\textup{(\alph*)}]
\item \emph{Differential} $D\colon
 \mathfrak{C}^p\to\mathfrak{C}^{p+1}$, defined on $\phi\in\mathfrak{C}^p$
 and $w\in B^{\mathrm{ord}}_n$ by
 \[
 (D\phi)(w) \;=\; d_A\bigl(\phi(w)\bigr)
 \;-\; (-1)^{p}\,\phi\bigl(d_B w\bigr).
 \]
\item \emph{Convolution bracket} $[\cdot,\cdot]\colon
 \mathfrak{C}^p\otimes\mathfrak{C}^q\to\mathfrak{C}^{p+q}$, defined
 on $\phi\in\mathfrak{C}^p$, $\psi\in\mathfrak{C}^q$, $w\in
 B^{\mathrm{ord}}$ by
 \[
 [\phi,\psi](w) \;=\; \mu\circ(\phi\otimes\psi)\circ \bar\Delta(w)
 \;-\; (-1)^{pq}\,
 \mu\circ(\psi\otimes\phi)\circ \bar\Delta(w).
 \]
 Equivalently, $[\phi,\psi] = \phi\smile\psi - (-1)^{pq}\psi\smile\phi$
 with convolution cup product
 $\phi\smile\psi := \mu\circ(\phi\otimes\psi)\circ\bar\Delta$.
\end{enumerate}
This is the ordered convolution datum attached to the K3 Hall bar
coalgebra; Proposition~\ref{prop:dgla-axioms-k3-convolution} verifies
that it is a differential graded Lie algebra.
\end{definition}

\begin{proposition}[dGLA axioms for
$\mathfrak{C}_{\mathbf{H}_{\Delta_5}}$; \ClaimStatusProvedHere]
\label{prop:dgla-axioms-k3-convolution}
The structure of Definition~\ref{def:convolution-dgla-k3} satisfies:
\textup{(i)} $D^2 = 0$; \textup{(ii)} $[\cdot,\cdot]$ is super-antisymmetric,
$[\phi,\psi] = -(-1)^{pq}[\psi,\phi]$; \textup{(iii)} the super-Jacobi
identity holds; \textup{(iv)} $D$ is a derivation of the bracket.
\end{proposition}

\begin{proof}
\emph{(i) $D^2 = 0$.} Expand:
$(D^2\phi)(w) = d_A^2\phi(w) - (-1)^{p+1}d_A\phi(d_B w)
- (-1)^p d_A\phi(d_B w) + \phi(d_B^2 w) = 0$,
where $d_A^2 = 0$ by Jacobi for
$(\mathfrak{g}_{\Delta_5},\delta_{\mathrm{GN}})$ in the fixed
completed Etingof--Kazhdan super-quasi-Hopf target
(Borcherds 1988 \S 4, Theorem 2; Kac 1990 \S 2.5), $d_B^2 = 0$ by
coassociativity of $\bar\Delta$ combined
with the Arnold relation
$\eta_{ij}\wedge\eta_{jk} + \eta_{jk}\wedge\eta_{ki} +
\eta_{ki}\wedge\eta_{ij} = 0$
on $\mathrm{Conf}_n(X)$ (Theorem~\ref{thm:arnold-three}), and the
cross-terms cancel by the Koszul sign rule.

\emph{(ii) Super-antisymmetry.} Immediate from the commutator form.

\emph{(iii) Super-Jacobi.} The cup product $\smile$ is associative
(both $\mu$ and $\bar\Delta$ are coassociative), so the graded
commutator $[\cdot,\cdot]$ is the super-commutator of an associative
product, and super-Jacobi is automatic
(Loday--Vallette 2012 \S 1.1.11).

\emph{(iv) Derivation.} Compute
$D[\phi,\psi](w) = d_A\bigl(\mu(\phi\otimes\psi)\bar\Delta w\bigr)
- \mu(\phi\otimes\psi)\bar\Delta(d_B w) \mp
d_A\bigl(\mu(\psi\otimes\phi)\bar\Delta w\bigr) \pm
\mu(\psi\otimes\phi)\bar\Delta(d_B w)$;
using $d_A$-compatibility with $\mu$ (it \emph{is} a chain map for
Etingof--Kazhdan) and $d_B$-compatibility with $\bar\Delta$
(Leibniz rule on deconcatenation, valid since $d_B$ is a
coderivation of $\bar\Delta$), the right-hand side regroups as
$[D\phi,\psi] + (-1)^p[\phi,D\psi]$.
\end{proof}

\begin{remark}[Weight completion for infinite-dimensional $A$]
\label{rem:c1-weight-completion-k3-convolution}
Because $\mathbf{H}_{\Delta_5}$ is infinite-dimensional (the imaginary
cone has one Fourier mode per positive Gritsenko--Nikulin coefficient
$c_{\Delta_5}(N,\ell)$), the naive
$\mathrm{Hom}(B^{\mathrm{ord}}(A),A)$ carries no sensible grading.
The correct ambient category is the weight-completed pro-category
$\mathrm{Hom}_{\widehat{\Lambda}_{\mathrm{Muk}}}$ with respect to the
Mukai lattice $\Lambda_{\mathrm{Muk}} = \mathrm{II}_{4,4}\oplus\mathrm{II}_{2,2}$
grading on $\mathbf{H}_{\Delta_5}$ (each Fourier mode carries a definite
Mukai-lattice weight; the bar-tensor is
weight-additive), and complete with respect to the filtration by
total weight $\ge N$. In every fixed weight piece
$\mathrm{Hom}^{\le W}(B^{\mathrm{ord}}(A),A)$, only finitely many
bar-degrees contribute (by the positive-cone hypothesis on
$\Lambda_{\mathrm{Muk}}$), so the dGLA is finite-dimensional per
degree per weight piece. The Maurer--Cartan equation is solved in the
weight-adic completion
$\mathfrak{C}_{\mathbf{H}_{\Delta_5}}[[\hbar]] =
\varprojlim_W \mathfrak{C}_{\mathbf{H}_{\Delta_5}}^{\le W}[[\hbar]]$.
This is the chain-level pro-object / weight-completed ambient in which
the class-M Maurer--Cartan tower is well-defined.
\end{remark}

\begin{definition}[K3 Maurer--Cartan input data; \ClaimStatusDefinitional]
\label{rem:theta-universal-MC-k3}
For the ordered K3 Hall model, write a degree-one series
\[
 \Theta_{\mathbf{H}_{\Delta_5}}
 \;=\; \sum_{n\ge 3} \hbar^n\,\phi^{(n)}
 \;\in\; \mathfrak{C}^1_{\mathbf{H}_{\Delta_5}}[[\hbar]],
\]
with $\phi^{(n)}\in \mathrm{Hom}^1(B^{\mathrm{ord}}_n(\mathbf{H}_{\Delta_5}),
\mathbf{H}_{\Delta_5})$ the order-$n$ pentagon cocycle in the
sense of Etingof--Kazhdan. The K3 construction supplies the leading
coefficient
\[
 \phi^{(3)}
 \;=\;
 \zeta(3)\,c_{\mathrm{symm}}
 \;+\;\tfrac{25}{3}\,c_{\mathrm{timelike}}
 \;+\; (\Phi_{10}^{\mathrm{un}}/\eta^{24})\, c_{\Phi_{10}^{\mathrm{un}}},
\]
and the next nontrivial correction
$\phi^{(6)} = \zeta(3)^2\, c_6^{(1)} + \tfrac{25}{3}\zeta(3)\, c_6^{(2)}
+ (25/3)^2\, c_6^{(3)} + \zeta(3)(\Phi_{10}^{\mathrm{un}}/\eta^{24})c_6^{(4)}
+ (\Phi_{10}^{\mathrm{un}}/\eta^{24})^2 c_6^{(5)} + \zeta(5)\, c_6^{(6)}$,
where $c_6^{(i)}$ are the six MZV--Borcherds-leg cohomology
representatives. These coefficients are input data for this section.
Only the ordered-convolution implications from explicit finite
identities among the $\phi^{(n)}$ are asserted here; no cross-volume
statement is used as a proof of those identities in this chapter.
\end{definition}

\begin{theorem}[Finite Maurer--Cartan window at $\hbar^3,\hbar^4$;
\ClaimStatusProvedHere]
\label{thm:MC-hbar3-hbar4-k3}
On the weight-completed dGLA
$\mathfrak{C}_{\mathbf{H}_{\Delta_5}}[[\hbar]]$ of
Definition~\ref{def:convolution-dgla-k3}, let
$\Theta_{\mathbf{H}_{\Delta_5}}$ of
Definition~\ref{rem:theta-universal-MC-k3} satisfy
\[
\phi^{(0)}=\phi^{(1)}=\phi^{(2)}=0,\qquad
D\phi^{(3)}=0,\qquad \phi^{(4)}=0.
\]
Then the Maurer--Cartan equation
\[
 D\,\Theta_{\mathbf{H}_{\Delta_5}} \;+\; \tfrac{1}{2}\,
 [\Theta_{\mathbf{H}_{\Delta_5}},\Theta_{\mathbf{H}_{\Delta_5}}]
 \;=\; 0
\]
holds modulo $\hbar^5$.
\end{theorem}

\begin{proof}
Collecting terms by powers of $\hbar$: the $\hbar^0,\hbar^1,\hbar^2$
terms vanish because $\Theta$ starts at $\hbar^3$ and $D,\ [\cdot,\cdot]$
preserve $\hbar$-degree. At $\hbar^3$:
$D\phi^{(3)} = 0$ by hypothesis.
At $\hbar^4$: $D\phi^{(4)} + \tfrac{1}{2}[\phi^{(3)},\phi^{(3)}]_{\hbar^4}
= 0$. The commutator has no $\hbar^4$ component because both inputs
begin at $\hbar^3$; the first possible bracket is at $\hbar^6$.
Thus the $\hbar^4$ equation reduces to $D\phi^{(4)}=0$, which follows
from $\phi^{(4)}=0$.
\end{proof}

\begin{theorem}[Finite Maurer--Cartan criterion through $\hbar^{12}$;
\ClaimStatusProvedHere]
\label{thm:MC-hbar7-hbar12-k3}
On the weight-completed pro-dGLA
$\mathfrak{C}_{\mathbf{H}_{\Delta_5}}[[\hbar]]$, suppose
$\Theta=\sum_{n\ge3}\hbar^n\phi^{(n)}$ has support in
$3\mathbb Z_{\ge1}$ through order~$12$:
\[
\phi^{(n)}=0
\quad
\text{for } n\in\{4,5,7,8,10,11\}.
\]
Assume the three finite identities
\[
D\phi^{(6)}+\tfrac12[\phi^{(3)},\phi^{(3)}]=0,
\]
\[
D\phi^{(9)}+[\phi^{(3)},\phi^{(6)}]=0,
\]
\[
D\phi^{(12)}+[\phi^{(3)},\phi^{(9)}]
+\tfrac12[\phi^{(6)},\phi^{(6)}]=0.
\]
Then $\Theta$ satisfies the Maurer--Cartan equation modulo
$\hbar^{13}$.
\end{theorem}

\begin{proof}
The coefficient of $\hbar^n$ in
$D\Theta+\tfrac12[\Theta,\Theta]$ is
\[
D\phi^{(n)}
+\tfrac12\sum_{\substack{i+j=n\\ i,j\ge3}}
[\phi^{(i)},\phi^{(j)}].
\]
For $n=3,4$ this reduces to
Theorem~\ref{thm:MC-hbar3-hbar4-k3}. For $n=5$ the coefficient is
$D\phi^{(5)}$ and vanishes by the support hypothesis.
For $n=6,9,12$ the displayed hypotheses are exactly the three
nontrivial coefficient equations. For
$n\in\{7,8,10,11\}$ each summand either contains
$\phi^{(4)}$, $\phi^{(5)}$, $\phi^{(7)}$, $\phi^{(8)}$,
$\phi^{(10)}$, or $\phi^{(11)}$, or has an impossible split into
positive multiples of~$3$; hence it vanishes. Therefore every
coefficient through $\hbar^{12}$ is zero.
\end{proof}

\section{Admissibility filtration on the pentagon tower}
\label{sec:bc-admissibility-filtration}

On the $A_\infty$-Humbert regime of the pentagon tower
$\{\phi^{(n)}\}_{n\ge 3}$, the strict $3\mathbb{Z}$-parity of
Theorem~\ref{thm:MC-hbar7-hbar12-k3} dissolves into a denser support on
which every $\phi^{(n)}$ admits a Fourier-theoretic primitive-count
description. The pentagon coherences index via the discriminant
$\mathfrak d_n = (n-3)/2$ against the Jacobi-index-$1$ lattice $4NM - \ell^2$
carried by the K3 elliptic genus $\phi^{K3}_{0,1}$. The question of
which $\phi^{(n)}$ survives is then not a motivic-depth question but a
polar-support question, and Eichler--Zagier $1985$ supplies the
polar cutoff.

\begin{lemma}[Polar support of the K3 elliptic genus]
\label{lem:bc-polar-support-phi-K3}
\ClaimStatusProvedElsewhere
The K3 elliptic genus $\phi^{K3}_{0,1}$ has Fourier coefficients
$c(N, \ell) = C(4N - \ell^2)$ depending only on the discriminant
$\Delta = 4N - \ell^2$, with
$C(-1) = 2$, $C(0) = 20$, and
$C(\Delta) = 0$ for all $\Delta < -1$. The primary source is
Eichler--Zagier $1985$ \textup{Prog.~Math.}~$55$ Theorem~$9.1$
(index-$m$ weak Jacobi forms vanish in discriminants below $-m^2$, here
$m = 1$).
\end{lemma}

\begin{theorem}[Polar depth-reduction of the pentagon tower on the K3
side]
\label{thm:bc-polar-depth-reduction}
\ClaimStatusConditional
Assume the pentagon $A_\infty$-tower $\{\phi^{(n)}\}_{n\ge 3}$ of
$\mathbf{H}_{\Delta_5}$ on the $A_\infty$-Humbert regime has coefficient
formula
\[
 \phi^{(n)} \;=\; \sum_{\substack{(N, \ell, M)\\ 4NM - \ell^2 = -\mathfrak d_n}}
 c_{\Phi_{10}^{\mathrm{un}}/\eta^{24}}(N, \ell, M)\,
 [\mathrm{generator}]^{\otimes n},
 \qquad \mathfrak d_n = (n-3)/2.
\]
Admissibility $\mathfrak d_n \bmod 4 \in \{0, 1\}$, equivalently
$n \equiv 3, 5 \pmod 8$, is necessary for the lattice sum to be
non-empty. On the admissible locus:
\begin{itemize}
\item $n = 3$: $\phi^{(3)} \ne 0$, the Etingof--Kazhdan pentagon
 cocycle.
\item $n = 5$: $\phi^{(5)} = 2\,[\mathrm{generator}]^{\otimes 5} \ne 0$,
 one-term lattice sum with Heegner coefficient
 $C(-1) = 2$.
\item $n \ge 11$ admissible: $\phi^{(n)} = 0$ because every lattice
 representative $(N, \ell, M)$ has $\Delta = -\mathfrak d_n \le -4 < -1$, and
 $C(\Delta) = 0$ for $\Delta < -1$ by
 Lemma~\ref{lem:bc-polar-support-phi-K3}.
\end{itemize}
The non-admissible $n$'s have $\phi^{(n)} = 0$ by empty lattice sum.
Under the displayed coefficient formula, the $A_\infty$-Humbert
pentagon tower is concentrated at
$n \in \{3, 5\}$: the strict quasi-Hopf obstruction plus a single
higher coherence at $n = 5$. The vanishing for $n \ge 11$ is
\emph{unconditional relative to the coefficient formula}: it follows
from Eichler--Zagier $1985$ Theorem~$9.1$
(index-$m$ polar support $\Delta \ge -m^2$, here $m = 1$), a
structural property of weak Jacobi forms that does not depend on
motivic-depth arithmetic.
\end{theorem}

\begin{proof}
The admissibility congruence $\mathfrak d_n \bmod 4 \in \{0, 1\}$ is the
discriminant form's signature: $4NM - \ell^2 \equiv -\ell^2 \pmod 4$
takes values in $\{0, -1\} \pmod 4$, so $-\mathfrak d_n$ is representable by the
form iff $\mathfrak d_n \in \{0, 1\} \pmod 4$. Combined with the odd-$n$
constraint (else $\mathfrak d_n \notin \mathbb{Z}$), this is $n \equiv 3, 5
\pmod 8$.

Case $n = 3$: $\mathfrak d_3 = 0$, the lattice sum has representative
$(N, \ell, M) = (0, 0, 0)$ or any $(N, \ell, M)$ with $4NM = \ell^2$;
Heegner coefficient $C(0) = 20 \ne 0$; $\phi^{(3)}$ is non-vanishing
and is the strict pentagon obstruction of Drinfeld $1990$ Problem
$3.16$ transported to the Manin-pair quantisation.

Case $n = 5$: $\mathfrak d_5 = 1$, lattice representative $(0, 1, 0)$ gives
$\Delta = -1$; $C(-1) = 2 \ne 0$; the lattice sum is one-term and
$\phi^{(5)} = 2\,[\mathrm{generator}]^{\otimes 5}$.

Case $n \ge 11$ admissible ($n \in \{11, 13, 19, 21, 27, 29, \ldots\}$):
$\mathfrak d_n \ge 4$, so every lattice representative $(N, \ell, M)$ of
$-\mathfrak d_n$ has $\Delta = -\mathfrak d_n \le -4$. By
Lemma~\ref{lem:bc-polar-support-phi-K3} (Eichler--Zagier $1985$
Theorem~$9.1$), $C(\Delta) = 0$ for $\Delta < -1$. Hence every
summand in the lattice sum vanishes and $\phi^{(n)} = 0$.

Case $n$ non-admissible: the congruence $\mathfrak d_n \in \{0, 1\} \pmod 4$
fails, so there exists no $(N, \ell, M)\in\mathbb{Z}_{\ge 0}\times
\mathbb{Z}\times\mathbb{Z}_{\ge 0}$ with $4NM - \ell^2 = -\mathfrak d_n$; the
lattice sum is empty and $\phi^{(n)} = 0$ by vacuity.

The independence statement is this: the finite-window
Maurer--Cartan criterion of
Theorem~\ref{thm:MC-hbar7-hbar12-k3} reduces the strict quasi-Hopf
calculation to explicit coefficient identities, while the present
vanishing does not use motivic-depth identities at all once the
$A_\infty$-Humbert coefficient formula is fixed. The vanishing of
$\phi^{(n)}$ for admissible $n \ge 11$ comes from Eichler--Zagier
$1985$ Theorem~$9.1$, which is unconditional and structural to the
theory of index-$1$ weak Jacobi forms, reaching all
$n$ with $\mathfrak d_n > 1$.
\end{proof}

\begin{remark}[Contrast with motivic-depth input]
\label{rem:bc-brown-horizon-contrast}
Theorem~\ref{thm:MC-hbar7-hbar12-k3} is a formal finite-window
criterion: it identifies the coefficient identities needed through
$\hbar^{12}$. The motivic verification of those identities is a
separate arithmetic input. Theorem~\ref{thm:bc-polar-depth-reduction}
is orthogonal: on the $A_\infty$-Humbert regime the vanishing is
polar-cutoff driven and reaches arbitrarily high $n$, without any
motivic-basis input. The
admissibility table through $n \le 30$ is obtained directly by listing
Jacobi-lattice representatives $(N,\ell,M)$ and evaluating the Heegner
coefficients $H_n=C(-\mathfrak d_n)$.
\end{remark}
